% Leçon 33 — Éléments socioculturels : le marketing en Chine et au Cameroun — Chinois Tle A — Cours signé OnBuch+
% Programme officiel MINESEC. Compiler avec : tectonic ZH33-socio-le-marketing-en-chine-et-au-cameroun.tex
\newcommand{\DOCMATIERE}{Chinois}
\newcommand{\DOCNIVEAU}{Tle A}
\newcommand{\DOCMODULE}{Module 4 : Activités économiques et monde du travail}
\newcommand{\DOCLECON}{ZH33}
\newcommand{\DOCTITRE}{Leçon 33 — Éléments socioculturels : le marketing en Chine et au Cameroun}
% OnBuch+ — préambule « cours signés » ENRICHI (compilé avec tectonic / XeLaTeX)
% Système visuel « Pop » d'OnBuch+ : fond crème (jamais blanc), bordures encre
% épaisses, ombres « dures » décalées, titres Archivo Black en capitales.
% Version MATHÉMATIQUES : graphes (pgfplots/tikz), boîtes pédagogiques
% (définition, propriété, méthode, attention, le savais-tu, exercice résolu,
% exercice, TP, bilan), page de garde et signature OnBuch+ ancrée en bas.
%
% Macros attendues dans main.tex AVANT \input{preamble} :
%   \DOCMATIERE  \DOCNIVEAU  \DOCMODULE  \DOCLECON (numéro)  \DOCTITRE
\documentclass[11pt]{article}

\usepackage{fontspec}
\usepackage[french]{babel} % français = langue principale ; le chinois via \zh{..} / textebox (xeCJK)
\usepackage{xeCJK}
\usepackage{pifont} % \ding{51} \ding{55} utilisés par les modèles
\usepackage[a4paper,margin=2cm,top=3.4cm,bottom=2.8cm,headheight=1.4cm,footskip=1.3cm]{geometry}
\usepackage{amsmath,amssymb,amsfonts}
\usepackage{mathtools} % \xmapsto, \coloneqq... souvent utilisés par les modèles
\usepackage{cancel} % simplifications barrées (bilans de forces)
% \abs{z} (module, valeur absolue) et \norm{u} (norme), utilisés par les modèles
\providecommand{\abs}{}\let\abs\relax\DeclarePairedDelimiter{\abs}{\lvert}{\rvert}
\providecommand{\norm}{}\let\norm\relax\DeclarePairedDelimiter{\norm}{\lVert}{\rVert}
\usepackage[table]{xcolor} % [table] : fournit aussi \rowcolors (alternance de lignes)
\usepackage{pagecolor}
\usepackage{tikz}
\usetikzlibrary{arrows.meta,positioning,calc,shapes.geometric,shapes.misc,shapes.arrows,decorations.pathmorphing,decorations.pathreplacing,decorations.markings,patterns,shadows,mindmap,trees,fit,backgrounds,matrix,intersections,angles,quotes}
\usepackage{pgfplots}
\usepackage[version=4]{mhchem} % équations nucléaires \ce{^{235}_{92}U}
\usepackage[european,siunitx]{circuitikz} % schémas électriques
\pgfplotsset{compat=1.18}
\usepgfplotslibrary{fillbetween,groupplots} % aires entre courbes (calcul intégral)
\usepackage{enumitem}
\usepackage{fancyhdr}
% [most] charge toutes les bibliothèques tcolorbox (listings, minted, raster,
% documentation...) et gonfle inutilement la mémoire TeX sur un document long
% (~300 boîtes) : on ne charge que ce qui est réellement utilisé.
\usepackage[skins,breakable]{tcolorbox}
\usepackage{graphicx}
\usepackage{colortbl}
\usepackage{booktabs}
\usepackage{tabularx}
\usepackage{diagbox} % case d'en-tête diagonale des tableaux à double entrée
% Colonnes extensibles centrée / gauche / droite (C, L, R), souvent utilisées par les modèles
\newcolumntype{C}{>{\centering\arraybackslash}X}
\newcolumntype{L}{>{\raggedright\arraybackslash}X}
\newcolumntype{R}{>{\raggedleft\arraybackslash}X}
\usepackage{array}
\usepackage{multicol}
\usepackage{multirow}
\usepackage{siunitx}
% parse-numbers=false : accepte aussi \SI{2,5 \times 10^{-3}}{...}
\sisetup{locale=FR,output-decimal-marker={,},group-minimum-digits=4,parse-numbers=false}
\DeclareSIUnit\ohm{\ensuremath{\symup{\Omega}}} % U+2126 absent de la police
\DeclareSIUnit\division{div} % réglages d'oscilloscope (ms/div, V/div)
\DeclareSIUnit\tr{tr} % tours (vitesse de rotation, ex. \SI{1500}{\tr\per\minute})
\DeclareSIUnit\molar{M} % concentration molaire (ex. \SI{3}{\molar}, \SI{1}{\micro\molar})
\DeclareSIUnit\an{an} % année (le modèle confond parfois avec \year, un compteur TeX) — ex. \SI{5}{\kilo\meter\per\an}
\DeclareSIUnit\FCFA{FCFA} % franc CFA (ex. \SI{500}{\FCFA\per\kilo\gram})
\DeclareSIUnit\degreeBrix{\textdegree Brix} % teneur en sucre d'un sirop/jus (réfractométrie)
\DeclareSIUnit\month{mois} % le modèle utilise parfois \month, un compteur TeX, comme unité de durée
\usepackage{qrcode}
% \degree : commande fréquemment utilisée par les modèles pour un angle ou une
% température en dehors de \SI{}{} (non fournie par les paquets chargés).
\providecommand{\degree}{^{\circ}}
% \qed (fin de démonstration), utilisé par les modèles sans amsthm
\providecommand{\qed}{\hfill$\square$}
% \barre : double barre des valeurs interdites dans un tableau de variations (tkz-tab)
\providecommand{\barre}{\ensuremath{\|}}
% environnement proof (amsthm non chargé) utilisé par les modèles
\ifdefined\proof\else\newenvironment{proof}[1][Démonstration]{\par\noindent\textit{#1.}\ }{\qed\par}\fi
\usepackage{lastpage}
\usepackage{hyperref}
\hypersetup{hidelinks}

% ============================================================
% Palette « Pop » OnBuch+ — mêmes hex que la classe P (plus_ui.dart)
% ============================================================
\definecolor{poporange}{HTML}{F59321}
\definecolor{popdark}{HTML}{E07A0C}
\definecolor{popcream}{HTML}{FFF6E8}
\definecolor{popink}{HTML}{1C1714}
\definecolor{poppurple}{HTML}{7A5AE0}
\definecolor{popgreen}{HTML}{2E9E6A}
\definecolor{poppink}{HTML}{E0457B}
\definecolor{popblue}{HTML}{2F7DE1}
\definecolor{popgold}{HTML}{C98A12}
\definecolor{popmuted}{HTML}{8A7F76}
\definecolor{popline}{HTML}{EADFD2}
\definecolor{popgray}{HTML}{8A7F76}
\colorlet{popgrey}{popgray}
% Alias tolérants (le modèle nomme parfois les couleurs différemment)
\colorlet{popwhite}{white}
\colorlet{popblack}{popink}
\colorlet{popred}{poppink}
\colorlet{popyellow}{popgold}
% Teintes claires (fonds de tableaux/encadrés) — le modèle nomme parfois
% les couleurs avec un suffixe L (ex. popblueL) sans les avoir définies.
\colorlet{poporangeL}{poporange!20}
\colorlet{popdarkL}{popdark!20}
\colorlet{poppurpleL}{poppurple!20}
\colorlet{popgreenL}{popgreen!20}
\colorlet{poppinkL}{poppink!20}
\colorlet{popblueL}{popblue!20}
\colorlet{popgoldL}{popgold!20}
\colorlet{popredL}{popred!20}
\colorlet{popgrayL}{popgray!20}
% Teintes claires pour fonds de tableaux / zones
\colorlet{popblueL}{popblue!10!white}
\colorlet{poporangeL}{poporange!14!white}
\colorlet{popgreenL}{popgreen!12!white}
\colorlet{poppurpleL}{poppurple!12!white}
\colorlet{poppinkL}{poppink!12!white}
\colorlet{popgoldL}{popgold!14!white}

\pagecolor{popcream}

% ============================================================
% Typographie
% ============================================================
\defaultfontfeatures{Ligatures=TeX}
\setmainfont{PlusJakartaSans-Medium.ttf}[Path=../../../fonts/,BoldFont=PlusJakartaSans-Bold.ttf]
% Symboles, API et emojis absents de la police principale : repli sur DejaVu Sans (caractères actifs)
\newfontface\symfont{DejaVuSans.ttf}[Path=../../../fonts/]
\makeatletter
\newcommand\symactive[1]{\catcode#1=13 \begingroup\lccode`\~=#1 \lowercase{\endgroup\def~}{{\symfont\char#1}}}
\newcommand\symblank[1]{\catcode#1=13 \begingroup\lccode`\~=#1 \lowercase{\endgroup\def~}{}}
\newcommand\symhyph[1]{\catcode#1=13 \begingroup\lccode`\~=#1 \lowercase{\endgroup\def~}{\mbox{-}}}
\newcount\symc
\newcommand\symrange[2]{\symc=#1 \@whilenum\symc<#2 \do{\expandafter\symactive\expandafter{\the\symc}\advance\symc by 1 }}
\newcommand\symblankrange[2]{\symc=#1 \@whilenum\symc<#2 \do{\expandafter\symblank\expandafter{\the\symc}\advance\symc by 1 }}
\makeatother
\symrange{"0250}{"02FF}\symactive{"2713}\symactive{"2714}\symactive{"2717}\symactive{"2718}\symactive{"2610}\symactive{"2611}\symactive{"2612}
\symrange{"2600}{"27BF}
\catcode"2705=13 \begingroup\lccode`\~="2705 \lowercase{\endgroup\def~}{{\symfont\char"2713}}\symblank{"26BD}\symblank{"26A1}
\symrange{"2460}{"257F}\symactive{"223C}\symactive{"2248}\symactive{"2260}
\symactive{"01F8}\symactive{"01F9}\symactive{"1E3E}\symactive{"1E3F}
\symhyph{"2011}\symhyph{"2E1A}\symblankrange{"1F300}{"1FAFF}\symblank{"FE0F}
% Caractères chinois : WenQuanYi Zen Hei (sous-ensemble GB2312 + pinyin), gras/italique simulés
\setCJKmainfont{WenQuanYiZenHei-subset.ttf}[Path=../../../fonts/,AutoFakeBold=3,AutoFakeSlant=0.2]
\xeCJKsetup{CJKspace=false}
% Pinyin : voyelles à caron (ǎ ǐ ǒ ǔ ǖ ǘ ǚ ǜ) absentes de la police principale -> caractères actifs
\newfontface\pyfont{WenQuanYiZenHei-subset.ttf}[Path=../../../fonts/]
\newcommand\pyactive[1]{\catcode#1=13 \begingroup\lccode`\~=#1 \lowercase{\endgroup\def~}{{\pyfont\char#1}}}
\pyactive{"01CD}\pyactive{"01CE}\pyactive{"01CF}\pyactive{"01D0}\pyactive{"01D1}\pyactive{"01D2}\pyactive{"01D3}\pyactive{"01D4}\pyactive{"01D5}\pyactive{"01D6}\pyactive{"01D7}\pyactive{"01D8}\pyactive{"01D9}\pyactive{"01DA}\pyactive{"01DB}\pyactive{"01DC}
% Maths en sans-serif (Fira Math) avec chiffres et lettres droites de Plus
% Jakarta, pour que les formules se fondent dans le texte.
\usepackage[math-style=ISO,bold-style=ISO]{unicode-math}
\setmathfont{FiraMath-Regular.otf}[Path=../../../fonts/]
\setmathfont{PlusJakartaSans-Medium.ttf}[Path=../../../fonts/,range={up/num,up/{Latin,latin},\mathup}]
\usepackage{icomma} % virgule décimale sans espace en mode math
% Caractères mathématiques Unicode absents des polices de texte (Plus Jakarta,
% Archivo Black) : rendus par leur équivalent mathématique, en texte comme en
% formule — sinon ils s'affichent en carré vide (ex. « Arithmétique dans ℤ »).
\usepackage{newunicodechar}
\newunicodechar{ℝ}{\ensuremath{\mathbb{R}}}
\newunicodechar{ℤ}{\ensuremath{\mathbb{Z}}}
\newunicodechar{ℂ}{\ensuremath{\mathbb{C}}}
\newunicodechar{ℕ}{\ensuremath{\mathbb{N}}}
\newunicodechar{ℚ}{\ensuremath{\mathbb{Q}}}
\newunicodechar{𝔻}{\ensuremath{\mathbb{D}}}
\newunicodechar{α}{\ensuremath{\alpha}}
\newunicodechar{β}{\ensuremath{\beta}}
\newunicodechar{γ}{\ensuremath{\gamma}}
\newunicodechar{δ}{\ensuremath{\delta}}
\newunicodechar{ε}{\ensuremath{\varepsilon}}
\newunicodechar{θ}{\ensuremath{\theta}}
\newunicodechar{λ}{\ensuremath{\lambda}}
\newunicodechar{μ}{\ensuremath{\mu}}
\newunicodechar{σ}{\ensuremath{\sigma}}
\newunicodechar{τ}{\ensuremath{\tau}}
\newunicodechar{φ}{\ensuremath{\varphi}}
\newunicodechar{ψ}{\ensuremath{\psi}}
\newunicodechar{ω}{\ensuremath{\omega}}
\newunicodechar{Σ}{\ensuremath{\Sigma}}
% majuscules produites par \MakeUppercase dans les titres (α-aminés -> Α)
\newunicodechar{Α}{\ensuremath{\alpha}}
\newunicodechar{Β}{\ensuremath{\beta}}
\newunicodechar{Γ}{\ensuremath{\Gamma}}
\newunicodechar{Δ}{\ensuremath{\Delta}}
\newunicodechar{Ε}{\ensuremath{\varepsilon}}
\newunicodechar{Θ}{\ensuremath{\Theta}}
\newunicodechar{Λ}{\ensuremath{\Lambda}}
\newunicodechar{Μ}{\ensuremath{\mu}}
\newunicodechar{Τ}{\ensuremath{\tau}}
\newunicodechar{Φ}{\ensuremath{\Phi}}
\newunicodechar{Ψ}{\ensuremath{\Psi}}
\newunicodechar{Ω}{\ensuremath{\Omega}}
\newunicodechar{↦}{\ensuremath{\mapsto}}
\newunicodechar{∈}{\ensuremath{\in}}
\newunicodechar{∉}{\ensuremath{\notin}}
\newunicodechar{∧}{\ensuremath{\wedge}}
\newunicodechar{⇔}{\ensuremath{\Leftrightarrow}}
\newunicodechar{⟺}{\ensuremath{\Longleftrightarrow}}
\newunicodechar{⇒}{\ensuremath{\Rightarrow}}
\newunicodechar{⟹}{\ensuremath{\Longrightarrow}}
\newunicodechar{∘}{\ensuremath{\circ}}
\newunicodechar{∪}{\ensuremath{\cup}}
\newunicodechar{∩}{\ensuremath{\cap}}
\newunicodechar{≡}{\ensuremath{\equiv}}
\newunicodechar{⊂}{\ensuremath{\subset}}
\newunicodechar{⊥}{\ensuremath{\perp}}
\newunicodechar{∃}{\ensuremath{\exists}}
\newunicodechar{∀}{\ensuremath{\forall}}
\newunicodechar{∛}{\ensuremath{\sqrt[3]{\,}}}
\newunicodechar{∜}{\ensuremath{\sqrt[4]{\,}}}
\newunicodechar{⃗}{\ensuremath{{}^{\to}}}
\newunicodechar{ⁿ}{\ifmmode^{n}\else\textsuperscript{\ensuremath{n}}\fi}
\newunicodechar{ˣ}{\ifmmode^{x}\else\textsuperscript{\ensuremath{x}}\fi}
\newunicodechar{ᵇ}{\ifmmode^{b}\else\textsuperscript{\ensuremath{b}}\fi}
\newunicodechar{ʳ}{\ifmmode^{r}\else\textsuperscript{\ensuremath{r}}\fi}
\newunicodechar{ᵘ}{\ifmmode^{u}\else\textsuperscript{\ensuremath{u}}\fi}
\newunicodechar{ʸ}{\ifmmode^{y}\else\textsuperscript{\ensuremath{y}}\fi}
\newunicodechar{ᵃ}{\ifmmode^{a}\else\textsuperscript{\ensuremath{a}}\fi}
\newunicodechar{ᵉ}{\ifmmode^{e}\else\textsuperscript{\ensuremath{e}}\fi}
\newunicodechar{ᵏ}{\ifmmode^{k}\else\textsuperscript{\ensuremath{k}}\fi}
\newunicodechar{ᵖ}{\ifmmode^{p}\else\textsuperscript{\ensuremath{p}}\fi}
\newunicodechar{ᴺ}{\ifmmode^{N}\else\textsuperscript{\ensuremath{N}}\fi}
\newunicodechar{ᶜ}{\ifmmode^{c}\else\textsuperscript{\ensuremath{c}}\fi}
\newunicodechar{ʰ}{\ifmmode^{h}\else\textsuperscript{\ensuremath{h}}\fi}
\newunicodechar{ᵠ}{\ifmmode^{q}\else\textsuperscript{\ensuremath{q}}\fi}
\newunicodechar{ₙ}{\ifmmode_{n}\else\textsubscript{\ensuremath{n}}\fi}
\newunicodechar{ₐ}{\ifmmode_{a}\else\textsubscript{\ensuremath{a}}\fi}
\newunicodechar{ᵢ}{\ifmmode_{i}\else\textsubscript{\ensuremath{i}}\fi}
\newunicodechar{ᵣ}{\ifmmode_{r}\else\textsubscript{\ensuremath{r}}\fi}
\newunicodechar{ₖ}{\ifmmode_{k}\else\textsubscript{\ensuremath{k}}\fi}
\newunicodechar{ᵦ}{\ifmmode_{\beta}\else\textsubscript{\ensuremath{\beta}}\fi}
\newunicodechar{ₓ}{\ifmmode_{x}\else\textsubscript{\ensuremath{x}}\fi}
\newfontfamily\popdisplay{ArchivoBlack-Regular.ttf}[Path=../../../fonts/]
\newfontfamily\poplogo{SpaceGrotesk-Bold.ttf}[Path=../../../fonts/]
\newfontfamily\popsemi{PlusJakartaSans-SemiBold.ttf}[Path=../../../fonts/]
\newfontfamily\popxbold{PlusJakartaSans-ExtraBold.ttf}[Path=../../../fonts/]
\newfontfamily\popxboldls{PlusJakartaSans-ExtraBold.ttf}[Path=../../../fonts/,LetterSpace=3.0]

\setlist[enumerate]{leftmargin=1.7em,itemsep=5pt,topsep=4pt}
\setlist[itemize]{leftmargin=1.7em,itemsep=4pt,topsep=4pt,label={\color{poporange}\textbullet}}
\setlength{\parindent}{0pt}
\setlength{\parskip}{5pt}
\renewcommand{\arraystretch}{1.25}

% pgfplots : style Pop par défaut
\pgfplotsset{
  every axis/.append style={axis line style={popink,line width=1pt},tick style={popink},
    grid style={popline,line width=0.5pt},label style={font=\small},tick label style={font=\footnotesize}},
  popcurve/.style={poporange,line width=1.8pt,no markers,smooth},
}
% Même style disponible dans un \draw TikZ simple (hors pgfplots)
\tikzset{popcurve/.style={poporange,line width=1.8pt}}
\tikzset{popgrid/.style={popline,very thin}}
\colorlet{popcurve}{poporange} % le nom du style est parfois utilisé comme couleur

% ============================================================
% Pill / squiggle
% ============================================================
\newcommand{\poppill}[3]{%
  \tikz[baseline=-0.55ex]\node[draw=popink,line width=1.2pt,rounded corners=2.6pt,%
    fill=#2,inner xsep=6.5pt,inner ysep=3pt]{%
    \popxboldls\fontsize{7}{7}\selectfont\color{#3}\MakeUppercase{#1}};%
}
\newcommand{\popsquiggle}[1]{%
  \tikz[baseline=-0.4ex]{
    \draw[#1,line width=2.6pt,line cap=round]
      plot [smooth] coordinates {(0,0.09) (0.34,0.01) (0.68,0.21) (1.02,0.01) (1.36,0.10)};
  }%
}
% Mot-clé surligné (marqueur orange sous le texte)
\newcommand{\cle}[1]{{\popxbold\color{popdark}#1}}

% ============================================================
% En-tête / pied
% ============================================================
\pagestyle{fancy}
\fancyhf{}
\renewcommand{\headrulewidth}{0pt}
\renewcommand{\footrulewidth}{0pt}
\lhead{\raisebox{-0.15ex}{\poplogo\fontsize{13}{13}\selectfont\color{popink}OnBuch\color{poporange}+}}
\rhead{\poppill{\DOCMATIERE\ \textbullet\ \DOCNIVEAU\ \textbullet\ Leçon \DOCLECON}{white}{popink}}
\lfoot{\popxbold\fontsize{7.6}{7.6}\selectfont\color{popmuted}\MakeUppercase{Cours signé OnBuch+}\ \textbullet\ {\popsemi\DOCTITRE}}
\rfoot{\tikz[baseline=-0.6ex]\node[rounded corners=8.5pt,fill=popink,minimum height=17pt,minimum width=17pt,inner xsep=5pt,inner sep=0]{\popxbold\fontsize{8}{8}\selectfont\color{popcream}\thepage\,/\,\pageref*{LastPage}};}

% ============================================================
% Titres
% ============================================================
\newcommand{\chaptitre}[2]{%
  \begin{tcolorbox}[enhanced,colback=poporange,colframe=popink,boxrule=2.6pt,arc=11pt,
      drop shadow=popink,left=15pt,right=15pt,top=12pt,bottom=11pt,before skip=3pt,after skip=18pt]
    \poppill{#2}{popink}{popcream}\par\vspace{9pt}
    {\raggedright\hyphenpenalty=10000\exhyphenpenalty=10000\popdisplay\fontsize{22}{24}\selectfont\color{popcream}\MakeUppercase{#1}\par}
  \end{tcolorbox}
}
% Section numérotée : pastille orange avec le numéro + titre Archivo
\newcounter{popsec}
\newcommand{\coursec}[1]{%
  \stepcounter{popsec}\setcounter{popsub}{0}%
  \par\vspace{16pt}\noindent
  \begin{minipage}{\linewidth}
  \tikz[baseline=-0.7ex]\node[draw=popink,line width=1.6pt,fill=poporange,rounded corners=4pt,minimum size=22pt,inner sep=0]{\popdisplay\fontsize{12}{12}\selectfont\color{popcream}\thepopsec};%
  \hspace{8pt}{\popdisplay\fontsize{15}{17}\selectfont\color{popink}\MakeUppercase{#1}}\par
  \vspace{4pt}\noindent{\color{poporange}\rule{36pt}{3.6pt}}
  \end{minipage}\par\nopagebreak\vspace{8pt}
}
\newcounter{popsub}
\newcommand{\courssub}[1]{%
  \stepcounter{popsub}%
  \par\vspace{9pt}\noindent{\popxbold\fontsize{11.5}{13}\selectfont\color{popink}{\color{poporange}\thepopsec.\thepopsub}\hspace{6pt}#1}\par\nopagebreak\vspace{4pt}
}

% ============================================================
% Boîtes pédagogiques — même recette : fond blanc, bordure épaisse colorée,
% ombre dure encre, badge plein. Titre optionnel [#1].
% ============================================================
\tcbset{popbase/.style={breakable,enhanced,colback=white,boxrule=2.3pt,arc=9pt,
  shadow={3pt}{-3pt}{0pt}{popink},left=14pt,right=14pt,top=11pt,bottom=11pt,before skip=11pt,after skip=15pt,
  fontupper=\popsemi\fontsize{10.6}{14.5}\selectfont\color{popink},
  fontlower=\popsemi\fontsize{10.6}{14.5}\selectfont\color{popink}}}
% Coche dessinée (le glyphe ✓ n'existe pas dans les polices du document)
\newcommand{\popcheck}[1]{\tikz[baseline=-0.5ex]\draw[#1,line width=1.6pt,line cap=round,line join=round](-0.13,0.0)--(-0.04,-0.09)--(0.13,0.1);}
\providecommand{\coche}{\popcheck{popgreen}}
\providecommand{\resultat}[1]{\par\smallskip\noindent\fcolorbox{popgreen}{popgreenL}{\parbox{\dimexpr\linewidth-2\fboxsep-2\fboxrule\relax}{\textbf{Résultat.} #1}}\par\smallskip}
\newcommand{\popboxhead}[3]{\poppill{#1}{#2}{white}\ifx\relax#3\relax\else\hspace{6pt}{\popxbold\fontsize{10.6}{12}\selectfont\color{popink}#3}\fi\par\vspace{7pt}}

\newtcolorbox{aretenir}[1][]{popbase,colframe=popgreen,before upper={\popboxhead{À retenir}{popgreen}{#1}}}
% Texte en chinois (document de lecture, dialogue, extrait) : cadre
% d'étude. Un titre optionnel (source, auteur) s'affiche dans l'en-tête.
\newtcolorbox{textebox}[1][]{popbase,colframe=popblue,before upper={\popboxhead{文本 · Texte}{popblue}{#1}},after upper={}}
% Marques ✓ / ✗ (forme correcte / fautive) souvent utilisées par les modèles
\providecommand{\cross}{\textcolor{poppink}{\ensuremath{\times}}}
\providecommand{\xmark}{\cross}\providecommand{\crossmark}{\cross}\providecommand{\cmark}{\ensuremath{\checkmark}}
% Ligne de réponse / justification à compléter (inventée par les modèles)
\providecommand{\justification}[1]{\par\noindent\textit{Justification~:} \dotfill\par}
% \textipa (package tipa non chargé) : la police ne couvre pas l'API -> simple passage
\providecommand{\textipa}[1]{#1}
% Soulignement (ulem non chargé) : \uline, \ul
\providecommand{\uline}[1]{\underline{#1}}\providecommand{\ul}[1]{\underline{#1}}
% \glossaire{\item ...} inventée par les modèles : liste de glossaire
\providecommand{\glossaire}[1]{\begin{itemize}#1\end{itemize}}
% \alert (beamer) inventée par les modèles : texte mis en évidence
\providecommand{\alert}[1]{\textbf{\textcolor{poppink}{#1}}}
% variantes ASCII d'ordinaux écrites en macro
\providecommand{\iere}{\textsuperscript{re}}\providecommand{\ieme}{\textsuperscript{e}}\providecommand{\ere}{\textsuperscript{re}}\providecommand{\eme}{\textsuperscript{e}}\providecommand{\er}{\textsuperscript{er}}\providecommand{\ie}{\textsuperscript{e}}
% Ordinaux français écrits en macro par les modèles : 1\ière, 3\ième
\newcommand{\ière}{\textsuperscript{re}}\newcommand{\ième}{\textsuperscript{e}}
% \exemple (inventée par les modèles) : étiquette « Exemple : »
\providecommand{\exemple}{\textbf{Exemple~:}\ }
% \square utilisable aussi hors mode math (cases à cocher)
\AtBeginDocument{\let\squareMath\square\renewcommand{\square}{\ensuremath{\squareMath}}}
% Mot ou phrase en chinois à l'intérieur d'un paragraphe français
\newcommand{\zh}[1]{\ifmmode\text{#1}\else{#1}\fi}
\let\eng\zh \let\esp\zh \let\all\zh % alias tolérés
% flèches d'intonation inventées par les modèles
\providecommand{\textsearrow}{}\renewcommand{\textsearrow}{\ensuremath{\searrow}}\providecommand{\textnearrow}{}\renewcommand{\textnearrow}{\ensuremath{\nearrow}}
% \fr{..} : mot français (inventé par les modèles)
\providecommand{\fr}[1]{\textit{#1}}
% zhbox : encadré de texte chinois inventé par les modèles
\newenvironment{zhbox}{\begin{quote}}{\end{quote}}
% \courssubsub : titre de niveau 3 inventé par les modèles
\providecommand{\courssubsub}[1]{\par\medskip\noindent\textbf{#1}\par\smallskip}
% \zhbf : texte chinois en gras inventé par les modèles
\providecommand{\zhbf}[1]{\textbf{#1}}
% \vs : « versus » inventé par les modèles
\providecommand{\vs}{\textit{vs}\ }
% \blank : case à compléter inventée par les modèles
\providecommand{\blank}[1][]{\underline{\hspace{1.6em}}}
\newtcolorbox{exemplebox}[1][]{popbase,colframe=poporange,before upper={\popboxhead{Exemple}{poporange}{#1}}}
\newtcolorbox{definition}[1][]{popbase,colframe=popblue,before upper={\popboxhead{Définition}{popblue}{#1}}}
\newtcolorbox{propriete}[1][]{popbase,colframe=poppurple,before upper={\popboxhead{Propriété}{poppurple}{#1}}}
\newtcolorbox{methode}[1][]{popbase,colframe=popgold,before upper={\popboxhead{Méthode}{popgold}{#1}}}
\newtcolorbox{attention}[1][]{popbase,colframe=poppink,before upper={\popboxhead{Attention}{poppink}{#1}}}
\newtcolorbox{experience}[1][]{popbase,colframe=popblue,colback=popblueL,before upper={\popboxhead{Expérience}{popblue}{#1}}}
% « Le savais-tu ? » : carte violette pleine (contexte, culture, Cameroun)
\newtcolorbox{savaistu}[1][]{popbase,colback=poppurple,colframe=popink,
  fontupper=\popsemi\fontsize{10.4}{14.2}\selectfont\color{white},
  before upper={\poppill{Le savais-tu\,?}{white}{poppurple}\ifx\relax#1\relax\else\hspace{6pt}{\popxbold\color{white}#1}\fi\par\vspace{7pt}}}
% Exercice résolu : énoncé en haut, \tcblower puis la solution sur fond crème
\newcounter{popexo}\newcounter{popexr}
\newtcolorbox{exoresolu}[1][]{popbase,colframe=poporange,colbacklower=popcream,
  segmentation style={popink,line width=1.2pt,dashed},
  before upper={\stepcounter{popexr}\popboxhead{Exercice résolu \thepopexr}{poporange}{#1}},
  before lower={\poppill{Solution}{popink}{popcream}\par\vspace{6pt}}}
% Exercice (non résolu en place) — la correction est dans la partie Corrigés
\newtcolorbox{exercice}[1][]{popbase,colframe=popink,
  before upper={\stepcounter{popexo}\popboxhead{Exercice \thepopexo}{popink}{#1}}}
% Corrigé
\newtcolorbox{corrige}[1][]{popbase,colframe=popgreen,colback=popgreenL,
  before upper={\popboxhead{Corrigé}{popgreen}{#1}}}
% Objectifs / prérequis / bilan
\newtcolorbox{objectifs}{popbase,colframe=popink,colback=poporangeL,
  before upper={\popboxhead{Objectifs de la leçon}{popink}{}}}
\newtcolorbox{prerequis}{popbase,colframe=popink,colback=popblueL,
  before upper={\popboxhead{Prérequis}{popblue}{}}}
\newtcolorbox{bilan}{popbase,colframe=popink,colback=white,boxrule=2.8pt,
  before upper={\popboxhead{Fiche bilan}{poporange}{L'essentiel à connaître pour l'examen}}}
% Figure encadrée avec légende
\newtcolorbox{popfigure}[1][]{enhanced,breakable=false,colback=white,colframe=popline,boxrule=1.4pt,arc=8pt,
  left=10pt,right=10pt,top=10pt,bottom=8pt,before skip=10pt,after skip=12pt,halign=center,
  lower separated=false,halign lower=center,
  fontlower=\popsemi\fontsize{9}{11.5}\selectfont\color{popmuted},
  #1}
\newcommand{\legende}[1]{\par\vspace{6pt}{\popsemi\fontsize{9}{11.5}\selectfont\color{popmuted}#1}}

% ============================================================
% Signature OnBuch+
% ============================================================
\newcommand{\poplogoinline}[1]{{\poplogo\fontsize{#1}{#1}\selectfont\color{popink}OnBuch\color{poporange}+}}
\newcommand{\signatureonbuch}{%
  \begin{tcolorbox}[enhanced,colback=popink,colframe=popink,boxrule=2.2pt,arc=10pt,
      drop shadow=poporange,left=14pt,right=14pt,top=10pt,bottom=10pt,before skip=0pt,after skip=0pt]
    \tikz[baseline=-0.7ex]\node[circle,fill=popgreen,draw=popcream,line width=1.3pt,minimum size=22pt,inner sep=0]{\popcheck{white}};%
    \hspace{9pt}\begin{minipage}[c]{0.62\linewidth}
      {\popxbold\fontsize{12}{13}\selectfont\color{popcream}Signé\ \textbullet\ {\poplogo OnBuch\color{poporange}+}}\par\vspace{2pt}
      {\raggedright\popsemi\fontsize{8}{10.5}\selectfont\color{popline}\MakeUppercase{Équipe pédagogique OnBuch+}\\\MakeUppercase{Conforme au programme officiel MINESEC}\par}
    \end{minipage}\hfill
    {\poplogo\fontsize{22}{22}\selectfont\color{popcream}OnBuch\color{poporange}+}
  \end{tcolorbox}%
}

% ============================================================
% Page de garde
% #1 titre · #2 matière · #3 niveau · #4 module · #5 n° leçon · #6 durée ·
% #7 description / objectifs (texte ou itemize)
% ============================================================
\newcommand{\pagegarde}[7]{%
  \thispagestyle{empty}
  \newgeometry{a4paper,margin=1.7cm,top=1.6cm,bottom=1.4cm}%
  \begin{tikzpicture}[remember picture,overlay]
    \fill[poporange!18] ([xshift=-3.2cm,yshift=-3.6cm]current page.north east) circle (4.6cm);
    \fill[poppurple!14] ([xshift=2.4cm,yshift=7.6cm]current page.south west) circle (3.8cm);
  \end{tikzpicture}%
  {\poplogo\fontsize{30}{30}\selectfont\color{popink}OnBuch\color{poporange}+}\hfill
  \raisebox{0.6ex}{\poppill{Cours signé \textbullet\ Premium}{popink}{popcream}}\par
  \vspace{1pt}\popsquiggle{poppurple}\par
  \vspace{8pt}
  \begin{tcolorbox}[enhanced,colback=poporange,colframe=popink,boxrule=2.8pt,arc=13pt,
      drop shadow=popink,left=18pt,right=18pt,top=12pt,bottom=13pt,after skip=10pt]
    \poppill{#2 \textbullet\ #3}{popink}{popcream}\hspace{5pt}\poppill{#4}{white}{popink}\par\vspace{8pt}
    {\popdisplay\fontsize{14}{14}\selectfont\color{popink}LEÇON #5}\par\vspace{4pt}
    {\raggedright\hyphenpenalty=10000\exhyphenpenalty=10000\popdisplay\fontsize{25}{27}\selectfont\color{popcream}\MakeUppercase{#1}\par}
    \vspace{8pt}
    {\popxbold\fontsize{9.5}{9.5}\selectfont\color{popink}\MakeUppercase{Durée conseillée : #6}}
  \end{tcolorbox}
  \begin{tcolorbox}[enhanced,colback=white,colframe=popink,boxrule=2.2pt,arc=10pt,
      drop shadow=popink,left=16pt,right=16pt,top=10pt,bottom=10pt,after skip=8pt]
    \poppill{Dans cette leçon}{poppurple}{white}\par\vspace{5pt}
    \popsemi\fontsize{9.3}{12.6}\selectfont\color{popink}#7
  \end{tcolorbox}
  \noindent\begin{minipage}[t]{0.487\linewidth}
    \begin{tcolorbox}[enhanced,colback=popgreen,colframe=popink,boxrule=2.2pt,arc=10pt,
        drop shadow=popink,left=11pt,right=11pt,top=10pt,bottom=10pt,height=3.1cm,valign=center]
      \tcbox[colback=white,colframe=popink,boxrule=1.3pt,arc=5pt,boxsep=0pt,nobeforeafter,
        left=3pt,right=3pt,top=3pt,bottom=3pt,valign=center]{\qrcode[height=1.9cm]{https://play.google.com/store/apps/details?id=com.luvvix.onbuch&hl=fr}}%
      \hspace{8pt}\begin{minipage}[c]{0.5\linewidth}
        {\popdisplay\fontsize{11}{12.5}\selectfont\color{white}\MakeUppercase{Télécharge OnBuch}}\par\vspace{4pt}
        {\raggedright\popsemi\fontsize{8.4}{11}\selectfont\color{white}Gratuit \textbullet\ Google Play\\Tuteur IA, annales, concours, résultats\par}
      \end{minipage}
    \end{tcolorbox}
  \end{minipage}\hfill
  \begin{minipage}[t]{0.487\linewidth}
    \begin{tcolorbox}[enhanced,colback=poppurple,colframe=popink,boxrule=2.2pt,arc=10pt,
        drop shadow=popink,left=11pt,right=11pt,top=10pt,bottom=10pt,height=3.1cm,valign=center]
      \tikz[baseline=-0.7ex]\node[circle,fill=white,draw=popink,line width=1.3pt,minimum size=1.6cm,inner sep=0]{\popdisplay\fontsize{24}{24}\selectfont\color{poppurple}+};%
      \hspace{8pt}\begin{minipage}[c]{0.6\linewidth}
        {\popdisplay\fontsize{11}{12.5}\selectfont\color{white}\MakeUppercase{Passe à OnBuch+}}\par\vspace{4pt}
        {\raggedright\popsemi\fontsize{8.4}{11}\selectfont\color{white}Tous les cours signés, Tuteur IA illimité, fiches PDF — onglet OnBuch+ de l'app\par}
      \end{minipage}
    \end{tcolorbox}
  \end{minipage}
  \vspace{8pt}
  \popcontenu
  \vfill
  \signatureonbuch
  \restoregeometry
  \newpage
}

% Rangée « Ce cours contient » (page de garde)
\newcommand{\poptile}[3]{% #1 couleur #2 gros texte #3 légende
  \begin{tcolorbox}[enhanced,colback=#1,colframe=popink,boxrule=2pt,arc=9pt,shadow={2.5pt}{-2.5pt}{0pt}{popink},
      left=6pt,right=6pt,top=9pt,bottom=9pt,halign=center,valign=center,height=1.75cm,width=\linewidth,nobeforeafter]
    {\popdisplay\fontsize{12.5}{13}\selectfont\color{white}\MakeUppercase{#2}}\par\vspace{3pt}
    {\centering\popsemi\fontsize{7.8}{9.5}\selectfont\color{white}#3\par}
  \end{tcolorbox}}
\newcommand{\popcontenu}{%
  \par\noindent{\popxboldls\fontsize{7.5}{7.5}\selectfont\color{popmuted}CE COURS SIGNÉ CONTIENT}\par\vspace{7pt}
  \noindent\begin{minipage}[t]{0.235\linewidth}\poptile{popblue}{Cours}{complet, illustré, pas à pas}\end{minipage}\hfill
  \begin{minipage}[t]{0.235\linewidth}\poptile{poporange}{Méthodes}{et exercices résolus}\end{minipage}\hfill
  \begin{minipage}[t]{0.235\linewidth}\poptile{poppink}{Exercices}{type examen + corrigés}\end{minipage}\hfill
  \begin{minipage}[t]{0.235\linewidth}\poptile{popgreen}{Fiche bilan}{l'essentiel à retenir}\end{minipage}\par}

% Sommaire
\newtcolorbox{sommaire}{enhanced,colback=white,colframe=popink,boxrule=2.2pt,arc=10pt,
  drop shadow=popink,left=18pt,right=18pt,top=13pt,bottom=13pt,before skip=6pt,after skip=20pt,
  before upper={\poppill{Sommaire}{poppurple}{white}\par\vspace{9pt}\popsemi\fontsize{10.6}{14.5}\selectfont\color{popink}}}

% Signature de fin de cours — poussée tout en bas de la dernière page
\newcommand{\findecours}{%
  \par\vfill\nopagebreak
  \begin{center}\popsquiggle{poporange}\end{center}\vspace{4pt}
  \signatureonbuch
}
\AtBeginDocument{\def\llbracket{\mathopen{[\mkern-3.5mu[}}\def\rrbracket{\mathclose{]\mkern-3.5mu]}}} % intervalles d'entiers (glyphe ⟦ absent de la police)
% Glyphes absents de Fira Math : redéfinis à partir de glyphes disponibles
\AtBeginDocument{%
  \def\setminus{\mathbin{\backslash}}\let\smallsetminus\setminus
  \def\vdots{\mathord{\vbox{\baselineskip3.2pt\lineskiplimit0pt\kern4pt\hbox{.}\hbox{.}\hbox{.}}}}%
  \def\ddots{\mathinner{\mkern1mu\raise6.5pt\hbox{.}\mkern2mu\raise3.6pt\hbox{.}\mkern2mu\raise0.7pt\hbox{.}\mkern1mu}}%
  \def\triangle{\mathord{\scalebox{0.9}{$\symup{\Delta}$}}}%
  \def\nabla{\mathord{\raisebox{\depth}{\scalebox{1}[-1]{$\symup{\Delta}$}}}}%
  \def\checkmark{\popcheck{popdark}}%
  \def\star{\mathbin{\ast}}\def\bigstar{\mathord{\ast}}%
  \def\diamond{\mathbin{\scalebox{0.6}{$\Diamond$}}}%
  \def\bigcup{\mathop{\mathchoice{\raisebox{-0.3em}{\scalebox{1.7}{$\cup$}}}{\scalebox{1.25}{$\cup$}}{\cup}{\cup}}\displaylimits}%
  \def\bigcap{\mathop{\mathchoice{\raisebox{-0.3em}{\scalebox{1.7}{$\cap$}}}{\scalebox{1.25}{$\cap$}}{\cap}{\cap}}\displaylimits}%
  \def\lesseqqgtr{\mathrel{\lessgtr}}\def\bigoplus{\mathop{\oplus}\displaylimits}%
}
\providecommand{\ssi}{si et seulement si} % abréviation fréquente
\AtBeginDocument{\providecommand{\wideparen}{\overparen}} % arc de cercle

%% FIGFIX-BEGIN v5 — correctifs globaux des figures (docs/onbuch-generation/tools/figfix)
\usepackage{adjustbox}\usepackage{etoolbox}\tcbuselibrary{raster}
% toute figure TikZ/pgfplots plus large que la ligne est réduite à la largeur disponible
\BeforeBeginEnvironment{tikzpicture}{\begin{adjustbox}{max width=\linewidth}}
\AfterEndEnvironment{tikzpicture}{\end{adjustbox}}
% légendes pgfplots lisibles par-dessus les courbes
\pgfplotsset{every axis/.append style={legend style={fill=white,fill opacity=0.92,text opacity=1,draw=black!15,inner sep=3pt}}}
% étiquettes d'arêtes (syntaxe edge["..."]) avec fond blanc
\tikzset{every edge quotes/.append style={fill=white,inner sep=1.2pt}}
%% FIGFIX-END
\begin{document}

\pagegarde{Leçon 33 — Éléments socioculturels : le marketing en Chine et au Cameroun}{Chinois}{Tle A}{Module 4 : Activités économiques et monde du travail}{ZH33}{2 h}{Cette leçon socioculturelle compare de façon factuelle et respectueuse les pratiques de marketing en Chine et au Cameroun : marchés traditionnels, publicité, réseaux sociaux et commerce en ligne. Elle fournit le vocabulaire chinois associé (caractères et pinyin) et prépare à un mini-exposé comparatif.
\vspace{4pt}
\begin{multicols}{2}\small
\begin{itemize}
\item À la fin de cette leçon, je sais définir le marketing et citer ses formes principales au Cameroun et en Chine
\item À la fin de cette leçon, je sais employer au moins 15 mots de vocabulaire chinois liés au marketing (caractères et pinyin)
\item À la fin de cette leçon, je sais comparer des pratiques commerciales camerounaises et chinoises avec les structures de comparaison (比, 和……一样)
\item À la fin de cette leçon, je sais décrire une publicité ou une pratique de vente en quelques phrases simples en chinois
\item À la fin de cette leçon, je sais exprimer un point de vue simple sur une pratique de marketing (我觉得……)
\item À la fin de cette leçon, je sais présenter un mini-exposé comparatif Cameroun/Chine de 8 à 10 phrases
\end{itemize}\end{multicols}}

\begin{sommaire}
\begin{multicols}{2}
\begin{enumerate}
\item Introduction : qu'est-ce que le marketing ?
\item Le marketing au Cameroun : faits et pratiques
\item Le marketing en Chine : faits et pratiques
\item Vocabulaire associé : glossaire et caractères
\item Comparer et donner son avis en chinois
\item Questions de réflexion et préparation du mini-exposé
\item Activité d'intégration
\item Méthodes et exercices résolus
\item Exercices
\item Corrigés des exercices
\item Fiche bilan
\end{enumerate}
\end{multicols}
\end{sommaire}

\begin{objectifs}
\begin{itemize}
\item À la fin de cette leçon, je sais définir le marketing et citer ses formes principales au Cameroun et en Chine
\item À la fin de cette leçon, je sais employer au moins 15 mots de vocabulaire chinois liés au marketing (caractères et pinyin)
\item À la fin de cette leçon, je sais comparer des pratiques commerciales camerounaises et chinoises avec les structures de comparaison (比, 和……一样)
\item À la fin de cette leçon, je sais décrire une publicité ou une pratique de vente en quelques phrases simples en chinois
\item À la fin de cette leçon, je sais exprimer un point de vue simple sur une pratique de marketing (我觉得……)
\item À la fin de cette leçon, je sais présenter un mini-exposé comparatif Cameroun/Chine de 8 à 10 phrases
\item À la fin de cette leçon, je sais adopter une attitude factuelle et respectueuse en comparant deux cultures
\end{itemize}
\end{objectifs}

\begin{prerequis}
\begin{itemize}
\item Structure de comparaison avec 比 (A 比 B + adjectif) vue en Première
\item Vocabulaire de l'argent, des prix et des achats (钱, 贵, 便宜, 买, 卖)
\item Expression de l'opinion avec 我觉得 / 我认为
\item Vocabulaire des moyens de communication (电视, 手机, 网络)
\item Nombres et pourcentages simples en chinois
\end{itemize}
\end{prerequis}

\newpage

\coursec{Introduction : qu'est-ce que le marketing ?}

\courssub{Situation d'accroche : vendre, ici et là-bas}

À Yaoundé, une vendeuse du marché Mokolo attire les clients en criant ses prix, tandis qu'à Douala, une entreprise de jus de fruits sponsorise un match de football. En Chine, des millions de consommateurs achètent en direct sur Douyin grâce au « livestreaming ». Comment vend-on ses produits au Cameroun ? Et en Chine ? Aujourd'hui, nous comparons les pratiques de marketing des deux pays et nous apprenons à en parler en chinois.

Ces trois scènes ont un point commun : dans chacune, quelqu'un cherche à \textbf{faire connaître} un produit et à \textbf{donner envie de l'acheter}. C'est exactement ce que recouvre la notion de marketing. La problématique de cette leçon est donc la suivante : \emph{qu'est-ce que le marketing, et comment s'exprime-t-il au Cameroun et en Chine ?}

\courssub{Observation : un texte pour découvrir la notion}

Lisons d'abord un court texte chinois qui donne une définition simple et vivante du marketing.

\begin{textebox}[Le marketing, en une phrase]
营销就是让别人知道并且想买你的产品。\\
\emph{Yíngxiāo jiùshì ràng biérén zhīdào bìngqiě xiǎng mǎi nǐ de chǎnpǐn.}\\
« Le marketing, c'est faire connaître votre produit et donner envie de l'acheter. »

比如说，一个卖水果的人大声喊价格，这就是一种很简单的营销。\\
\emph{Bǐrú shuō, yí ge mài shuǐguǒ de rén dàshēng hǎn jiàgé, zhè jiùshì yì zhǒng hěn jiǎndān de yíngxiāo.}\\
« Par exemple, un vendeur de fruits qui crie ses prix à voix haute pratique déjà une forme très simple de marketing. »
\end{textebox}

\textbf{Glossaire du texte}

\begin{center}
\begin{tabularx}{\linewidth}{|X|X|X|X|}
\hline
\rowcolor{popblueL} Caractères & Pinyin & Nature & Traduction \\
\hline
营销 & yíngxiāo & nom & marketing, commercialisation \\
\hline
让 & ràng & verbe & faire, laisser, amener (quelqu'un à) \\
\hline
别人 & biérén & nom & les autres, autrui \\
\hline
知道 & zhīdào & verbe & savoir, connaître \\
\hline
并且 & bìngqiě & conjonction & et de plus, et aussi \\
\hline
想买 & xiǎng mǎi & verbe + verbe & vouloir acheter \\
\hline
产品 & chǎnpǐn & nom & produit \\
\hline
比如说 & bǐrú shuō & expression & par exemple \\
\hline
大声 & dàshēng & adverbe & à voix haute \\
\hline
喊 & hǎn & verbe & crier \\
\hline
\end{tabularx}
\end{center}

\textbf{Questions de compréhension}

\begin{enumerate}
\item Selon le texte, quels sont les deux objectifs du marketing ?
\item Quel exemple de marketing simple le texte donne-t-il ?
\item Connaissez-vous des vendeurs qui font la même chose dans votre ville ?
\end{enumerate}

\courssub{Brainstorming : où voit-on de la publicité au Cameroun ?}

Avant de définir le mot chinois, observons notre environnement. Où voit-on de la publicité à Yaoundé, Douala ou Bafoussam ?

\begin{itemize}
\item à la \textbf{radio} : les spots publicitaires entre les émissions ;
\item sur les \textbf{affiches} et les panneaux au bord des routes ;
\item sur les \textbf{t-shirts} et les casquettes offerts par les entreprises de téléphonie mobile ou de boissons ;
\item aux \textbf{carrefours} et dans les marchés : haut-parleurs, crieurs, dégustations ;
\item sur les \textbf{réseaux sociaux} : Facebook, WhatsApp, TikTok, où les vendeurs montrent leurs produits en photo ou en vidéo.
\end{itemize}

\begin{experience}[Brainstorming en classe]
En groupes de quatre, listez en deux minutes tous les endroits où vous avez vu de la publicité cette semaine. Puis essayez de dire en chinois : \zh{我在……看到了广告。} \emph{Wǒ zài … kàndào le guǎnggào.} « J'ai vu de la publicité à/sur … » (par exemple : \zh{我在市场上看到了广告。} « J'ai vu de la publicité au marché. »)
\end{experience}

\courssub{Les mots clés : 营销 yíngxiāo et 广告 guǎnggào}

\begin{definition}[营销 yíngxiāo]
\zh{营销} \emph{yíngxiāo} est un nom qui signifie « marketing, commercialisation ». Il est composé de deux caractères :
\begin{itemize}
\item \zh{营} \emph{yíng} : gérer, exploiter, faire fonctionner (on le retrouve dans \zh{经营} \emph{jīngyíng}, « gérer une entreprise ») ;
\item \zh{销} \emph{xiāo} : vendre (on le retrouve dans \zh{销售} \emph{xiāoshòu}, « vente »).
\end{itemize}
Le mot chinois dit donc littéralement « gérer-vendre » : le marketing, c'est l'art d'organiser la vente.
\end{definition}

\begin{definition}[广告 guǎnggào]
\zh{广告} \emph{guǎnggào} est un nom qui signifie « publicité, annonce publicitaire ». Il est composé de \zh{广} \emph{guǎng} (« large, vaste ») et \zh{告} \emph{gào} (« annoncer, informer ») : littéralement « annoncer largement », c'est-à-dire diffuser un message au plus grand nombre.
\end{definition}

\begin{propriete}[Marketing et publicité : deux notions à distinguer]
Le \textbf{marketing} \zh{营销} est la \textbf{stratégie globale} : étudier le marché, connaître les clients, fixer les prix, choisir les réseaux de vente. La \textbf{publicité} \zh{广告} n'est qu'\textbf{un outil} du marketing, parmi d'autres. Schéma en français :

marketing = étude du marché + produit + prix + publicité + réseaux de vente.

En chinois, on dira : \zh{广告是营销的一部分。} \emph{Guǎnggào shì yíngxiāo de yí bùfen.} « La publicité est une partie du marketing. »
\end{propriete}

\begin{attention}[Ne pas confondre 广告 et 市场]
Les francophones qui traduisent « marketing » mot à mot tombent souvent sur deux pièges :
\begin{itemize}
\item \zh{广告} \emph{guǎnggào} = la \textbf{publicité} (le message, l'annonce), pas le marketing ;
\item \zh{市场} \emph{shìchǎng} = le \textbf{marché} (le lieu ou l'espace économique où l'on achète et vend), pas le marketing.
\end{itemize}
On ne dit donc jamais \zh{市场学} pour « marketing » : le mot correct est \zh{营销}. Retenez le trio : \zh{营销} (stratégie) > \zh{广告} (outil) ; \zh{市场} (lieu de vente).
\end{attention}

\courssub{Le vocabulaire du marketing : carte mentale}

Autour du mot central \zh{营销}, voici les cinq mots essentiels de cette leçon.

\begin{popfigure}
\begin{center}\tcbox[colback=popink,colframe=popink,coltext=white,arc=9pt,boxsep=4pt,left=6pt,right=6pt]{\bfseries\small 营销 yíngxiāo}\end{center}
\vspace{-2pt}
\begin{tcbraster}[raster columns=3,raster equal height=rows,raster column skip=6pt,raster row skip=6pt,enhanced,blankest]
\begin{tcolorbox}[enhanced,colback=poporange,colframe=poporange,coltext=white,boxrule=0.9pt,arc=3pt,left=4pt,right=4pt,top=3pt,bottom=3pt,boxsep=1pt,halign=left]\bfseries\scriptsize 广告 guǎnggào publicité\end{tcolorbox}
\begin{tcolorbox}[enhanced,colback=popgreen,colframe=popgreen,coltext=white,boxrule=0.9pt,arc=3pt,left=4pt,right=4pt,top=3pt,bottom=3pt,boxsep=1pt,halign=left]\bfseries\scriptsize 市场 shìchǎng marché\end{tcolorbox}
\begin{tcolorbox}[enhanced,colback=poppurple,colframe=poppurple,coltext=white,boxrule=0.9pt,arc=3pt,left=4pt,right=4pt,top=3pt,bottom=3pt,boxsep=1pt,halign=left]\bfseries\scriptsize 顾客 gùkè client\end{tcolorbox}
\begin{tcolorbox}[enhanced,colback=poppink,colframe=poppink,coltext=white,boxrule=0.9pt,arc=3pt,left=4pt,right=4pt,top=3pt,bottom=3pt,boxsep=1pt,halign=left]\bfseries\scriptsize 价格 jiàgé prix\end{tcolorbox}
\begin{tcolorbox}[enhanced,colback=popgold,colframe=popgold,coltext=white,boxrule=0.9pt,arc=3pt,left=4pt,right=4pt,top=3pt,bottom=3pt,boxsep=1pt,halign=left]\bfseries\scriptsize 网络 wǎngluò réseau, internet\end{tcolorbox}
\end{tcbraster}

\legende{Carte mentale lexicale : le mot 营销 et ses cinq compagnons}
\end{popfigure}

\begin{center}
\begin{tabularx}{\linewidth}{|X|X|X|X|}
\hline
\rowcolor{popblueL} Caractères & Pinyin & Nature & Traduction \\
\hline
营销 & yíngxiāo & nom & marketing, commercialisation \\
\hline
广告 & guǎnggào & nom & publicité, annonce \\
\hline
市场 & shìchǎng & nom & marché \\
\hline
顾客 & gùkè & nom & client \\
\hline
价格 & jiàgé & nom & prix \\
\hline
网络 & wǎngluò & nom & réseau, internet \\
\hline
产品 & chǎnpǐn & nom & produit \\
\hline
卖 & mài & verbe & vendre \\
\hline
买 & mǎi & verbe & acheter \\
\hline
公司 & gōngsī & nom & entreprise, société \\
\hline
\end{tabularx}
\end{center}

\begin{aretenir}[Tableau comparatif : français et chinois]
\begin{center}
\begin{tabularx}{\linewidth}{|X|X|X|}
\hline
\rowcolor{popblueL} Notion française & Chinois & Piège à éviter \\
\hline
le marketing (stratégie) & 营销 yíngxiāo & ne pas dire 市场 \\
\hline
la publicité (un outil) & 广告 guǎnggào & ce n'est qu'une partie du marketing \\
\hline
le marché (lieu, espace) & 市场 shìchǎng & ne veut pas dire « marketing » \\
\hline
le client & 顾客 gùkè & plus poli que 客人 dans le commerce \\
\hline
le prix & 价格 jiàgé & ne pas confondre avec 钱 « argent » \\
\hline
\end{tabularx}
\end{center}
\end{aretenir}

\courssub{Du produit au client : le chemin du marketing}

Le marketing est le pont entre le produit et le client. Observons le schéma suivant.

\begin{popfigure}
\begin{tikzpicture}[node distance=1.6cm, >=Stealth]
\node[draw=popblue, fill=popblueL, rounded corners, thick, align=center, minimum width=2.6cm, minimum height=1.1cm] (prod) {产品 chǎnpǐn\\ produit};
\node[draw=poporange, fill=poporangeL, rounded corners, thick, align=center, minimum width=2.6cm, minimum height=1.1cm, right=of prod] (mkt) {营销 yíngxiāo\\ marketing};
\node[draw=popgreen, fill=popgreenL, rounded corners, thick, align=center, minimum width=2.6cm, minimum height=1.1cm, right=of mkt] (cli) {顾客 gùkè\\ client};
\draw[->, very thick, popdark] (prod) -- (mkt);
\draw[->, very thick, popdark] (mkt) -- (cli);
\end{tikzpicture}
\legende{Le marketing relie le produit au client}
\end{popfigure}

En chinois, on peut décrire ce schéma avec la structure \zh{通过……，……} « grâce à …, … » :

\zh{公司通过营销，把产品卖给顾客。} \emph{Gōngsī tōngguò yíngxiāo, bǎ chǎnpǐn mài gěi gùkè.} « Grâce au marketing, l'entreprise vend ses produits aux clients. »

Notez la structure en \zh{把} : \zh{把 + objet + verbe + complément} permet de mettre l'objet (\zh{产品}) avant le verbe (\zh{卖给}).

\courssub{Exercice résolu et entraînement}

\begin{exoresolu}[Traduire en chinois]
Énoncé : traduisez la phrase « La publicité fait connaître le produit aux clients. »
\tcblower
Solution détaillée :
\begin{enumerate}
\item Sujet : « la publicité » = \zh{广告}.
\item « faire connaître » = \zh{让……知道} (structure causative : 让 + personne + verbe).
\item « aux clients » = \zh{让顾客} ; « le produit » = \zh{产品}, placé après le verbe.
\end{enumerate}
Traduction : \zh{广告让顾客知道产品。} \emph{Guǎnggào ràng gùkè zhīdào chǎnpǐn.}

Erreur à éviter : ne pas traduire « faire » par \zh{做} ; la structure causative française « faire + infinitif » se rend en chinois par \zh{让} + personne + verbe.
\end{exoresolu}

\begin{exercice}[À vous de jouer]
\begin{enumerate}
\item Traduisez en français : \zh{营销不只是广告。} \emph{Yíngxiāo bù zhǐshì guǎnggào.}
\item Complétez avec le bon mot (营销 / 广告 / 市场) : \zh{我在电视上看到了一个手机……。} « J'ai vu une … de téléphone à la télévision. »
\item Traduisez en chinois : « Le marketing aide l'entreprise à vendre ses produits. » (utilisez \zh{帮助} « aider »).
\end{enumerate}
\end{exercice}

\begin{corrige}[Exercice « À vous de jouer »]
\begin{enumerate}
\item « Le marketing, ce n'est pas seulement la publicité. » (\zh{不只是} = « n'est pas seulement ».)
\item \zh{广告} : on voit une \emph{publicité} à la télévision, pas un marché ni une stratégie.
\item \zh{营销帮助公司卖产品。} \emph{Yíngxiāo bāngzhù gōngsī mài chǎnpǐn.}
\end{enumerate}
\end{corrige}

\begin{savaistu}[La Chine, géant du commerce en ligne]
La Chine est le premier marché mondial du commerce en ligne. Des plateformes comme \zh{淘宝} \emph{Táobǎo} (groupe Alibaba) ou JD.com permettent d'acheter presque tout, du téléphone au panier de légumes, livré à domicile. Une pratique typiquement chinoise est le \zh{直播带货} \emph{zhíbō dàihuò}, le « livestreaming shopping » : un animateur présente des produits en direct vidéo et les spectateurs achètent en un clic. Au Cameroun, les ventes par WhatsApp et Facebook se développent aussi rapidement : les deux pays montrent que le marketing suit les habitudes des consommateurs.
\end{savaistu}

\courssub{Annonce du plan de la leçon et de la tâche finale}

Maintenant que la notion de marketing est claire, voici le chemin que nous allons suivre :

\begin{enumerate}
\item \textbf{Introduction} : qu'est-ce que le marketing ? (notion acquise !)
\item Le marketing \textbf{au Cameroun} : faits et pratiques.
\item Le marketing \textbf{en Chine} : faits et pratiques.
\item Le \textbf{vocabulaire} associé : glossaire et caractères.
\item \textbf{Comparer} et donner son avis en chinois.
\item Questions de réflexion et préparation du \textbf{mini-exposé}.
\end{enumerate}

\begin{methode}[Tâche finale : le mini-exposé comparatif]
À la fin de la leçon, vous présenterez un mini-exposé de deux minutes comparant le marketing au Cameroun et en Chine. Préparez-vous dès maintenant :
\begin{enumerate}
\item Notez deux exemples de marketing observés au Cameroun (marché, radio, réseaux sociaux).
\item Retenez deux pratiques chinoises (commerce en ligne, livestreaming).
\item Mémorisez les cinq mots de la carte mentale : \zh{营销、广告、市场、顾客、价格}.
\item Préparez une phrase d'opinion : \zh{我觉得……} \emph{Wǒ juéde …} « Je pense que … ».
\end{enumerate}
\end{methode}

\coursec{Le marketing au Cameroun : faits et pratiques}

Dans la section précédente, nous avons défini le \cle{marketing} (\zh{市场营销} shìchǎng yíngxiāo) comme l'ensemble des moyens utilisés par une entreprise pour faire connaître et vendre ses produits. Nous allons maintenant observer comment ces moyens se concrétisent au Cameroun, notre terrain d'observation quotidien : les marchés de Yaoundé et de Douala, la radio, les affiches, le parrainage sportif et, de plus en plus, le téléphone portable. L'objectif est double : décrire ces réalités économiques de façon factuelle et apprendre à en parler en chinois.

\courssub{Texte support : le marketing au quotidien au Cameroun}

\begin{textebox}[Le marketing au Cameroun]
在喀麦隆，很多人在市场买东西，也喜欢在手机上看广告。\\
\emph{Zài Kāmàilóng, hěn duō rén zài shìchǎng mǎi dōngxi, yě xǐhuan zài shǒujī shang kàn guǎnggào.}\\
« Au Cameroun, beaucoup de gens achètent au marché et aiment aussi regarder des publicités sur leur téléphone. »\\[4pt]
雅温得和杜阿拉的大市场很有名。在市场里，卖东西的人常常大声介绍自己的商品，价格也可以谈。\\
\emph{Yǎwēndé hé Dù'ālā de dà shìchǎng hěn yǒumíng. Zài shìchǎng lǐ, mài dōngxi de rén chángcháng dà shēng jièshào zìjǐ de shāngpǐn, jiàgé yě kěyǐ tán.}\\
« Les grands marchés de Yaoundé et de Douala sont très connus. Au marché, les vendeurs présentent souvent leurs marchandises à voix haute, et on peut aussi discuter le prix. »\\[4pt]
很多公司在广播和电视上做广告，也在路边放很大的广告画。一些电话公司和啤酒公司还给足球比赛和节日活动出钱，因为这样可以让更多人知道他们的名字。\\
\emph{Hěn duō gōngsī zài guǎngbō hé diànshì shang zuò guǎnggào, yě zài lù biān fàng hěn dà de guǎnggào huà. Yìxiē diànhuà gōngsī hé píjiǔ gōngsī hái gěi zúqiú bǐsài hé jiérì huódòng chū qián, yīnwèi zhèyàng kěyǐ ràng gèng duō rén zhīdào tāmen de míngzi.}\\
« Beaucoup d'entreprises font de la publicité à la radio et à la télévision, et placent aussi de grandes affiches au bord des routes. Certaines compagnies de téléphone et de brasserie financent aussi des matchs de football et des festivals, car cela permet à plus de gens de connaître leur nom. »\\[4pt]
现在，很多年轻人用脸书和微信卖东西，网上商店也越来越多。可是，朋友介绍和口耳相传还是很重要的营销方式。\\
\emph{Xiànzài, hěn duō niánqīngrén yòng liǎnshū hé Wēixìn mài dōngxi, wǎngshang shāngdiàn yě yuè lái yuè duō. Kěshì, péngyou jièshào hé kǒu'ěr xiāng chuán háishi hěn zhòngyào de yíngxiāo fāngshì.}\\
« Aujourd'hui, beaucoup de jeunes vendent sur Facebook et WeChat, et les boutiques en ligne sont de plus en plus nombreuses. Pourtant, la recommandation entre amis et le bouche-à-oreille restent des moyens de marketing très importants. »
\end{textebox}

\textbf{Glossaire du texte}

\begin{center}
\begin{tabularx}{\linewidth}{|X|X|X|X|}
\hline
\rowcolor{popblueL} Caractères & Pinyin & Nature & Traduction \\
\hline
市场 & shìchǎng & nom & marché \\
\hline
广告 & guǎnggào & nom & publicité, annonce \\
\hline
商品 & shāngpǐn & nom & marchandise, produit \\
\hline
价格 & jiàgé & nom & prix \\
\hline
广播 & guǎngbō & nom & radio (diffusion) \\
\hline
电视 & diànshì & nom & télévision \\
\hline
公司 & gōngsī & nom & entreprise, société \\
\hline
比赛 & bǐsài & nom & match, compétition \\
\hline
出钱 & chū qián & verbe & financer, mettre de l'argent \\
\hline
口耳相传 & kǒu'ěr xiāng chuán & expression & bouche-à-oreille \\
\hline
方式 & fāngshì & nom & manière, moyen \\
\hline
\end{tabularx}
\end{center}

\textbf{Questions de compréhension}
\begin{enumerate}
\item Où les Camerounais achètent-ils souvent leurs produits ?
\item Que font les compagnies de téléphone et les brasseries pour se faire connaître ?
\item Quels moyens de marketing les jeunes utilisent-ils aujourd'hui ?
\end{enumerate}

\courssub{Les marchés traditionnels : cœur du commerce}

\begin{definition}[Le marché traditionnel]
Au Cameroun, le \cle{marché traditionnel} (\zh{传统市场} chuántǒng shìchǎng) est un lieu de vente directe : le vendeur et l'acheteur se rencontrent face à face. Les grands marchés comme Mokolo et le Marché Central à Yaoundé, ou le Marché Mboppi à Douala, concentrent des milliers de vendeurs. On y pratique la \cle{criée} (le vendeur annonce ses produits à voix haute) et la \cle{négociation des prix} (\zh{谈价格} tán jiàgé, « discuter le prix »).
\end{definition}

Dans ce contexte, le « marketing » n'est pas une affiche : c'est la voix du vendeur, la présentation des étals, la qualité visible des produits et la relation de confiance. Le prix affiché n'est souvent qu'un point de départ de la discussion.

\begin{exemplebox}[Décrire le marché en chinois]
这个市场很有名。\\
\emph{Zhège shìchǎng hěn yǒumíng.}\\
« Ce marché est très connu. »\\[4pt]
市场里的人很多。\\
\emph{Shìchǎng lǐ de rén hěn duō.}\\
« Il y a beaucoup de monde dans le marché. »\\[4pt]
这个价格可以谈吗？\\
\emph{Zhège jiàgé kěyǐ tán ma ?}\\
« Ce prix est-il négociable ? »
\end{exemplebox}

\courssub{La publicité par les médias classiques}

La \cle{radio} (\zh{广播} guǎngbō) est un canal publicitaire majeur au Cameroun : elle est très écoutée, y compris dans les zones rurales où la télévision et l'internet sont moins accessibles. La \cle{télévision} (\zh{电视} diànshì), avec la chaîne publique CRTV et les chaînes privées, diffuse des spots publicitaires pour les grandes marques. Enfin, les \cle{affiches et bâches publicitaires} (\zh{广告画} guǎnggào huà) bordent les grandes artères de Douala et de Yaoundé.

\begin{exemplebox}[Parler de publicité à la télévision]
他们在电视上做广告。\\
\emph{Tāmen zài diànshì shang zuò guǎnggào.}\\
« Ils font de la publicité à la télévision. »\\[4pt]
这家公司在广播里做广告。\\
\emph{Zhè jiā gōngsī zài guǎngbō lǐ zuò guǎnggào.}\\
« Cette entreprise fait de la publicité à la radio. »
\end{exemplebox}

\begin{savaistu}[La radio, média de proximité]
Au Cameroun, la radio reste le média le plus touché en zone rurale : un simple poste à piles suffit, et les émissions en langues locales touchent des publics que ni la télévision ni l'internet n'atteignent. C'est pourquoi de nombreuses campagnes publicitaires — et aussi des messages d'intérêt public — passent d'abord par la radio. En Chine, à l'inverse, la publicité passe aujourd'hui massivement par les applications mobiles comme WeChat.
\end{savaistu}

\courssub{Le marketing par le parrainage}

Le \cle{parrainage} ou \cle{sponsoring} (\zh{赞助} zànzhù) consiste pour une entreprise à financer un événement en échange de visibilité. Au Cameroun, les entreprises de téléphonie mobile et les brasseries parrainent régulièrement le football — sport le plus populaire du pays — ainsi que des festivals culturels et des concerts. Leur nom apparaît alors sur les maillots, les banderoles et les affiches.

\begin{exemplebox}[Exprimer le parrainage]
这家公司赞助足球比赛。\\
\emph{Zhè jiā gōngsī zànzhù zúqiú bǐsài.}\\
« Cette entreprise sponsorise des matchs de football. »\\[4pt]
他们出钱，因为这样可以让更多人知道他们的名字。\\
\emph{Tāmen chū qián, yīnwèi zhèyàng kěyǐ ràng gèng duō rén zhīdào tāmen de míngzi.}\\
« Ils financent, car cela permet à plus de gens de connaître leur nom. »
\end{exemplebox}

\courssub{L'essor du marketing digital}

Depuis quelques années, le \cle{marketing digital} (\zh{数字营销} shùzì yíngxiāo) se développe rapidement au Cameroun, porté par le téléphone mobile :

\begin{itemize}
\item \textbf{Facebook} (\zh{脸书} liǎnshū) : les petits commerçants y présentent leurs produits avec des photos ;
\item \textbf{WhatsApp Business} : catalogue de produits et prise de commandes par message ;
\item les \textbf{influenceurs} (\zh{网红} wǎnghóng, litt. « célébrités d'internet ») camerounais recommandent des marques à leurs abonnés ;
\item la \textbf{vente en ligne} (\zh{网上购物} wǎngshang gòuwù) avec des plateformes comme Jumia, qui livrent à domicile dans les grandes villes.
\end{itemize}

\begin{exemplebox}[Le commerce en ligne]
很多年轻人在网上卖东西。\\
\emph{Hěn duō niánqīngrén zài wǎngshang mài dōngxi.}\\
« Beaucoup de jeunes vendent des produits en ligne. »\\[4pt]
网上商店越来越多。\\
\emph{Wǎngshang shāngdiàn yuè lái yuè duō.}\\
« Les boutiques en ligne sont de plus en plus nombreuses. »
\end{exemplebox}

\courssub{Le bouche-à-oreille et les réseaux relationnels}

Enfin, une forme de promotion très ancienne reste centrale : le \cle{bouche-à-oreille} (\zh{口耳相传} kǒu'ěr xiāng chuán, « transmis de bouche en oreille »). Les \cle{tontines} et les associations de quartier ou d'anciens élèves fonctionnent comme des réseaux de confiance : on y recommande un bon tailleur, un fournisseur sérieux, une vendeuse honnête. La recommandation personnelle (\zh{朋友介绍} péngyou jièshào) y a souvent plus de poids qu'une affiche.

\begin{propriete}[La structure descriptive : lieu + 很多人 + verbe + complément]
Pour décrire une pratique locale, le chinois utilise une structure simple :\\[2pt]
\textbf{Complément de lieu (在 + lieu) + 很多人 + verbe + complément}\\[2pt]
\zh{在市场，很多人买东西。} \emph{Zài shìchǎng, hěn duō rén mǎi dōngxi.} « Au marché, beaucoup de gens achètent des choses. »\\[2pt]
\zh{在网上，很多年轻人做广告。} \emph{Zài wǎngshang, hěn duō niánqīngrén zuò guǎnggào.} « Sur internet, beaucoup de jeunes font de la publicité. »\\[2pt]
Remarque : le complément de lieu avec \zh{在} se place avant le verbe, contrairement au français où « au marché » vient souvent après (« les gens achètent au marché »).
\end{propriete}

\begin{methode}[Décrire une pratique de marketing en quatre étapes]
\begin{enumerate}
\item \textbf{Identifier le lieu} : 在市场 (au marché), 在电视上 (à la télévision), 在网上 (en ligne).
\item \textbf{Dire qui agit} : 很多人 (beaucoup de gens), 公司 (les entreprises), 年轻人 (les jeunes).
\item \textbf{Choisir le verbe} : 买 (acheter), 卖 (vendre), 做广告 (faire de la publicité), 赞助 (sponsoriser).
\item \textbf{Ajouter le complément} : 东西 (des choses), 商品 (des produits), 足球比赛 (des matchs de football).
\end{enumerate}
Résultat : \zh{在电视上，很多公司做广告。} « À la télévision, beaucoup d'entreprises font de la publicité. »
\end{methode}

\begin{attention}[Erreurs fréquentes des francophones]
\begin{itemize}
\item \textbf{Place du complément de lieu} : ne dites pas \zh{很多人买东西在市场} (calque du français « les gens achètent au marché »). Dites \zh{很多人在市场买东西。}
\item \textbf{做广告} est un verbe à complément inséparable : on dit \zh{在电视上做广告}, jamais \zh{广告做}.
\item \textbf{有名} signifie « célèbre, connu » ; ne le confondez pas avec \zh{名字} (míngzi, « nom »).
\item Attention aux tons de \zh{买} (mǎi, 3\up{e} ton, « acheter ») et \zh{卖} (mài, 4\up{e} ton, « vendre ») : un seul ton les oppose !
\end{itemize}
\end{attention}

\begin{popfigure}
\begin{tikzpicture}
\node[draw=popblue, fill=popblueL, rounded corners, thick, minimum width=3.6cm, align=center, font=\bfseries] (titre) at (0,0) {Le marketing\\ au Cameroun};
\node[draw=poporange, fill=poporangeL, rounded corners, thick, minimum width=4.6cm, align=center] (trad) at (-4.2,-2.2) {\textbf{Marketing traditionnel}\\ 市场 marchés (Mokolo, Mboppi)\\ 广播 radio, 电视 télévision\\ 广告画 affiches, bâches\\ 赞助 parrainage (football)\\ 口耳相传 bouche-à-oreille};
\node[draw=popgreen, fill=popgreenL, rounded corners, thick, minimum width=4.6cm, align=center] (dig) at (4.2,-2.2) {\textbf{Marketing digital}\\ 脸书 Facebook\\ WhatsApp Business\\ 网红 influenceurs\\ 网上购物 vente en ligne (Jumia)\\ 手机 téléphone mobile};
\draw[-{Stealth}, thick, poporange] (titre) -- (trad);
\draw[-{Stealth}, thick, popgreen] (titre) -- (dig);
\end{tikzpicture}
\legende{Carte mentale : les deux grandes familles du marketing au Cameroun}
\end{popfigure}

\begin{center}
\begin{tabularx}{\linewidth}{|X|X|X|}
\hline
\rowcolor{popblueL} Pratique & Exemple en chinois + pinyin & Traduction \\
\hline
Marché traditionnel & 这个市场很有名。Zhège shìchǎng hěn yǒumíng. & Ce marché est très connu. \\
\hline
Négociation & 价格可以谈。Jiàgé kěyǐ tán. & On peut discuter le prix. \\
\hline
Publicité TV & 他们在电视上做广告。Tāmen zài diànshì shang zuò guǎnggào. & Ils font de la publicité à la télévision. \\
\hline
Parrainage & 公司赞助足球比赛。Gōngsī zànzhù zúqiú bǐsài. & L'entreprise sponsorise des matchs de football. \\
\hline
Vente en ligne & 年轻人在网上卖东西。Niánqīngrén zài wǎngshang mài dōngxi. & Les jeunes vendent en ligne. \\
\hline
Bouche-à-oreille & 朋友介绍很重要。Péngyou jièshào hěn zhòngyào. & La recommandation entre amis est très importante. \\
\hline
\end{tabularx}
\end{center}

\begin{exoresolu}[Traduire en chinois]
Énoncé : Traduisez la phrase suivante en chinois : « Au marché, beaucoup de gens achètent des choses et discutent les prix. »
\tcblower
Solution détaillée :
\begin{enumerate}
\item Complément de lieu d'abord : \zh{在市场} (zài shìchǎng).
\item Sujet : \zh{很多人} (hěn duō rén).
\item Premier verbe + complément : \zh{买东西} (mǎi dōngxi).
\item Deuxième action reliée par \zh{也} ou par la coordination : \zh{谈价格} (tán jiàgé).
\end{enumerate}
Phrase complète : \zh{在市场，很多人买东西，也谈价格。}\\
\emph{Zài shìchǎng, hěn duō rén mǎi dōngxi, yě tán jiàgé.}\\
Piège évité : le complément de lieu \zh{在市场} est bien placé avant le verbe, et non à la fin comme en français.
\end{exoresolu}

\begin{exercice}[À vous de décrire]
En vous inspirant de la méthode en quatre étapes, écrivez deux phrases en chinois :
\begin{enumerate}
\item une phrase sur la publicité à la radio au Cameroun ;
\item une phrase sur les jeunes qui vendent en ligne.
\end{enumerate}
\end{exercice}

\begin{corrige}[Exercice « À vous de décrire »]
Propositions de correction :
\begin{enumerate}
\item \zh{在广播里，很多公司做广告。} \emph{Zài guǎngbō lǐ, hěn duō gōngsī zuò guǎnggào.} « À la radio, beaucoup d'entreprises font de la publicité. »
\item \zh{在网上，很多年轻人卖东西。} \emph{Zài wǎngshang, hěn duō niánqīngrén mài dōngxi.} « En ligne, beaucoup de jeunes vendent des produits. »
\end{enumerate}
\end{corrige}

\begin{aretenir}[L'essentiel de la section]
\begin{itemize}
\item Le marketing camerounais repose sur cinq piliers : les marchés traditionnels (vente directe, négociation), les médias classiques (radio, télévision, affiches), le parrainage sportif et culturel, le marketing digital en plein essor, et le bouche-à-oreille.
\item Ces pratiques sont des réalités économiques locales : on les décrit de façon factuelle, sans jugement de valeur.
\item Structure clé : \textbf{在 + lieu + 很多人 + verbe + complément}.
\item Vocabulaire clé : \zh{市场} marché, \zh{广告} publicité, \zh{赞助} sponsoriser, \zh{口耳相传} bouche-à-oreille.
\end{itemize}
\end{aretenir}

\coursec{Leçon 33 — Éléments socioculturels : le marketing en Chine et au Cameroun}

\courssub{Introduction : qu'est-ce que le marketing ?}

\begin{definition}[营销  — Marketing]
Le terme chinois pour « marketing » est \zh{营销} \emph{yíngxiāo}. Il est composé de \zh{营} \emph{yíng} (« gérer, exploiter, camp ») et \zh{销} \emph{xiāo} (« vendre, écouler »). Il désigne l'ensemble des techniques et actions visant à \textbf{comprendre les besoins des consommateurs} (\zh{消费者} \emph{xiāofèizhě}), à \textbf{créer de la valeur} et à \textbf{conclure des échanges} sur un \zh{市场} \emph{shìchǎng} (marché).
\end{definition}

\begin{exemplebox}[Mots-clés du marketing]
\zh{市场调研} \emph{shìchǎng diyán} — étude de marché \\
\zh{品牌} \emph{pǐnpái} — marque \\
\zh{广告} \emph{guǎnggào} — publicité \\
\zh{促销} \emph{cùxiāo} — promotion des ventes \\
\zh{消费者行为} \emph{xiāofèizhě xíngwéi} — comportement du consommateur \\
\zh{策略} \emph{cèlüè} — stratégie
\end{exemplebox}

\begin{textebox}[Qu'est-ce que le marketing ?]
营销不是单纯的“卖东西”，而是“让产品遇见需要它的人”。\\
Yíngxiāo bùshì dānchún de « mài dōngxi », érshì « ràng chǎnpǐn yùjiàn xūyào tā de rén ».\\
« Le marketing n'est pas simplement "vendre des choses", c'est "faire en sorte que le produit rencontre la personne qui en a besoin". »\\[0.4em]
无论在中国还是在喀麦隆，营销的核心都是：了解需求、传递价值、建立信任。\\
Wúlùn zài Zhōngguó háishì zài Kāmàilóng, yíngxiāo de héxīn dōu shì: liǎojiě xūqiú, chuándí jiàzhí, jiànlì xìnrèn.\\
« Que ce soit en Chine ou au Cameroun, le cœur du marketing reste : comprendre les besoins, transmettre la valeur, bâtir la confiance. »
\end{textebox}

\courssub{Le marketing au Cameroun : faits et pratiques}

\begin{textebox}[喀麦隆的营销现场 — Le marketing sur le terrain camerounais]
在喀麦隆，传统市场非常热闹。比如在雅温得的莫库洛市场或杜阿拉的桑加市场，小贩大声吆喝，顾客讨价还价，人情味很浓。\\
Zài Kāmàilóng, chuántǒng shìchǎng fēicháng rènao. Bǐrú zài Yǎwēndé de Mòkùluò Shìchǎng huò Dūālā de Sāngjiā Shìchǎng, xiǎofàn dàshēng yāohe, gùkè tǎojià huánjià, rénqíngwèi hěn nóng.\\
« Au Cameroun, les marchés traditionnels sont très animés. Par exemple au marché Mokolo à Yaoundé ou au marché Sandaga à Douala, les vendeurs crient pour attirer le client, les clients marchent, l'ambiance est très humaine. »\\[0.4em]
广告主要靠收音机、电视、路边的大幅海报和车身广告。最近几年，Facebook 和 WhatsApp 成了小商贩最重要的“店铺”，大家在群里发图片、发语音卖衣服、卖鞋、卖吃的。\\
Guǎnggào zhǔyào kào shōuyīnjī, diànshì, lùbiān de dàfú hǎibào hé chēshēn guǎnggào. Zuìjìn jǐ nián, Facebook hé WhatsApp chéngle xiǎoshāngfàn zuì zhòngyào de « diànpù », dàjiā zài qún li fā túpiàn, fā yǔyīn mài yīfu, mài xié, mài chī de.\\
« La publicité passe surtout par la radio, la télé, les grandes affiches au bord de la route et la pub sur les véhicules. Ces dernières années, Facebook et WhatsApp sont devenus la "boutique" la plus importante des petits commerçants : on y poste photos et messages vocaux pour vendre vêtements, chaussures, nourriture. »\\[0.4em]
付钱方面，现金还是王道，但 Orange Money 和 MTN MoMo 让手机付钱变得普遍，不用找零，很方便。\\
Fù qián fāngmiàn, xiànjīn hái shì wángdào, dàn Orange Money hé MTN MoMo ràng shǒujī fù qián biànde pǔbiàn, bùyòng zhǎolíng, hěn fāngbiàn.\\
« Côté paiement, le cash reste roi, mais Orange Money et MTN MoMo ont généralisé le paiement mobile, pas de monnaie à rendre, très pratique. »\\[0.4em]
喀麦隆人买东西看重“关系”和“信任”，认识老板、熟人介绍往往比广告更管用。\\
Kāmàilóng rén mǎi dōngxi kànzhòng « guānxì » hé « xìnrèn », rènshi lǎobǎn, shúrén jièshào wǎngwǎng bǐ guǎnggào gèng guǎnyòng.\\
« Les Camerounais achètent en misant sur la "relation" et la "confiance" : connaître le patron, une recommandation d'un proche valent souvent mieux que la pub. »
\end{textebox}

\begin{center}
\begin{tabularx}{\linewidth}{|X|X|X|X|}
\hline
\rowcolor{popblueL} Caractères & Pinyin & Nature & Traduction \\
\hline
热闹 & rènao & adj. & animé, vivant \\
\hline
小贩 & xiǎofàn & nom & petit vendeur ambulant \\
\hline
吆喝 & yāohe & verbe & crier (pour vendre), héler le client \\
\hline
讨价还价 & tǎojià huánjià & verbe & marchander, négocier le prix \\
\hline
人情味 & rénqíngwèi & nom & chaleur humaine, convivialité \\
\hline
收音机 & shōuyīnjī & nom & radio \\
\hline
车身广告 & chēshēn guǎnggào & nom & publicité sur véhicule \\
\hline
小商贩 & xiǎoshāngfàn & nom & petit commerçant \\
\hline
语音 & yǔyīn & nom & message vocal \\
\hline
现金 & xiànjīn & nom & argent liquide \\
\hline
找零 & zhǎolíng & verbe & rendre la monnaie \\
\hline
关系 & guānxì & nom & relations, réseau, connexions \\
\hline
信任 & xìnrèn & nom/verbe & confiance / faire confiance \\
\hline
管用 & guǎnyòng & adj. & efficace, qui marche \\
\hline
\end{tabularx}
\end{center}

\begin{attention}[讨价还价  vs  砍价]
\zh{讨价还价} \emph{tǎojià huánjià} (litt. « demander le prix, rendre le prix ») est le terme standard pour « marchander, négocier ». \zh{砍价} \emph{kǎnjià} (litt. « couper le prix ») est plus familier, très utilisé à l'oral. Les deux sont corrects. Exemple : \zh{在市场买东西要会讨价还价。} \emph{Zài shìchǎng mǎi dōngxi yào huì tǎojià huánjià.} « Au marché, il faut savoir marchander. »
\end{attention}

\courssub{Le marketing en Chine : faits et pratiques}

Dans la section précédente, nous avons observé comment le marketing fonctionne au Cameroun : marchés vivants, crieurs, radio, affiches et réseaux sociaux en pleine croissance. Traversons maintenant le monde et découvrons la Chine, où le marketing a pris un virage spectaculaire : celui du \cle{tout-numérique}. Comme nous allons le voir, les deux modèles sont différents, mais aucun n'est « meilleur » que l'autre : chacun répond à l'histoire et aux habitudes de son pays.

\courssub{Observation : un texte pour découvrir}

Lisons d'abord ce court texte original. Soulignez mentalement les mots qui désignent des moyens de paiement et des plateformes de vente.

\begin{textebox}[在中国买东西 — Acheter en Chine]
在中国，很多人用手机付钱，也在直播里买东西。\\
Zài Zhōngguó, hěn duō rén yòng shǒujī fù qián, yě zài zhíbō lǐ mǎi dōngxi.\\
« En Chine, beaucoup de gens paient avec leur téléphone et achètent aussi pendant des directs vidéo. »\\[0.4em]
中国人很喜欢在网上买东西。淘宝、京东和拼多多是最有名的购物平台。买东西的时候，大家很少用现金：他们打开支付宝或者微信支付，扫一下二维码，钱就付了。非常方便！\\
Zhōngguó rén hěn xǐhuan zài wǎng shàng mǎi dōngxi. Táobǎo, Jīngdōng hé Pīnduōduō shì zuì yǒumíng de gòuwù píngtái. Mǎi dōngxi de shíhou, dàjiā hěn shǎo yòng xiànjīn: tāmen dǎkāi Zhīfùbǎo huòzhě Wēixìn Zhīfù, sǎo yí xià èrwéimǎ, qián jiù fù le. Fēicháng fāngbiàn!\\
« Les Chinois aiment beaucoup acheter sur Internet. Taobao, JD.com et Pinduoduo sont les plateformes d'achat les plus célèbres. Quand on achète quelque chose, on utilise rarement du liquide : on ouvre Alipay ou WeChat Pay, on scanne un code QR, et c'est payé. Très pratique ! »\\[0.4em]
最近几年，直播卖东西非常流行。主播在手机前面介绍产品，几百万观众一边看，一边下单。\\
Zuìjìn jǐ nián, zhíbō mài dōngxi fēicháng liúxíng. Zhǔbō zài shǒujī qiánmiàn jièshào chǎnpǐn, jǐ bǎi wàn guānzhòng yìbiān kàn, yìbiān xiàdān.\\
« Ces dernières années, vendre en direct vidéo est devenu très populaire. Le présentateur présente les produits devant son téléphone, et des millions de spectateurs commandent tout en regardant. »\\[0.4em]
每年十一月十一日是“双十一”，这是中国最大的购物节。那一天，很多东西都打折，人们在网上花很多钱。\\
Měi nián shíyī yuè shíyī rì shì « Shuāng Shíyī », zhè shì Zhōngguó zuì dà de gòuwù jié. Nà yī tiān, hěn duō dōngxi dōu dǎzhé, rénmen zài wǎng shàng huā hěn duō qián.\\
« Chaque année, le 11 novembre, c'est le "Double 11", la plus grande fête du shopping de Chine. Ce jour-là, beaucoup d'articles sont soldés et les gens dépensent beaucoup d'argent en ligne. »\\[0.4em]
当然，传统广告也还存在：电视上有广告，地铁里有海报，大城市的高楼上有很大的电子屏幕。\\
Dāngrán, chuántǒng guǎnggào yě hái cúnzài: diànshì shàng yǒu guǎnggào, dìtiě lǐ yǒu hǎibào, dà chéngshì de gāolóu shàng yǒu hěn dà de diànzǐ píngmù.\\
« Bien sûr, la publicité traditionnelle existe toujours : des pubs à la télévision, des affiches dans le métro, de grands écrans électroniques sur les immeubles des grandes villes. »
\end{textebox}

\textbf{Questions de compréhension}\par
\begin{enumerate}
\item Quelles sont les trois plateformes de vente en ligne citées dans le texte ?
\item Comment les Chinois paient-ils leurs achats, selon le texte ?
\item Que font les spectateurs pendant un \zh{直播} ?
\item Quelle date correspond au « Double 11 » ?
\item La publicité traditionnelle a-t-elle disparu en Chine ? Justifiez par une phrase du texte.
\end{enumerate}

\begin{corrige}[Questions de compréhension]
1) \zh{淘宝} (Taobao), \zh{京东} (JD.com), \zh{拼多多} (Pinduoduo). 2) Ils paient avec leur téléphone via \zh{支付宝} (Alipay) ou \zh{微信支付} (WeChat Pay) en scannant un \zh{二维码} (code QR). 3) Ils regardent et commandent en même temps (\zh{一边看，一边下单}). 4) Le 11 novembre (11/11). 5) Non, elle existe encore : \zh{传统广告也还存在}.
\end{corrige}

\courssub{Le poids du commerce en ligne et du paiement mobile}

En Chine, le commerce en ligne (\zh{网购} \emph{wǎnggòu}) occupe une place considérable dans la vie quotidienne. Trois plateformes (\zh{平台} \emph{píngtái}) dominent le marché :

\begin{itemize}
\item \zh{淘宝} \emph{Táobǎo} : immense place de marché (\zh{集市} \emph{jíshì}) où des millions de petits vendeurs (\zh{小商家} \emph{xiǎo shāngjiā}) proposent une variété infinie de produits ;
\item \zh{京东} \emph{Jīngdōng} (JD.com) : plateforme réputée pour sa logistique propre et sa livraison ultra-rapide (\zh{物流快} \emph{wùliú kuài}), forte en électronique et électroménager ;
\item \zh{拼多多} \emph{Pīnduōduō} : plateforme d'achats groupés (\zh{拼团} \emph{pīntuán}) à prix réduits, très populaire dans les villes de rang inférieur et zones rurales.
\end{itemize}

Mais la vraie révolution, c'est le \cle{paiement mobile} (\zh{移动支付} \emph{yídòng zhīfù}). Deux applications dominent : \zh{支付宝} \emph{Zhīfùbǎo} (Alipay, né d'Alibaba) et \zh{微信支付} \emph{Wēixìn Zhīfù} (WeChat Pay, intégré à WeChat). Dans les villes chinoises, on paie le restaurant, le taxi, le marché et même le petit vendeur de rue (\zh{路边摊} \emph{lùbiāntān}) en scannant un \zh{二维码} \emph{èrwéimǎ} (code QR) avec son téléphone. Le liquide (\zh{现金} \emph{xiànjīn}) devient rare ; beaucoup de jeunes ne portent plus de portefeuille.

\begin{definition}[付钱  — Payer (une transaction)]
\zh{付钱} \emph{fù qián} signifie « payer, régler une somme due ». Structure : \textbf{Sujet + 用 + Moyen + 付钱} ou \textbf{Sujet + 付 + (somme) + 钱}.
\zh{付} \emph{fù} = verser, payer (contexte transactionnel) ; \zh{钱} \emph{qián} = argent.
Exemple : \zh{我用微信付钱。} \emph{Wǒ yòng Wēixìn fù qián.} « Je paie avec WeChat. »
\end{definition}

\begin{exemplebox}[Le paiement mobile en phrases]
\zh{他扫了一下二维码，钱就付了。} \emph{Tā sǎo le yí xià èrwéimǎ, qián jiù fù le.} « Il a scanné le code QR, et c'était payé. »\\
\zh{在中国，用现金的人越来越少。} \emph{Zài Zhōngguó, yòng xiànjīn de rén yuè lái yuè shǎo.} « En Chine, les gens qui utilisent du liquide sont de moins en moins nombreux. »\\
\zh{你用支付宝还是微信支付？} \emph{Nǐ yòng Zhīfùbǎo háishi Wēixìn Zhīfù?} « Tu utilises Alipay ou WeChat Pay ? »\\
\zh{这个小摊也支持扫码付钱。} \emph{Zhège xiǎotān yě zhīchí sǎomǎ fù qián.} « Ce petit stand accepte aussi le paiement par scan. »
\end{exemplebox}

\courssub{Le livestreaming : vendre en direct (\zh{直播带货})}

\begin{definition}[直播  — Diffusion en direct / Livestreaming]
\zh{直播} \emph{zhíbō} (nom/verbe) = \zh{直} \emph{zhí} (direct) + \zh{播} \emph{bō} (diffuser).
\zh{主播} \emph{zhǔbō} = \zh{主} \emph{zhǔ} (principal) + \zh{播} \emph{bō} = présentateur/animateur du direct.
\zh{直播带货} \emph{zhíbō dài huò} = « direct vidéo qui emporte les marchandises » = vente par livestreaming.
\end{definition}

Le \zh{直播带货} est devenu un phénomène majeur : un vendeur (\zh{主播}) présente des produits en direct, répond aux questions en temps réel, offre des coupons limités (\zh{优惠券} \emph{yōuhuìquàn}), et les spectateurs (\zh{观众} \emph{guānzhòng}) commandent (\zh{下单} \emph{xià dān}) d'un clic sans quitter la vidéo. Les têtes d'affiche (\zh{头部主播} \emph{tóubù zhǔbō}) comme Li Jiaqi ou Viya ont réalisé des milliards de yuans de ventes en une soirée.

\begin{popfigure}
\begin{tikzpicture}[node distance=1.8cm, every node/.style={font=\small}]
\node[draw=popblue, fill=popblueL, rounded corners, thick, align=center, minimum width=3.5cm, minimum height=1.2cm] (vendeur) {\zh{主播}\\ Présentateur en direct};
\node[draw=poporange, fill=poporangeL, rounded corners, thick, align=center, minimum width=3.5cm, minimum height=1.2cm, right=of vendeur] (public) {\zh{观众}\\ Spectateurs / Clients};
\node[draw=popgreen, fill=popgreenL, rounded corners, thick, align=center, minimum width=3.5cm, minimum height=1.2cm, right=of public] (commande) {\zh{下单}\\ Passer commande};
\draw[-{Stealth}, very thick, popdark] (vendeur) -- node[above, font=\footnotesize, align=center]{\zh{介绍产品}\\ présente produits} (public);
\draw[-{Stealth}, very thick, popdark] (public) -- node[above, font=\footnotesize, align=center]{\zh{互动提问 / 下单}\\ interagit / commande} (commande);
\node[draw=poppink, fill=poppinkL, rounded corners, thick, align=center, minimum width=3.5cm, minimum height=1cm, below=1.2cm of public] (platform) {\zh{平台}\\ (抖音/快手/淘宝直播)};
\draw[-{Stealth}, thick, popmuted] (vendeur) -- (platform);
\draw[-{Stealth}, thick, popmuted] (public) -- (platform);
\draw[-{Stealth}, thick, popmuted] (commande) -- (platform);
\end{tikzpicture}
\legende{L'écosystème de la vente en direct : le présentateur, le public, la commande, tous sur une plateforme unique.}
\end{popfigure}

\begin{propriete}[Structure : 一边……一边…… — Actions simultanées]
Pour exprimer que deux actions se déroulent en même temps, on utilise : \textbf{Sujet + 一边 + Verbe 1 + 一边 + Verbe 2}.
\zh{观众一边看，一边下单。} \emph{Guānzhòng yìbiān kàn, yìbiān xiàdān.} « Les spectateurs regardent et commandent en même temps. »
\emph{Attention :} Les deux verbes partagent le même sujet. Ne pas confondre avec \zh{又……又……} (adjectifs/états simultanés).
\end{propriete}

\courssub{Les super-applications : WeChat, le couteau suisse numérique}

\zh{微信} \emph{Wēixìn} (WeChat) n'est pas une simple messagerie : c'est une \cle{super-application} (\zh{超级应用} \emph{chāojí yìngyòng}). Avec un seul compte, un utilisateur peut : discuter (\zh{聊天} \emph{liáotiān}), publier sur son fil (\zh{朋友圈} \emph{péngyouquān}), payer (\zh{微信支付}), prendre un taxi (\zh{打车} \emph{dǎchē}), lire des articles (\zh{公众号文章} \emph{gōngzhònghào wénzhāng}), faire des achats (\zh{小程序商店} \emph{xiǎochéngxù shāngdiàn}).

Pour les marques, c'est un outil de marketing intégré : comptes officiels (\zh{公众号} \emph{gōngzhònghào}), mini-programmes boutiques (\zh{小程序} \emph{xiǎochéngxù}), publicité ciblée (\zh{广告投放} \emph{guǎnggào tóufàng}), service client.

\begin{exemplebox}[Autour de WeChat]
\zh{我在微信上看到一条广告。} \emph{Wǒ zài Wēixìn shàng kàndào yì tiáo guǎnggào.} « J'ai vu une pub sur WeChat. »\\
\zh{这个品牌有自己的小程序商城。} \emph{Zhège pǐnpái yǒu zìjǐ de xiǎochéngxù shāngchéng.} « Cette marque a sa propre boutique mini-programme. »\\
\zh{我们在群里讨价还价。} \emph{Wǒmen zài qún li tǎojià huánjià.} « On marchande dans le groupe (WeChat). »
\end{exemplebox}

\courssub{Le Double 11 : la plus grande journée de shopping au monde}

Le \zh{双十一} \emph{Shuāng Shíyī} (« Double Onze »), le 11 novembre, est la plus grande journée de vente en ligne au monde (\zh{全球最大的网购狂欢节} \emph{quánqiú zuì dà de wǎnggòu kuánghuàn jié}).

\begin{itemize}
\item \textbf{Origine :} Années 1990, fête informelle des célibataires (\zh{光棍节} \emph{guānggùn jié}) : les quatre « 1 » de 11/11 symbolisent des personnes seules.
\item \textbf{Commercialisation :} 2009, Alibaba lance la première promo \zh{双十一} sur Taobao.
\item \textbf{Aujourd'hui :} Toutes les plateformes participent. Préventes (\zh{预售} \emph{yùshòu}) dès octobre, galas télévisés (\zh{晚会} \emph{wǎnhuì}) avec stars, compte à rebours jusqu'à minuit. Volume de transactions (\zh{成交额} \emph{chéngjiāo'é}) record : > 1000 milliards de yuans (≈ 130 Mds €) en 2023 sur l'ensemble de la période.
\end{itemize}

\begin{popfigure}
\begin{tikzpicture}[every node/.style={font=\small}]
\draw[-{Stealth}, very thick, popdark] (0,0) -- (12,0);
\foreach \x/\annee in {0.5/2009, 3/2012, 5.5/2016, 8/2020, 10.5/2023, 11.5/今天}
  \draw[thick, popdark] (\x,0.12) -- (\x,-0.12) node[below, align=center]{\annee};
\node[draw=popblue, fill=popblueL, rounded corners, thick, align=center, above=0.4cm] at (0.5,0) {Premier \zh{双十一}\\ (Taobao, 2009)};
\node[draw=poporange, fill=poporangeL, rounded corners, thick, align=center, above=0.4cm] at (5.5,0) {Toutes plateformes\\ \zh{全民购物节}};
\node[draw=popgreen, fill=popgreenL, rounded corners, thick, align=center, above=0.4cm] at (10.5,0) {Record mondial\\ > 1000 Mds ¥};
\node[draw=poppurple, fill=poppurpleL, rounded corners, thick, align=center, above=0.4cm] at (11.5,0) {Écosystème mature\\ Préventes + Direct};
\end{tikzpicture}
\legende{Frise : du premier Double 11 (2009) à l'écosystème mature d'aujourd'hui.}
\end{popfigure}

\begin{savaistu}[Le Double 11, né en 2009]
Le « Double 11 » a été créé en 2009 par le groupe Alibaba (Daniel Zhang). En seulement 24 heures, cette journée génère aujourd'hui plus de ventes que le Black Friday américain. Des émissions télévisées géantes, avec des stars internationales, accompagnent le compte à rebours jusqu'à minuit, heure du début des promotions ! Le concept s'est exporté : on trouve maintenant des « 11.11 sales » en Asie du Sud-Est et en Europe.
\end{savaistu}

\courssub{La publicité traditionnelle existe toujours}

Malgré ce monde numérique, la publicité traditionnelle n'a pas disparu en Chine : spots à la télévision (\zh{电视广告} \emph{diànshì guǎnggào}), affiches dans le métro (\zh{地铁海报} \emph{dìtiě hǎibào}) et écrans géants lumineux sur les immeubles des grandes villes comme Shanghai (\zh{外滩} \emph{Wàitān}), Pékin (\zh{王府井} \emph{Wángfǔjǐng}) ou Shenzhen. Le marketing chinois combine donc l'ancien (\zh{传统} \emph{chuántǒng}) et le nouveau (\zh{新媒体} \emph{xīn méitǐ}).

\begin{center}
\begin{tabularx}{\linewidth}{|X|X|X|X|}
\hline
\rowcolor{popblueL} Caractères & Pinyin & Nature & Traduction \\
\hline
营销 & yíngxiāo & nom/verbe & marketing \\
\hline
消费者 & xiāofèizhě & nom & consommateur \\
\hline
市场 & shìchǎng & nom & marché \\
\hline
品牌 & pǐnpái & nom & marque \\
\hline
促销 & cùxiāo & verbe/nom & promotion / promouvoir \\
\hline
付钱 & fù qián & verbe & payer (transaction) \\
\hline
现金 & xiànjīn & nom & argent liquide, espèces \\
\hline
二维码 & èrwéimǎ & nom & code QR \\
\hline
扫码 & sǎomǎ & verbe & scanner un code \\
\hline
直播 & zhíbō & nom/verbe & diffusion en direct, livestreaming \\
\hline
主播 & zhǔbō & nom & présentateur en direct \\
\hline
观众 & guānzhòng & nom & spectateurs, public \\
\hline
下单 & xià dān & verbe & passer commande \\
\hline
平台 & píngtái & nom & plateforme (en ligne) \\
\hline
打折 & dǎzhé & verbe & faire une réduction, solder \\
\hline
广告 & guǎnggào & nom & publicité \\
\hline
海报 & hǎibào & nom & affiche, poster \\
\hline
购物节 & gòuwù jié & nom & fête du shopping \\
\hline
双十一 & Shuāng Shíyī & nom propre & Double 11 (11 novembre) \\
\hline
预售 & yùshòu & nom/verbe & prévente \\
\hline
成交额 & chéngjiāo'é & nom & volume de transactions (CA) \\
\hline
超级应用 & chāojí yìngyòng & nom & super-application \\
\hline
小程序 & xiǎochéngxù & nom & mini-programme (dans WeChat) \\
\hline
\end{tabularx}
\end{center}

\courssub{Comparaison respectueuse : deux modèles, aucune hiérarchie}

Comment situer les deux pays ? De façon factuelle : la Chine est plus avancée dans le paiement mobile, le commerce en ligne et le livestreaming ; le Cameroun, lui, garde des marchés physiques très vivants, où le contact humain, la négociation (\zh{讨价还价}) et le lien social (\zh{关系}) restent au cœur de l'achat. Le \cle{mobile money} (Orange Money, MTN MoMo) y connaît une croissance rapide, rapprochant les usages du paiement mobile chinois. Ce sont \textbf{deux modèles différents, pas hiérarchisés} : chacun a ses forces adaptées à son contexte.

\begin{center}
\begin{tabularx}{\linewidth}{|X|X|X|}
\hline
\rowcolor{popblueL} Aspect & Chine & Cameroun \\
\hline
Paiement dominant & Mobile (\zh{支付宝}, \zh{微信支付}) — quasi sans numéraire & Liquide (\zh{现金}) dominant, \zh{Mobile Money} en forte croissance \\
\hline
Lieu d'achat principal & Plateformes en ligne (\zh{淘宝}, \zh{京东}, \zh{拼多多}) + \zh{直播} & Marchés physiques (\zh{市场}) très vivants + Réseaux sociaux (FB, WhatsApp) \\
\hline
Nouveauté marketing & \zh{直播带货} (vente en direct), \zh{短视频} marketing & Commerce via WhatsApp/FB Groups, influenceurs locaux naissants \\
\hline
Relation client & Algorithmes, avis notés, service client standardisé & Relation personnelle (\zh{关系}), confiance, recommandation bouche-à-oreille \\
\hline
Point fort & Rapidité, choix immense, traçabilité, promo massives (\zh{双十一}) & Contact humain, négociation, flexibilité, ancrage local \\
\hline
Défi actuel & Saturation, confiance (contrefaçons), protection données & Inclusion numérique, logistique, formalisation \\
\hline
\end{tabularx}
\end{center}

\begin{exoresolu}[Traduire et compléter]
Énoncé : 1) Traduisez en chinois : « En Chine, beaucoup de gens paient avec leur téléphone. » 2) Complétez : \zh{主播在直播里介绍产品，观众一边看，一边\_\_\_\_。} (passer commande)
\tcblower
Solution : 1) \zh{在中国，很多人用手机付钱。} \emph{Zài Zhōngguó, hěn duō rén yòng shǒujī fù qián.} Structure : \zh{用} \emph{yòng} (moyen) + \zh{付钱} \emph{fù qián}. 2) Le mot manquant est \zh{下单} \emph{xià dān} : \zh{主播在直播里介绍产品，观众一边看，一边下单。} \emph{Zhǔbō zài zhíbō lǐ jièshào chǎnpǐn, guānzhòng yìbiān kàn, yìbiān xiàdān.} La structure \zh{一边……一边……} exprime deux actions simultanées du même sujet.
\end{exoresolu}

\begin{attention}[付钱  vs  给钱 — Piège classique francophone]
\zh{付钱} \emph{fù qián} = « payer » dans une transaction commerciale (achat, facture, loyer). Verbe technique : \zh{付} = verser, régler.
\zh{给钱} \emph{gěi qián} = « donner de l'argent » (à quelqu'un). \zh{给} = donner (direction vers un bénéficiaire).
\textbf{Erreur :} \zh{*我给钱买这个。} (Je donne de l'argent pour acheter ça — maladroit).
\textbf{Correct :} \zh{我付钱买这个。} / \zh{我给钱给老板。} (Je donne l'argent au patron — action de donner).
Retenez : \textbf{Transaction = 付钱} ; \textbf{Don / Transfert à personne = 给钱}.
\end{attention}

\begin{methode}[Réussir un mini-exposé comparatif (oral / écrit)]
\begin{enumerate}
\item \textbf{Choisir un angle précis} : ex. « Le paiement », « La place du marché physique », « Le rôle des influenceurs ».
\item \textbf{Structurer en 3 temps} : \zh{首先} (Premièrement) — Présenter le fait chinois ; \zh{其次} (Deuxièmement) — Présenter le fait camerounais ; \zh{最后} (Enfin) — Donner votre avis / synthèse.
\item \textbf{Utiliser les connecteurs logiques} : \zh{但是} (mais), \zh{相比之下} (en comparaison), \zh{因此} (donc), \zh{比如} (par exemple).
\item \textbf{Varier le vocabulaire} : n'utilisez pas seulement \zh{有} / \zh{没有}, mais \zh{普及} (généralisé), \zh{流行} (populaire), \zh{习惯} (habitude), \zh{方便} (pratique).
\item \textbf{S'entraîner à l'oral} : chronométrez-vous (1 min 30 -- 2 min), enregistrez-vous, vérifiez les tons (\zh{zhíbō}, \zh{zhǔbō}, \zh{xiàdān}, \zh{dǎzhé}).
\end{enumerate}
\end{methode}

\begin{experience}[Débat guidé : « Le marché physique va-t-il disparaître ? »]
\textbf{Consigne :} En groupes de 4. Deux élèves défendent l'idée « Le tout-numérique va remplacer les marchés physiques » (argument : efficacité, prix, choix). Deux élèves défendent « Le marché physique a un avenir » (argument : lien social, vérification produit, négociation, emploi).
\textbf{Vocabulaire utile :} \zh{替代} (remplacer), \zh{共存} (coexister), \zh{体验} (expérience), \zh{信任度} (degré de confiance), \zh{就业} (emploi).
\textbf{Production attendue :} 3 arguments par camp en chinois simple, conclusion synthétique.
\end{experience}

\begin{exercice}[À vous — Entraînement complet]
\begin{enumerate}
\item \textbf{Traduction (Version) :} Traduisez en français : \zh{双十一那天，很多东西都打折，人们在网上花很多钱。}
\item \textbf{Traduction (Thème) :} Traduisez en chinois : « Au Cameroun, les marchés sont très animés, on aime marchander. » (Utilisez \zh{热闹} et \zh{讨价还价}).
\item \textbf{Grammaire :} Construisez deux phrases avec \zh{一边……一边……} sur le thème du marketing (ex: regarder une pub / acheter ; faire un direct / vendre).
\item \textbf{Expression écrite :} Rédigez un court paragraphe (5-6 phrases) en chinois pour présenter \zh{双十一} à un ami camerounais qui ne connaît pas. Utilisez : \zh{每年}, \zh{最大}, \zh{打折}, \zh{很方便}, \zh{超级应用}.
\item \textbf{Caractères à écrire (ordre des traits) :} \zh{付}, \zh{直}, \zh{播}, \zh{主}, \zh{折}. Pour chacun : nombre de traits, radical (clé), ordre principal (ex: \zh{付} : 5 traits, radical \zh{亻} (ren), ordre : trait horizontal, trait vertical, puis le côté droit \zh{寸}).
\end{enumerate}
\end{exercice}

\begin{corrige}[Exercice « À vous »]
1) « Le jour du Double 11, beaucoup d'articles sont soldés, les gens dépensent beaucoup d'argent en ligne. »\\
2) \zh{在喀麦隆，市场很热闹，大家喜欢讨价还价。} \emph{Zài Kāmàilóng, shìchǎng hěn rènao, dàjiā xǐhuan tǎojià huánjià.}\\
3) Exemples : \zh{消费者一边看广告，一边下单。} \emph{Xiāofèizhě yìbiān kàn guǎnggào, yìbiān xiàdān.} / \zh{主播一边介绍，一边回答问题。} \emph{Zhǔbō yìbiān jièshào, yìbiān huídá wèntí.}\\
4) Modèle : \zh{双十一是每年十一月十一日，是中国最大的购物节。那天东西都打折，很多人在手机上买东西，非常方便。微信是超级应用，也可以买东西。} \emph{Shuāng Shíyī shì měi nián shíyī yuè shíyī rì, shì Zhōngguó zuì dà de gòuwù jié. Nà tiān dōngxi dōu dǎzhé, hěn duō rén zài shǒujī shàng mǎi dōngxi, fēicháng fāngbiàn. Wēixìn shì chāojí yìngyòng, yě kěyǐ mǎi dōngxi.}\\
5) \zh{付} (5 traits, \zh{亻}) : \zh{亻} → \zh{寸} (horizontal, vertical, horizontal, vertical, hook). \zh{直} (8 traits, \zh{目}) : \zh{十} → \zh{目} (vertical, horizontal×3, vertical, horizontal×2). \zh{播} (12 traits, \zh{扌}) : \zh{扌} → \zh{播} (côté droit : \zh{土}, \zh{田}, \zh{艹}...). \zh{主} (5 traits, \zh{丶}) : point, horizontal, vertical, horizontal, horizontal. \zh{折} (7 traits, \zh{扌}) : \zh{扌} → \zh{斤} (horizontal, vertical, horizontal, vertical, hook, horizontal, horizontal).
\end{corrige}

\coursec{Vocabulaire associé : glossaire et caractères}

Après avoir observé les pratiques du marketing au Cameroun et en Chine — des marchés de Mokolo aux plateformes de vente en ligne chinoises —, il est temps de nous arrêter sur les \cle{mots} qui permettent de parler de tout cela en chinois. Cette section est le cœur lexical de la leçon : nous allons d'abord lire un court texte qui emploie naturellement ce vocabulaire, puis construire le glossaire, explorer les clés (radicaux) qui relient ces caractères entre eux, apprendre à écrire deux caractères importants, et enfin mémoriser les collocations utiles.

\courssub{Texte support : une vendeuse de Douala}

\begin{textebox}[Au marché de Douala]
在杜阿拉的一个大市场里，有一位卖衣服的商人，她叫玛丽。她的产品价格不贵，很便宜，所以每天都有很多顾客。\\\\
Zài Dù'ālā de yí ge dà shìchǎng lǐ, yǒu yí wèi mài yīfu de shāngrén, tā jiào Mǎlì. Tā de chǎnpǐn jiàgé bú guì, hěn piányi, suǒyǐ měi tiān dōu yǒu hěn duō gùkè.\\\\
« Dans un grand marché de Douala, il y a une commerçante qui vend des vêtements ; elle s'appelle Marie. Ses produits ne sont pas chers, ils sont bon marché, c'est pourquoi elle a beaucoup de clients chaque jour. »\\\\
为了吸引更多的顾客，玛丽常常做广告：她在网络上发照片，也在市场门口大声说：“今天打八折！”\\\\
Wèile xīyǐn gèng duō de gùkè, Mǎlì chángcháng zuò guǎnggào : tā zài wǎngluò shang fā zhàopiàn, yě zài shìchǎng ménkǒu dàshēng shuō : “Jīntiān dǎ bā zhé !”\\\\
« Pour attirer encore plus de clients, Marie fait souvent de la publicité : elle publie des photos sur internet, et elle crie aussi à l'entrée du marché : "Aujourd'hui, réduction de 20 \% !" »\\\\
在中国，很多品牌用直播卖东西：一个人在网上介绍产品，顾客看了以后马上付钱。玛丽觉得这种营销方法很有意思，她想明年也试一试。\\\\
Zài Zhōngguó, hěn duō pǐnpái yòng zhíbō mài dōngxi : yí ge rén zài wǎngshang jièshào chǎnpǐn, gùkè kàn le yǐhòu mǎshàng fù qián. Mǎlì juéde zhè zhǒng yíngxiāo fāngfǎ hěn yǒu yìsi, tā xiǎng míngnián yě shì yi shì.\\\\
« En Chine, beaucoup de marques vendent leurs produits en direct vidéo : une personne présente le produit sur internet, et les clients, après avoir regardé, paient aussitôt. Marie trouve cette méthode de marketing très intéressante ; elle veut l'essayer elle aussi l'année prochaine. »
\end{textebox}

\textbf{Questions de compréhension}\par
\begin{enumerate}
\item 玛丽卖什么？(Mǎlì mài shénme ? — Que vend Marie ?)
\item 她怎么吸引顾客？(Tā zěnme xīyǐn gùkè ? — Comment attire-t-elle les clients ?)
\item 中国的品牌用什么方法卖东西？(Zhōngguó de pǐnpái yòng shénme fāngfǎ mài dōngxi ? — Quelle méthode les marques chinoises utilisent-elles pour vendre ?)
\end{enumerate}

\courssub{Glossaire du marketing}

\begin{center}
\begin{tabularx}{\linewidth}{|X|X|X|X|}
\hline
\rowcolor{popblueL} Caractères & Pinyin & Nature & Traduction \\
\hline
营销 & yíngxiāo & n. & marketing \\
\hline
广告 & guǎnggào & n. & publicité \\
\hline
市场 & shìchǎng & n. & marché \\
\hline
顾客 & gùkè & n. & client \\
\hline
产品 & chǎnpǐn & n. & produit \\
\hline
价格 & jiàgé & n. & prix \\
\hline
便宜 & piányi & adj. & bon marché \\
\hline
贵 & guì & adj. & cher \\
\hline
打折 & dǎzhé & v. & faire une réduction \\
\hline
促销 & cùxiāo & v./n. & promotion, promouvoir \\
\hline
品牌 & pǐnpái & n. & marque \\
\hline
网络 & wǎngluò & n. & internet, réseau \\
\hline
直播 & zhíbō & n./v. & direct vidéo, diffuser en direct \\
\hline
付钱 & fù qián & v. & payer \\
\hline
影响 & yǐngxiǎng & v./n. & influencer, influence \\
\hline
\end{tabularx}
\end{center}

\begin{definition}[Deux familles de mots]
Ce glossaire se divise naturellement en deux familles de sens :
\begin{itemize}
\item la famille de l'\cle{argent} et du prix : \zh{价格} (prix), \zh{贵} (cher), \zh{便宜} (bon marché), \zh{付钱} (payer) ;
\item la famille de la \cle{vente} et de la communication : \zh{促销} (promotion), \zh{打折} (faire une réduction), \zh{广告} (publicité), \zh{直播} (direct vidéo).
\end{itemize}
\end{definition}

\begin{exemplebox}[Le vocabulaire en contexte]
\zh{这个产品打折。} \emph{Zhège chǎnpǐn dǎzhé.} « Ce produit est en promotion. »\\
\zh{这个品牌很有名。} \emph{Zhège pǐnpái hěn yǒumíng.} « Cette marque est très connue. »\\
\zh{市场上的东西很便宜。} \emph{Shìchǎng shang de dōngxi hěn piányi.} « Les articles du marché sont bon marché. »\\
\zh{广告影响顾客的选择。} \emph{Guǎnggào yǐngxiǎng gùkè de xuǎnzé.} « La publicité influence le choix des clients. »\\
\zh{他在网络上付钱。} \emph{Tā zài wǎngluò shang fù qián.} « Il paie sur internet. »
\end{exemplebox}

\courssub{Les clés (radicaux) : lire le sens dans le caractère}

En chinois, la plupart des caractères contiennent une \cle{clé} (radical, \zh{部首} bùshǒu) qui donne un indice sur le sens. Le vocabulaire du commerce en est un exemple remarquable.

\begin{propriete}[La clé 贝, le coquillage-monnaie]
La clé \zh{贝} (bèi) représente un \cle{coquillage}. Dans la Chine ancienne, les coquillages servaient de monnaie : c'est pourquoi \zh{贝} apparaît dans de très nombreux caractères liés à l'argent et au commerce :
\begin{itemize}
\item \zh{贵} guì : cher (ce qui coûte beaucoup de « coquillages ») ;
\item \zh{购} gòu : acheter ;
\item \zh{贩} fàn : vendre, revendre (comme dans \zh{小贩} xiǎofàn, petit vendeur) ;
\item \zh{货} huò : marchandise.
\end{itemize}
Quand tu rencontres un caractère inconnu contenant \zh{贝}, tu peux deviner qu'il a un rapport avec l'argent.
\end{propriete}

\begin{popfigure}
\begin{tikzpicture}[grow=down, level distance=1.6cm,
  edge from parent/.style={draw=popblue, thick, -{Stealth}},
  every node/.style={align=center}]
\node[draw=popdark, fill=popgoldL, rounded corners, thick, font=\large, align=center]{贝 bèi\\ coquillage = argent}
  child { node[draw=popblue, fill=popblueL, rounded corners, align=center]{贵 guì\\ cher} }
  child { node[draw=popblue, fill=popblueL, rounded corners, align=center]{购 gòu\\ acheter} }
  child { node[draw=popblue, fill=popblueL, rounded corners, align=center]{贩 fàn\\ vendre} }
  child { node[draw=popblue, fill=popblueL, rounded corners, align=center]{货 huò\\ marchandise} };
\end{tikzpicture}
\legende{La famille du radical 贝 : des caractères liés à l'argent et au commerce.}
\end{popfigure}

Les autres clés utiles de cette leçon :

\begin{center}
\begin{tabularx}{\linewidth}{|X|X|X|}
\hline
\rowcolor{popblueL} Clé & Sens de la clé & Exemples de la leçon \\
\hline
钅 (金) & métal & 销 xiāo (vendre, écouler) \\
\hline
亻 & personne & 促 cù (promouvoir), 付 fù (payer) \\
\hline
囗 & enceinte, cadre & 品 pǐn (produit, article) \\
\hline
扌 & main & 打 dǎ (frapper, faire), 播 bō (diffuser) \\
\hline
\end{tabularx}
\end{center}

\courssub{Écriture des caractères : 贵 et 播}

\begin{methode}[Écrire 贵 (9 traits)]
Le caractère \zh{贵} (guì, cher) s'écrit en 9 traits, de haut en bas :
\begin{enumerate}
\item 丨 : le trait vertical court en haut ;
\item (trait brisé) : le cadre supérieur (horizontal-vertical) ;
\item 一 : l'horizontal qui ferme le haut ;
\item 丨 : le vertical central ;
\item 一 : l'horizontal médian ;
\item 丨 : le vertical du \zh{贝} ;
\item (trait brisé) : le horizontal-vertical du \zh{贝} ;
\item 丿 : le trait tombant à gauche ;
\item 丶 : le point final à droite.
\end{enumerate}
Règle générale respectée : on écrit de \cle{haut en bas}, et l'extérieur avant l'intérieur pour la partie \zh{贝}.
\end{methode}

\begin{methode}[Écrire 播 (15 traits)]
Le caractère \zh{播} (bō, diffuser), que l'on retrouve dans \zh{直播} (direct vidéo), porte la clé de la main \zh{扌} à gauche. On écrit d'abord la clé (3 traits : horizontal, vertical-crochet, trait montant), puis la partie droite de haut en bas, de gauche à droite, en 12 traits. Règle : \cle{gauche avant droite}, \cle{haut avant bas}.
\end{methode}

\begin{attention}[Caractères proches à ne pas confondre]
\begin{itemize}
\item \zh{销} xiāo (clé \zh{钅}, métal) : vendre, écouler — dans \zh{营销} (marketing), \zh{促销} (promotion). Ne pas confondre avec \zh{消} xiāo (clé \zh{氵}, eau) : disparaître — dans \zh{消息} (nouvelles). Même prononciation, sens différents : c'est la clé qui tranche.
\item \zh{场} chǎng (clé \zh{土}, terre) : lieu, terrain — dans \zh{市场} (marché). Ne pas confondre avec \zh{杨} yáng (clé \zh{木}, bois) : peuplier. La partie droite est identique, seule la clé change.
\end{itemize}
\end{attention}

\courssub{Collocations : les mots qui vont ensemble}

En chinois comme en français, certains verbes appellent certains noms. Voici les combinaisons à connaître par cœur :

\begin{center}
\begin{tabularx}{\linewidth}{|X|X|X|}
\hline
\rowcolor{popblueL} Collocation & Pinyin & Traduction \\
\hline
做广告 & zuò guǎnggào & faire de la publicité \\
\hline
打八折 & dǎ bā zhé & réduction de 20 \% \\
\hline
网上购物 & wǎngshàng gòuwù & acheter en ligne \\
\hline
吸引顾客 & xīyǐn gùkè & attirer les clients \\
\hline
促销活动 & cùxiāo huódòng & opération promotionnelle \\
\hline
付现金 & fù xiànjīn & payer en espèces \\
\hline
\end{tabularx}
\end{center}

\begin{attention}[打折 : la logique inverse du français]
\zh{打折} fonctionne à l'envers du français : le chiffre annoncé indique la \cle{part du prix que l'on paie}, pas la réduction. Ainsi \zh{打八折} (dǎ bā zhé) signifie « payer 80 \% du prix », donc une réduction de 20 \%, et non de 80 \% ! De même, \zh{打五折} (dǎ wǔ zhé) = payer 50 \% = moitié prix. Erreur typique du francophone : traduire \zh{打八折} par « réduction de 80 \% ».
\end{attention}

\begin{methode}[Mémoriser par familles de sens]
Pour retenir durablement ce vocabulaire, ne l'apprends pas mot par mot dans le désordre :
\begin{enumerate}
\item Regroupe les mots par famille : argent (\zh{价格}, \zh{贵}, \zh{便宜}, \zh{付钱}) et vente (\zh{促销}, \zh{打折}, \zh{广告}, \zh{直播}).
\item Associe chaque caractère à sa clé (\zh{贝} → argent ; \zh{亻} → action humaine).
\item Construis une phrase personnelle avec chaque mot, par exemple sur ton marché préféré à Yaoundé ou à Douala.
\item Révise les collocations comme des blocs uniques : \zh{做广告}, \zh{网上购物}.
\end{enumerate}
\end{methode}

\begin{exoresolu}[Traduction et vocabulaire]
Énoncé : traduis en chinois les phrases suivantes, en utilisant le glossaire de la leçon. 1) « Ce marché est très grand, les prix ne sont pas chers. » 2) « Cette marque fait de la publicité sur internet pour attirer les clients. »
\tcblower
Solution détaillée :\\
1) « Marché » = \zh{市场} ; « prix » = \zh{价格} ; « cher » = \zh{贵}. On obtient : \zh{这个市场很大，价格不贵。} \emph{Zhège shìchǎng hěn dà, jiàgé bú guì.} Attention : l'adjectif \zh{贵} fonctionne directement comme prédicat, sans verbe « être ».\\
2) « Marque » = \zh{品牌} ; collocation « faire de la publicité » = \zh{做广告} ; « sur internet » = \zh{在网络上} (le complément de lieu se place avant le verbe) ; « pour attirer les clients » = \zh{为了吸引顾客} ou simplement \zh{吸引顾客}. On obtient : \zh{这个品牌在网络上做广告，吸引顾客。} \emph{Zhège pǐnpái zài wǎngluò shang zuò guǎnggào, xīyǐn gùkè.} Piège évité : ne pas placer \zh{在网络上} après le verbe comme en français (« fait de la publicité sur internet ») ; en chinois, le complément de lieu précède le verbe.
\end{exoresolu}

\begin{exercice}[À toi de jouer]
1) Classe les mots suivants dans les deux familles « argent » et « vente » : \zh{促销}, \zh{价格}, \zh{广告}, \zh{便宜}, \zh{付钱}, \zh{打折}.\\
2) Trouve la clé commune de \zh{贵}, \zh{购}, \zh{货} et explique son sens.\\
3) Traduis : « Aujourd'hui, ce produit est à moitié prix » (utilise \zh{打五折}).
\end{exercice}

\begin{corrige}[Exercice « À toi de jouer »]
1) Famille argent : \zh{价格}, \zh{便宜}, \zh{付钱} ; famille vente : \zh{促销}, \zh{广告}, \zh{打折}.\\
2) La clé commune est \zh{贝} (bèi), le coquillage, ancienne monnaie : elle signale un rapport avec l'argent ou le commerce.\\
3) \zh{今天，这个产品打五折。} \emph{Jīntiān, zhège chǎnpǐn dǎ wǔ zhé.} « Aujourd'hui, ce produit est à moitié prix » (on paie 50 \% du prix).
\end{corrige}

\begin{aretenir}[L'essentiel de la section]
Le vocabulaire du marketing s'organise autour de deux familles (argent et vente). La clé \zh{贝} (coquillage-monnaie) relie \zh{贵}, \zh{购}, \zh{贩}, \zh{货}. Les collocations \zh{做广告}, \zh{网上购物}, \zh{吸引顾客} s'apprennent comme des blocs. Enfin, retiens la logique de \zh{打折} : \zh{打八折} = payer 80 \% du prix = réduction de 20 \%. Ce vocabulaire te servira dans les sections suivantes pour comparer les pratiques commerciales du Cameroun et de la Chine et préparer ton mini-exposé.
\end{aretenir}

\coursec{Comparer et donner son avis en chinois}

Dans la section précédente, nous avons constitué un glossaire riche sur le marketing : \zh{广告} (guǎnggào, publicité), \zh{市场} (shìchǎng, marché), \zh{网购} (wǎnggòu, achat en ligne), \zh{顾客} (gùkè, client), \zh{品牌} (pǐnpái, marque), \zh{营销} (yíngxiāo, marketing), \zh{直播带货} (zhíbò dàihuò, vente en direct). Mais un vocabulaire ne sert à rien si l'on ne sait pas l'organiser en phrases. Or, parler du marketing au Cameroun et en Chine, c'est presque toujours \emph{comparer} : dire que les achats en ligne sont plus développés ici, que les marchés sont plus animés là, et \emph{donner son avis}. C'est exactement ce que cette section vous apprend à faire, avec quatre outils grammaticaux simples et puissants : la comparaison avec \zh{比}, la comparaison d'égalité avec \zh{和……一样}, l'expression du point de vue avec \zh{我觉得} et \zh{我认为}, et la nuance avec \zh{虽然……但是……}.

\courssub{Observer d'abord : quatre phrases à deviner}

Lisez ces quatre phrases et essayez de deviner la fonction de chaque mot souligné par le contexte.

\begin{exemplebox}[Phrases d'observation]
\zh{中国的网购比喀麦隆的发达。} \emph{Zhōngguó de wǎnggòu bǐ Kāmàilóng de fādá.} « Le commerce en ligne de la Chine est plus développé que celui du Cameroun. »\\[4pt]
\zh{两个国家的广告一样多吗？} \emph{Liǎng ge guójiā de guǎnggào yíyàng duō ma?} « Les publicités des deux pays sont-elles aussi nombreuses l'une que l'autre ? »\\[4pt]
\zh{我觉得广告很有意思。} \emph{Wǒ juéde guǎnggào hěn yǒu yìsi.} « Je trouve la publicité intéressante. »\\[4pt]
\zh{虽然中国的网购很发达，但是喀麦隆的市场也很热闹。} \emph{Suīrán Zhōngguó de wǎnggòu hěn fādá, dànshì Kāmàilóng de shìchǎng yě hěn rènao.} « Bien que le commerce en ligne chinois soit très développé, les marchés camerounais sont aussi très animés. »
\end{exemplebox}

Vous l'avez remarqué : le chinois ne dit jamais « plus que » avec deux mots séparés comme le français. Un seul mot, \zh{比} (bǐ), placé \emph{entre} les deux termes comparés, fait tout le travail. Observons maintenant chaque structure en détail.

\courssub{La comparaison avec 比 (bǐ)}

\begin{propriete}[Structure de la comparaison : A 比 B + adjectif]
\textbf{Schéma :} A + \cle{比} + B + adjectif\\[4pt]
Cette structure signifie « A est plus [adjectif] que B ». Règles essentielles :
\begin{enumerate}
\item L'adjectif se place \textbf{toujours après B}, jamais avant \zh{比}.
\item L'adjectif n'est \textbf{jamais précédé de} \zh{很} dans cette structure : on dit \zh{网购比市场方便} et non \zh{网购比市场很方便}.
\item Si B est un nom composé (ex: \zh{喀麦隆的网购}), on conserve \zh{的} après B pour marquer l'ellipse : \zh{中国的网购比喀麦隆的发达} (celui du Cameroun).
\item Pour préciser l'écart, on ajoute \textbf{après} l'adjectif : \zh{得多} (de duō, « beaucoup plus ») ou \zh{一点儿} (yìdiǎnr, « un peu plus »).
\end{enumerate}
\end{propriete}

\begin{popfigure}
\begin{tikzpicture}[node distance=0.35cm]
\node[draw=popblue, fill=popblueL, rounded corners, minimum width=2.4cm, minimum height=1cm, align=center] (A) {\textbf{A}\\ 网购\\ \emph{wǎnggòu}};
\node[draw=poporange, fill=poporangeL, rounded corners, minimum width=1.2cm, minimum height=1cm, align=center, right=of A] (bi) {\textbf{比}\\ \emph{bǐ}};
\node[draw=popblue, fill=popblueL, rounded corners, minimum width=2.4cm, minimum height=1cm, align=center, right=of bi] (B) {\textbf{B}\\ 传统市场\\ \emph{chuántǒng shìchǎng}};
\node[draw=poppurple, fill=poppurpleL, rounded corners, minimum width=2.2cm, minimum height=1cm, align=center, right=of B] (adj) {\textbf{adjectif}\\ 方便\\ \emph{fāngbiàn}};
\node[draw=popgreen, fill=popgreenL, rounded corners, minimum width=2.6cm, minimum height=1cm, align=center, right=of adj] (deg) {\textbf{degré (optionnel)}\\ 得多 / 一点儿};
\draw[-{Stealth}, thick, popdark] (A) -- (bi);
\draw[-{Stealth}, thick, popdark] (bi) -- (B);
\draw[-{Stealth}, thick, popdark] (B) -- (adj);
\draw[-{Stealth}, thick, popdark] (adj) -- (deg);
\end{tikzpicture}
\legende{Schéma de la structure comparative : A + 比 + B + adjectif (+ 得多 ou 一点儿).}
\end{popfigure}

\begin{exemplebox}[Exemples traduits]
\zh{网购比传统市场方便。} \emph{Wǎnggòu bǐ chuántǒng shìchǎng fāngbiàn.} « Acheter en ligne est plus pratique que le marché traditionnel. »\\[4pt]
\zh{中国的网购比喀麦隆的多得多。} \emph{Zhōngguó de wǎnggòu bǐ Kāmàilóng de duō de duō.} « Il y a beaucoup plus d'achats en ligne en Chine qu'au Cameroun. »\\[4pt]
\zh{雅温得的市场比杜阿拉的大一点儿。} \emph{Yǎwēndé de shìchǎng bǐ Dù'ālā de dà yìdiǎnr.} « Les marchés de Yaoundé sont un peu plus grands que ceux de Douala. »\\[4pt]
\zh{手机广告比报纸广告贵。} \emph{Shǒujī guǎnggào bǐ bàozhǐ guǎnggào guì.} « La publicité sur téléphone est plus chère que la publicité dans les journaux. »
\end{exemplebox}

\textbf{Comparaison avec le français.} Le français dit « A est \emph{plus} pratique \emph{que} B » : l'adjectif est encadré par deux mots-outils. Le chinois place \zh{比} devant B et l'adjectif tout à la fin : \zh{A 比 B 方便}. L'ordre est donc : A — « que » — B — adjectif. C'est l'inverse du français pour la place de l'adjectif !

\begin{attention}[Erreur typique des francophones]
Ne dites jamais \zh{网购很方便比市场} ni \zh{网购比市场很方便}. Deux fautes guettent le francophone : (1) placer l'adjectif avant \zh{比} par calque du français « plus pratique que » ; (2) garder le \zh{很} automatique devant l'adjectif. Retenez : dans la structure \zh{比}, l'adjectif est \textbf{nu} et \textbf{final} : \zh{网购比市场方便。}
\end{attention}

\courssub{La comparaison d'égalité : A 和 B 一样 + adjectif}

\begin{propriete}[Structure de l'égalité]
\textbf{Schéma :} A + \cle{和} + B + \cle{一样} + adjectif\\[4pt]
Elle signifie « A est aussi [adjectif] que B ». \zh{和} (hé) correspond à « et/avec », \zh{一样} (yíyàng) à « pareil, également ». La négation se fait avec \zh{不一样} : \zh{A 和 B 不一样 + adj.} = « A n'est pas aussi [adj.] que B ». \textbf{Attention :} le sujet A doit être répété ou sous-entendu clairement avant \zh{和} ; on ne dit pas \zh{和市场一样热闹} sans sujet.
\end{propriete}

\begin{exemplebox}[Exemples traduits]
\zh{两个国家的广告一样多吗？} \emph{Liǎng ge guójiā de guǎnggào yíyàng duō ma?} « Y a-t-il autant de publicité dans les deux pays ? »\\[4pt]
\zh{这个品牌和那个品牌一样有名。} \emph{Zhège pǐnpái hé nàge pǐnpái yíyàng yǒumíng.} « Cette marque est aussi connue que celle-là. »\\[4pt]
\zh{喀麦隆的网购和中国的不一样发达。} \emph{Kāmàilóng de wǎnggòu hé Zhōngguó de bù yíyàng fādá.} « Le commerce en ligne camerounais n'est pas aussi développé que celui de la Chine. »
\end{exemplebox}

\begin{savaistu}[Le \zh{不如} (bùrú) : une alternative élégante pour « moins que »]
Au niveau HSK 4, on utilise souvent \zh{不如} (bùrú) pour dire « A n'est pas aussi bon/... que B » (A est inférieur à B). Structure : A + \zh{不如} + B + adjectif. Exemple : \zh{喀麦隆的网购不如中国的发达} (Le e-commerce camerounais est moins développé que le chinois). C'est plus concis que \zh{不一样} pour exprimer l'infériorité.
\end{savaistu}

\courssub{Exprimer son point de vue : 我觉得 et 我认为}

\begin{propriete}[Introduire une opinion]
\cle{我觉得} (wǒ juéde, « je trouve, je pense ») et \cle{我认为} (wǒ rènwéi, « je considère, j'estime ») s'emploient comme « je pense que » en français : ils se placent en tête de phrase, suivis d'une proposition complète. \zh{我觉得} est courant à l'oral ; \zh{我认为} est un peu plus formel, idéal pour un exposé écrit ou oral structuré.
\end{propriete}

\begin{exemplebox}[Exemples traduits]
\zh{我觉得广告很有意思。} \emph{Wǒ juéde guǎnggào hěn yǒu yìsi.} « Je trouve la publicité intéressante. »\\[4pt]
\zh{我认为网购会越来越重要。} \emph{Wǒ rènwéi wǎnggòu huì yuè lái yuè zhòngyào.} « Je pense que les achats en ligne deviendront de plus en plus importants. »\\[4pt]
\zh{我觉得喀麦隆的市场比超市热闹。} \emph{Wǒ juéde Kāmàilóng de shìchǎng bǐ chāoshì rènao.} « Je trouve que les marchés camerounais sont plus animés que les supermarchés. »
\end{exemplebox}

Notez la dernière phrase : on peut \emph{emboîter} une comparaison avec \zh{比} à l'intérieur de \zh{我觉得}. C'est exactement le type de phrase attendu dans votre mini-exposé.

\courssub{Nuancez votre propos : 虽然……但是……}

\begin{propriete}[La concession]
\textbf{Schéma :} \cle{虽然} + phrase 1, \cle{但是} + phrase 2\\[4pt]
« Bien que …, (mais) … ». Contrairement au français, où « bien que » et « mais » ne se combinent pas, le chinois emploie \textbf{les deux mots ensemble} : \zh{虽然} introduit la concession, \zh{但是} (ou \zh{可是}) introduit la réponse. N'en supprimez aucun des deux. La virgule sépare les deux propositions.
\end{propriete}

\begin{exemplebox}[Exemples traduits]
\zh{虽然中国的网购很发达，但是喀麦隆的市场也很热闹。} \emph{Suīrán Zhōngguó de wǎnggòu hěn fādá, dànshì Kāmàilóng de shìchǎng yě hěn rènao.} « Bien que le commerce en ligne chinois soit très développé, les marchés camerounais sont aussi très animés. »\\[4pt]
\zh{虽然广告有时候很贵，但是它很有用。} \emph{Suīrán guǎnggào yǒu shíhou hěn guì, dànshì tā hěn yǒuyòng.} « Bien que la publicité soit parfois chère, elle est très utile. »\\[4pt]
\zh{虽然网购很方便，但是网购和市场不一样热闹。} \emph{Suīrán wǎnggòu hěn fāngbiàn, dànshì wǎnggòu hé shìchǎng bù yíyàng rènao.} « Bien que l'achat en ligne soit pratique, il n'est pas aussi animé que le marché. »
\end{exemplebox}

\courssub{Texte support : regards croisés sur le marketing}

\begin{textebox}[Comparaison marketing Cameroun -- Chine]
\zh{中国和喀麦隆的市场营销不一样。} \emph{Zhōngguó hé Kāmàilóng de shìchǎng yíngxiāo bù yíyàng.}\\[4pt]
\zh{中国的网购非常发达，直播带货很流行。} \emph{Zhōngguó de wǎnggòu fēicháng fādá, zhíbò dàihuò hěn liúxíng.}\\[4pt]
\zh{喀麦隆的传统市场比超市热闹，顾客喜欢讨价还价。} \emph{Kāmàilóng de chuántǒng shìchǎng bǐ chāoshì rènao, gùkè xǐhuan tǎojià huánjià.}\\[4pt]
\zh{虽然中国的广告多，但是喀麦隆的广告也很有创意。} \emph{Suīrán Zhōngguó de guǎnggào duō, dànshì Kāmàilóng de guǎnggào yě hěn yǒu chuàngyì.}\\[4pt]
\zh{我觉得两个国家的营销都很有意思。} \emph{Wǒ juéde liǎng ge guójiā de yíngxiāo dōu hěn yǒu yìsi.}
\end{textebox}

\courssub{Tableau récapitulatif des quatre structures}

\begin{center}
\begin{tabularx}{\linewidth}{|X|X|X|}
\hline
\rowcolor{popblueL} \textbf{Structure} & \textbf{Exemple (caractères + pinyin)} & \textbf{Traduction} \\
\hline
A 比 B + adj. & 网购比市场方便。Wǎnggòu bǐ shìchǎng fāngbiàn. & Acheter en ligne est plus pratique que le marché. \\
\hline
A 比 B + adj. + 得多 / 一点儿 & 中国的网购比喀麦隆的多得多。Zhōngguó de wǎnggòu bǐ Kāmàilóng de duō de duō. & Il y a beaucoup plus d'achats en ligne en Chine qu'au Cameroun. \\
\hline
A 和 B 一样 + adj. & 两个品牌一样有名。Liǎng ge pǐnpái yíyàng yǒumíng. & Les deux marques sont aussi connues l'une que l'autre. \\
\hline
A 和 B 不一样 + adj. & 网购和市场不一样热闹。Wǎnggòu hé shìchǎng bù yíyàng rènao. & L'achat en ligne n'est pas aussi animé que le marché. \\
\hline
我觉得 / 我认为 + phrase & 我觉得广告很有意思。Wǒ juéde guǎnggào hěn yǒu yìsi. & Je trouve la publicité intéressante. \\
\hline
虽然……但是…… & 虽然网购方便，但是市场更热闹。Suīrán wǎnggòu fāngbiàn, dànshì shìchǎng gèng rènao. & Bien que l'achat en ligne soit pratique, le marché est plus animé. \\
\hline
\end{tabularx}
\end{center}

\courssub{Exercice résolu}

\begin{exoresolu}[Traduisez en chinois]
Énoncé : traduisez les phrases suivantes. (1) « La publicité à la télévision est plus chère que la publicité à la radio. » (2) « Je pense que les clients camerounais aiment beaucoup les marchés. » (3) « Bien que les achats en ligne soient pratiques, ils ne sont pas aussi animés que le marché. »
\tcblower
Solution détaillée :
\begin{enumerate}
\item On identifie A = \zh{电视广告}, B = \zh{广播广告}, adjectif = \zh{贵}. L'adjectif va en finale, sans \zh{很} : \zh{电视广告比广播广告贵。} \emph{Diànshì guǎnggào bǐ guǎngbō guǎnggào guì.}
\item On introduit l'opinion avec \zh{我觉得} ou \zh{我认为}, puis une phrase complète : \zh{我觉得喀麦隆的顾客很喜欢市场。} \emph{Wǒ juéde Kāmàilóng de gùkè hěn xǐhuan shìchǎng.}
\item Concession double : \zh{虽然网购很方便，但是网购和市场不一样热闹。} \emph{Suīrán wǎnggòu hěn fāngbiàn, dànshì wǎnggòu hé shìchǎng bù yíyàng rènao.} Notez la répétition du sujet \zh{网购} dans la seconde proposition pour la clarté, et la négation \zh{不一样} pour « pas aussi … que ».
\end{enumerate}
\end{exoresolu}

\courssub{Entraînement oral rapide et préparation du mini-exposé}

\begin{experience}[Comparaisons éclair]
Par groupes de deux, comparez à voix haute, en utilisant obligatoirement \zh{比}, \zh{一样} ou \zh{虽然……但是……} :
\begin{itemize}
\item deux produits : le ndolé et un plat chinois (\zh{Ndolé 比饺子辣。}) ;
\item deux marchés : le marché Mokolo (Yaoundé) et un supermarché de Pékin (\zh{莫科罗市场比北京的超市热闹。}) ;
\item deux publicités : une affiche à Yaoundé et une publicité sur téléphone portable (\zh{手机广告比墙上的广告多。}) .
\end{itemize}
Terminez chaque échange par un avis personnel : \zh{我觉得……}
\end{experience}

\begin{methode}[Construire votre mini-exposé en cinq étapes]
\begin{enumerate}
\item \textbf{Présenter le thème} : \zh{今天我说一说中国和喀麦隆的市场营销。} (Jīntiān wǒ shuō yi shuō Zhōngguó hé Kāmàilóng de shìchǎng yíngxiāo.) « Aujourd'hui, je parle du marketing en Chine et au Cameroun. »
\item \textbf{Décrire le Cameroun} : deux ou trois phrases simples (\zh{喀麦隆的市场很热闹，顾客喜欢讨价还价。}).
\item \textbf{Décrire la Chine} : deux ou trois phrases simples (\zh{中国的网购很发达，直播带货很流行。}).
\item \textbf{Comparer avec} \zh{比} et \zh{一样} : \zh{中国的网购比喀麦隆的多得多，但是市场一样热闹。}
\item \textbf{Donner son avis avec} \zh{我觉得} et conclure avec une nuance : \zh{虽然两个国家不一样，但是我觉得都很有意思。}
\end{enumerate}
\end{methode}

\begin{aretenir}[L'essentiel de la section]
Comparer : A \zh{比} B + adjectif (sans \zh{很}, adjectif en finale ; \zh{得多} / \zh{一点儿} pour le degré ; garder \zh{的} après B si ellipse). Égalité : A \zh{和} B \zh{一样} + adjectif ; infériorité : A \zh{和} B \zh{不一样} + adjectif (ou \zh{不如}). Opinion : \zh{我觉得} / \zh{我认为} + phrase. Nuance : \zh{虽然……但是……}, les deux mots obligatoires. Avec ces quatre outils, vous pouvez comparer le Cameroun et la Chine et défendre votre point de vue en chinois correct.
\end{aretenir}

\begin{exercice}[Entraînement écrit]
\begin{enumerate}
\item Traduisez en chinois : « Les marques chinoises sont plus connues que les marques camerounaises. »
\item Traduisez : « Je trouve que la publicité à la radio n'est pas aussi intéressante que la publicité à la télé. »
\item Complétez : \zh{虽然喀麦隆的网购\_\_\_\_\_\_\_\_，\_\_\_\_\_\_\_\_喀麦隆的市场非常热闹。} (Bien que le e-commerce camerounais \_\_\_\_, \_\_\_\_ les marchés camerounais sont très animés.)
\item Rédigez 5 phrases pour votre mini-exposé en suivant la méthode en 5 étapes.
\end{enumerate}
\end{exercice}

\begin{corrige}[Exercice n]
\begin{enumerate}
\item \zh{中国的品牌比喀麦隆的品牌有名。} (Zhōngguó de pǐnpái bǐ Kāmàilóng de pǐnpái yǒumíng.)
\item \zh{我觉得广播广告和电视广告不一样有意思。} (Wǒ juéde guǎngbō guǎnggào hé diànshì guǎnggào bù yíyàng yǒu yìsi.)
\item \zh{虽然喀麦隆的网购不发达，但是喀麦隆的市场非常热闹。} (Suīrán Kāmàilóng de wǎnggòu bù fādá, dànshì Kāmàilóng de shìchǎng fēicháng rènao.)
\item Production personnelle attendue (exemple) : \zh{今天我说一说中国和喀麦隆的市场营销。喀麦隆的市场很热闹，顾客喜欢讨价还价。中国的网购很发达，直播带货很流行。中国的网购比喀麦隆的多得多，但是市场一样热闹。虽然两个国家不一样，但是我觉得都很有意思。}
\end{enumerate}
\end{corrige}

\coursec{Questions de réflexion et préparation du mini-exposé}

Dans la section précédente, nous avons appris à comparer le Cameroun et la Chine avec la structure \zh{比} et à donner notre avis avec \zh{我觉得}. Nous allons maintenant utiliser ces outils dans une tâche complète et authentique : réfléchir ensemble, débattre avec respect, puis préparer et présenter un \cle{mini-exposé} oral sur le marketing dans les deux pays. C'est exactement le type de production attendue au baccalauréat.

\courssub{Questions de réflexion guidée}

Avant de parler, il faut d'abord penser. Voici deux grandes questions. Pour chacune, observez d'abord les exemples de réponses possibles, puis construisez votre propre réponse.

\textbf{Question 1 : Le marketing influence-t-il nos choix ?}

\zh{营销影响我们的选择吗？} \emph{(Yíngxiāo yǐngxiǎng wǒmen de xuǎnzé ma ?)} « Le marketing influence-t-il nos choix ? »

Exemples de réponses courtes :

\begin{itemize}
\item \zh{我觉得广告影响我们的选择。} \emph{(Wǒ juéde guǎnggào yǐngxiǎng wǒmen de xuǎnzé.)} « Je pense que la publicité influence nos choix. »
\item \zh{有时候我们买不需要的东西。} \emph{(Yǒu shíhou wǒmen mǎi bù xūyào de dōngxi.)} « Parfois nous achetons des choses dont nous n'avons pas besoin. »
\item \zh{但是，聪明的人会比较价格。} \emph{(Dànshì, cōngming de rén huì bǐjiào jiàgé.)} « Mais les gens avisés comparent les prix. »
\end{itemize}

\textbf{Question 2 : Le marché traditionnel va-t-il disparaître ?}

\zh{传统市场会消失吗？} \emph{(Chuántǒng shìchǎng huì xiāoshī ma ?)} « Le marché traditionnel va-t-il disparaître ? »

Exemples de réponses :

\begin{itemize}
\item \zh{我觉得市场不会消失，因为人们喜欢聊天和讲价。} \emph{(Wǒ juéde shìchǎng bú huì xiāoshī, yīnwèi rénmen xǐhuan liáotiān hé jiǎngjià.)} « Je pense que le marché ne disparaîtra pas, car les gens aiment discuter et marchander. »
\item \zh{网购越来越方便，但是市场有人情味。} \emph{(Wǎnggòu yuè lái yuè fāngbiàn, dànshì shìchǎng yǒu rénqíngwèi.)} « Les achats en ligne sont de plus en plus pratiques, mais le marché a une chaleur humaine. »
\end{itemize}

\begin{attention}[Réfléchir avant de répondre]
En chinois comme en français, une bonne réponse suit trois étapes : (1) donner son avis (\zh{我觉得...}), (2) donner une raison (\zh{因为...}), (3) éventuellement nuancer (\zh{但是...}). Évitez les réponses d'un seul mot comme \zh{会} ou \zh{不会} : au baccalauréat, elles ne suffisent pas.
\end{attention}

\courssub{Débat respectueux : acheter au marché ou en ligne ?}

La classe se divise en deux groupes. Le groupe A défend le marché traditionnel, le groupe B défend les achats en ligne. Chaque élève doit présenter au moins un argument en chinois. Voici une banque d'arguments pour vous aider.

\begin{center}
\begin{tabularx}{\linewidth}{|X|X|X|}\hline
\rowcolor{popblueL} Position & Argument en chinois + pinyin & Traduction \\ \hline
Pour le marché & \zh{在市场可以讲价。} \emph{Zài shìchǎng kěyǐ jiǎngjià.} & Au marché, on peut marchander. \\ \hline
Pour le marché & \zh{可以看到和摸到商品。} \emph{Kěyǐ kàndào hé mōdào shāngpǐn.} & On peut voir et toucher les produits. \\ \hline
Pour le marché & \zh{市场帮助小商人。} \emph{Shìchǎng bāngzhù xiǎo shāngrén.} & Le marché aide les petits commerçants. \\ \hline
Pour le marché & \zh{买东西的时候可以和朋友聊天。} \emph{Mǎi dōngxi de shíhou kěyǐ hé péngyou liáotiān.} & On peut discuter avec des amis en faisant ses achats. \\ \hline
En ligne & \zh{网购比去市场方便。} \emph{Wǎnggòu bǐ qù shìchǎng fāngbiàn.} & Acheter en ligne est plus pratique qu'aller au marché. \\ \hline
En ligne & \zh{网上的选择比市场多。} \emph{Wǎngshàng de xuǎnzé bǐ shìchǎng duō.} & Le choix en ligne est plus grand qu'au marché. \\ \hline
En ligne & \zh{可以比较很多价格。} \emph{Kěyǐ bǐjiào hěn duō jiàgé.} & On peut comparer beaucoup de prix. \\ \hline
En ligne & \zh{直播让我们看到商品怎么用。} \emph{Zhíbō ràng wǒmen kàndào shāngpǐn zěnme yòng.} & Les directs nous montrent comment utiliser les produits. \\ \hline
\end{tabularx}
\end{center}

\begin{attention}[Neutralité et respect dans la comparaison culturelle]
Comparer ne signifie pas juger. Il est interdit de se moquer des pratiques d'un pays ou de les dévaloriser : ne dites jamais que le marché camerounais est « arriéré » ni que le commerce en ligne chinois est « inhumain ». Utilisez des formulations factuelles et neutres : \zh{中国的网购比喀麦隆发达} (le commerce en ligne chinois est plus développé que celui du Cameroun) est un fait ; ajouter un jugement de valeur est une faute de méthode et de respect. En chinois, préférez \zh{不一样} (différent) à tout adjectif péjoratif.
\end{attention}

\begin{experience}[Débat en classe]
Déroulement : (1) cinq minutes de préparation en groupe, chaque élève écrit deux arguments avec le tableau ci-dessus ; (2) chaque groupe présente ses arguments à tour de rôle, une phrase par élève ; (3) le professeur note les arguments au tableau ; (4) vote final : chaque élève dit \zh{我觉得...比较好，因为...} (je pense que... est mieux, parce que...). Règle d'or : on écoute sans interrompre, on répond avec \zh{我同意} (je suis d'accord) ou \zh{我不同意，因为...} (je ne suis pas d'accord, parce que...).
\end{experience}

\courssub{Glossaire du thème : le marketing}

\begin{aretenir}[Mots-clés obligatoires pour l'exposé]
Voici le vocabulaire de base de la leçon. Vous devez en utiliser au moins 5 dans votre exposé. Maîtrisez l'écriture, le pinyin (tons compris) et le sens.

\begin{center}
\begin{tabularx}{\linewidth}{|X|X|X|X|}\hline
\rowcolor{popblueL} Caractères & Pinyin & Nature & Traduction \\ \hline
营销 & yíngxiāo & nom/verbe & marketing, commercialiser \\ \hline
广告 & guǎnggào & nom & publicité, annonce \\ \hline
市场 & shìchǎng & nom & marché (lieu ou concept) \\ \hline
网购 & wǎnggòu & verbe/nom & achat en ligne, acheter en ligne \\ \hline
直播 & zhíbō & nom/verbe & direct (live), diffusion en direct \\ \hline
商品 & shāngpǐn & nom & produit, marchandise \\ \hline
价格 & jiàgé & nom & prix \\ \hline
讲价 & jiǎngjià & verbe & marchander, négocier le prix \\ \hline
发达 & fādá & adjectif & développé (économie, technologie) \\ \hline
影响 & yǐngxiǎng & verbe/nom & influencer, influence \\ \hline
\end{tabularx}
\end{center}
\end{aretenir}

\courssub{Consignes du mini-exposé}

\begin{definition}[Consignes officielles du mini-exposé]
Vous allez préparer un exposé oral intitulé \zh{中国和喀麦隆的营销} (le marketing en Chine et au Cameroun). Contraintes obligatoires :
\begin{itemize}
\item \cle{8 à 10 phrases} en chinois, rédigées par écrit puis présentées à l'oral ;
\item au moins \cle{5 mots du glossaire} ci-dessus (营销, 广告, 市场, 网购, 直播, 商品, 价格, 讲价, 发达, 影响...) ;
\item au moins \cle{1 comparaison avec 比} : A + 比 + B + adjectif ;
\item au moins \cle{1 opinion avec 我觉得} suivi d'une phrase complète.
\end{itemize}
\end{definition}

\courssub{Grille de préparation : plan, vocabulaire, structures, prononciation}

Un bon exposé se prépare en quatre étapes. Complétez cette grille avant de rédiger.

\begin{center}
\begin{tabularx}{\linewidth}{|X|X|X|}\hline
\rowcolor{popblueL} Étape & Question à se poser & Outil de la leçon \\ \hline
1. Plan & Quelles parties, dans quel ordre ? & Introduction → Cameroun → Chine → Comparaison → Avis → Conclusion \\ \hline
2. Vocabulaire & Quels 5 mots du glossaire vais-je employer ? & 营销, 广告, 网购, 直播, 市场, 讲价... \\ \hline
3. Structures & Quelles structures grammaticales vais-je montrer ? & A + 比 + B + adj. ; 我觉得 + phrase ; 因为...所以... \\ \hline
4. Prononciation & Quels mots ont des tons difficiles ? & 营销 (yíngxiāo, 2-1), 广告 (guǎnggào, 3-4), 直播 (zhíbō, 2-1) \\ \hline
\end{tabularx}
\end{center}

\begin{popfigure}
\begin{center}
\begin{tikzpicture}[
node distance=0.55cm,
etape/.style={rectangle, rounded corners, draw=popblue, fill=popblueL, thick, align=center, minimum width=4.6cm, minimum height=0.9cm, font=\small},
etapeF/.style={rectangle, rounded corners, draw=popgreen, fill=popgreenL, thick, align=center, minimum width=4.6cm, minimum height=0.9cm, font=\small},
fleche/.style={-{Stealth[length=3mm]}, thick, popdark}
]
\node[etape, align=center] (n1){\textbf{Introduction}\\ 今天我说一说...};
\node[etape, below=of n1, align=center] (n2){\textbf{Le Cameroun}\\ 在喀麦隆，很多人喜欢在市场...};
\node[etape, below=of n2, align=center] (n3){\textbf{La Chine}\\ 在中国，很多人在网上...};
\node[etape, below=of n3, align=center] (n4){\textbf{Comparaison}\\ 中国的网购比喀麦隆发达。};
\node[etape, below=of n4, align=center] (n5){\textbf{Mon avis}\\ 我觉得...};
\node[etapeF, below=of n5, align=center] (n6){\textbf{Conclusion}\\ 谢谢大家！};
\draw[fleche] (n1) -- (n2);
\draw[fleche] (n2) -- (n3);
\draw[fleche] (n3) -- (n4);
\draw[fleche] (n4) -- (n5);
\draw[fleche] (n5) -- (n6);
\end{tikzpicture}
\end{center}
\legende{Organigramme du plan d'exposé : six étapes, de l'introduction à la conclusion.}
\end{popfigure}

\begin{methode}[Réussir son exposé oral en cinq conseils]
\begin{enumerate}
\item \cle{Parlez lentement.} En chinois, la vitesse n'est pas une qualité : la clarté des tons est la qualité. Une phrase lente et correcte vaut mieux que trois phrases rapides et fausses.
\item \cle{Soignez les tons des mots-clés.} Entraînez-vous à voix haute sur : \zh{营销} \emph{yíngxiāo} (tons 2-1 : la voix monte puis reste haute), \zh{广告} \emph{guǎnggào} (tons 3-4 : la voix descend-remonte puis tombe franchement), \zh{直播} \emph{zhíbō} (tons 2-1 : la voix monte puis reste haute).
\item \cle{Regardez votre auditoire} entre deux phrases ; ne lisez pas votre feuille en continu.
\item \cle{Annoncez votre plan} dès la première phrase avec \zh{今天我说一说...} (aujourd'hui, je parle de...).
\item \cle{Terminez proprement} avec \zh{谢谢大家！} \emph{(Xièxie dàjiā !)} « Merci à tous ! », la formule de politesse standard à la fin d'un exposé en Chine.
\end{enumerate}
\end{methode}

\courssub{Relecture collective d'un exposé modèle}

Lisons ensemble ce modèle d'exposé. Il contient exactement 8 phrases, 6 mots du glossaire, une comparaison avec \zh{比} et une opinion avec \zh{我觉得}.

\begin{textebox}[Exposé modèle : 中国和喀麦隆的营销]
今天我说一说中国和喀麦隆的营销。\\
\emph{Jīntiān wǒ shuō yi shuō Zhōngguó hé Kāmàilóng de yíngxiāo.}\\
« Aujourd'hui, je parle du marketing en Chine et au Cameroun. »\\[0.4em]
在喀麦隆，很多人喜欢在市场买东西。\\
\emph{Zài Kāmàilóng, hěn duō rén xǐhuan zài shìchǎng mǎi dōngxi.}\\
« Au Cameroun, beaucoup de gens aiment acheter au marché. »\\[0.4em]
在市场上，人们可以讲价，也可以看到商品。\\
\emph{Zài shìchǎng shàng, rénmen kěyǐ jiǎngjià, yě kěyǐ kàndào shāngpǐn.}\\
« Au marché, on peut marchander et aussi voir les produits. »\\[0.4em]
在中国，很多人在网上买东西，也看直播。\\
\emph{Zài Zhōngguó, hěn duō rén zài wǎngshàng mǎi dōngxi, yě kàn zhíbō.}\\
« En Chine, beaucoup de gens achètent en ligne et regardent aussi des directs. »\\[0.4em]
中国的网购比喀麦隆发达。\\
\emph{Zhōngguó de wǎnggòu bǐ Kāmàilóng fādá.}\\
« Le commerce en ligne chinois est plus développé que celui du Cameroun. »\\[0.4em]
两个国家的广告都很多。\\
\emph{Liǎng gè guójiā de guǎnggào dōu hěn duō.}\\
« Dans les deux pays, il y a beaucoup de publicité. »\\[0.4em]
我觉得两个国家的方法都很有意思。\\
\emph{Wǒ juéde liǎng gè guójiā de fāngfǎ dōu hěn yǒu yìsi.}\\
« Je trouve que les méthodes des deux pays sont très intéressantes. »\\[0.4em]
谢谢大家！\\
\emph{Xièxie dàjiā !}\\
« Merci à tous ! »
\end{textebox}

Vérification collective des contraintes :

\begin{center}
\begin{tabularx}{\linewidth}{|X|X|X|}\hline
\rowcolor{popblueL} Contrainte & Où dans le modèle ? & Validé ? \\ \hline
8 à 10 phrases & 8 phrases & Oui \\ \hline
5 mots du glossaire & 营销, 市场, 讲价, 商品, 网购, 直播, 广告 (7 mots) & Oui \\ \hline
1 comparaison avec 比 & 中国的网购比喀麦隆发达。 & Oui \\ \hline
1 opinion avec 我觉得 & 我觉得两个国家的方法都很有意思。 & Oui \\ \hline
Conclusion polie & 谢谢大家！ & Oui \\ \hline
\end{tabularx}
\end{center}

\begin{exoresolu}[Repérer et corriger les erreurs dans un exposé]
Énoncé : un élève a écrit les trois phrases suivantes pour son exposé. Chacune contient une erreur. Trouvez-la et corrigez-la.\\
(1) \zh{喀麦隆的市场比中国发达。}\\
(2) \zh{我觉得网购很方便，因为不用去市场，所以我觉得网购很方便。}\\
(3) \zh{在中国，很多人买网上东西。}
\tcblower
Solution détaillée :\\
(1) Erreur de fait et de logique : c'est le commerce en ligne chinois qui est plus développé, pas le marché camerounais par rapport à la Chine. Correction factuelle et neutre : \zh{中国的网购比喀麦隆发达。} \emph{(Zhōngguó de wǎnggòu bǐ Kāmàilóng fādá.)} Attention : une comparaison avec \zh{比} doit être vraie et vérifiable.\\
(2) Erreur de répétition : la phrase tourne en rond. En chinois comme en français, on évite de répéter la même idée. Correction : \zh{我觉得网购很方便，因为不用去市场。} \emph{(Wǒ juéde wǎnggòu hěn fāngbiàn, yīnwèi bú yòng qù shìchǎng.)} Une opinion + une raison suffisent.\\
(3) Erreur d'ordre des mots : le complément de lieu \zh{在网上} se place avant le verbe, pas entre le verbe et l'objet. Correction : \zh{在中国，很多人在网上买东西。} \emph{(Zài Zhōngguó, hěn duō rén zài wǎngshàng mǎi dōngxi.)} Rappel : Sujet + 在 + lieu + verbe + objet.
\end{exoresolu}

\begin{attention}[Les trois pièges de l'exposé des francophones]
\begin{itemize}
\item \cle{Ordre des mots} : le lieu se dit avant le verbe (\zh{在网上买}, jamais \zh{买在网上}).
\item \cle{La comparaison} : avec \zh{比}, ne rajoutez pas \zh{很} après : on dit \zh{网购比市场方便}, pas \zh{网购比市场很方便}.
\item \cle{Les tons sous le stress} : à l'oral, les élèves nerveux « aplatissent » les tons. Répétez dix fois \zh{营销} (2-1), \zh{广告} (3-4) et \zh{直播} (2-1) avant de passer.
\end{itemize}
\end{attention}

\begin{savaistu}[Le baccalauréat aime les comparaisons]
Au baccalauréat camerounais, les questions socioculturelles demandent souvent de comparer un fait chinois et un fait camerounais en 3 à 5 phrases : par exemple « Comparez les achats au marché au Cameroun et les achats en ligne en Chine ». La méthode gagnante est toujours la même : une phrase sur le Cameroun, une phrase sur la Chine, une comparaison avec \zh{比}, une opinion avec \zh{我觉得}. Si vous maîtrisez ce schéma en quatre phrases, vous maîtrisez toutes les questions socioculturelles de l'examen.
\end{savaistu}

\begin{exercice}[À vous de préparer]
Préparez votre mini-exposé en suivant la grille de la section 6.4 : rédigez 8 à 10 phrases sur une feuille, soulignez vos 5 mots du glossaire, encadrez votre comparaison avec \zh{比} et votre opinion avec \zh{我觉得}, puis répétez à voix haute trois fois en veillant aux tons de \zh{营销}, \zh{广告} et \zh{直播}. Vous présenterez votre exposé à la séance suivante.
\end{exercice}

\newpage

\coursec{Activité d'intégration}

\courssub{Objectifs de l'activité}

À l'issue de cette activité d'intégration, tu seras capable de :
\begin{itemize}
\item comparer, de façon factuelle et respectueuse, une pratique de marketing camerounaise et une pratique chinoise ;
\item employer correctement au moins cinq mots du glossaire de la leçon (市场， 广告， 顾客， 直播， 打折， 手机支付， etc.) ;
\item construire une comparaison avec la structure \zh{比} et exprimer un point de vue avec \zh{我觉得} ;
\item présenter un mini-exposé oral de 8 à 10 phrases en chinois, clair et structuré ;
\item écouter activement l'exposé des autres groupes et remplir une fiche d'écoute.
\end{itemize}

\courssub{Situation}

Dans le cadre de la semaine « Cameroun--Chine » organisée par ton lycée à Yaoundé, le club de chinois prépare une exposition intitulée \zh{两个市场，一个世界} (\emph{Liǎng ge shìchǎng, yí ge shìjiè} — « Deux marchés, un monde »). Des élèves d'une école partenaire de Pékin participeront en visioconférence. Ton professeur de chinois te propose, avec ton groupe, de présenter un mini-exposé comparatif : « Le marketing chez moi et en Chine ».

Pour préparer l'exposé, le professeur distribue deux documents d'appui. Le premier est un extrait d'article sur le marketing au Cameroun ; le second décrit les pratiques de vente en ligne en Chine.

\begin{textebox}[Document 1 : le marketing au Cameroun]
在喀麦隆，很多公司在广播和电视上做广告。\\
\emph{Zài Kāmàilóng, hěn duō gōngsī zài guǎngbō hé diànshì shang zuò guǎnggào.}\\
« Au Cameroun, beaucoup d'entreprises font de la publicité à la radio et à la télévision. »\\[4pt]
在市场里，商人用大声的叫卖来吸引顾客。\\
\emph{Zài shìchǎng li, shāngrén yòng dà shēng de jiàomài lái xīyǐn gùkè.}\\
« Au marché, les commerçants attirent les clients en criant leurs marchandises. »\\[4pt]
现在，很多年轻人也用 WhatsApp 来卖东西。\\
\emph{Xiànzài, hěn duō niánqīngrén yě yòng WhatsApp lái mài dōngxi.}\\
« Aujourd'hui, beaucoup de jeunes utilisent aussi WhatsApp pour vendre. »
\end{textebox}

\begin{textebox}[Document 2 : le marketing en Chine]
在中国，网上购物非常流行。\\
\emph{Zài Zhōngguó, wǎngshang gòuwù fēicháng liúxíng.}\\
« En Chine, les achats en ligne sont très populaires. »\\[4pt]
十一月十一日是“双十一”，那一天很多商品打折。\\
\emph{Shíyī yuè shíyī rì shì ``Shuāng Shíyī'', nà yì tiān hěn duō shāngpǐn dǎzhé.}\\
« Le 11 novembre, c'est le ``Double 11'' : ce jour-là, beaucoup de produits sont soldés. »\\[4pt]
很多人看直播买东西，然后用手机支付。\\
\emph{Hěn duō rén kàn zhíbō mǎi dōngxi, ránhòu yòng shǒujī zhīfù.}\\
« Beaucoup de gens regardent du livestreaming pour acheter, puis paient par téléphone mobile. »
\end{textebox}

\courssub{Consignes}

Travail en groupes de 3 élèves. Durée totale : 2 heures (préparation + présentations).

\begin{enumerate}
\item \textbf{Lecture et compréhension (15 min).} Lis les deux documents à voix haute dans ton groupe. Souligne dans chaque document : un lieu, un moyen de communication et une action de vente. Vérifie le sens des mots inconnus dans le glossaire de la leçon.

\item \textbf{Choix des pratiques (10 min).} Choisissez ensemble : une pratique camerounaise (le marché, la radio, le parrainage d'une équipe de football, WhatsApp Business) et une pratique chinoise (le livestreaming, le Double 11, le paiement mobile). Chaque élève du groupe se charge d'une partie : présentation du Cameroun, présentation de la Chine, comparaison et conclusion.

\item \textbf{Rédaction (30 min).} Rédigez un exposé de 8 à 10 phrases en chinois simplifié. Contraintes obligatoires :
\begin{itemize}
\item au moins \cle{5 mots du glossaire} de la leçon (surlignez-les dans votre texte) ;
\item au moins \cle{une comparaison} avec la structure \zh{A 比 B + adjectif} ;
\item au moins \cle{une opinion} introduite par \zh{我觉得} ;
\item une phrase d'introduction et une phrase de conclusion polie (\zh{谢谢大家！}).
\end{itemize}

\item \textbf{Répétition (15 min).} Répétez la prononciation : attention aux tons des mots du glossaire (\zh{市场} shìchǎng, \zh{广告} guǎnggào, \zh{打折} dǎzhé). Répartissez la parole équitablement entre les trois membres.

\item \textbf{Présentation (30 min).} Chaque groupe présente devant la classe (et la webcam pour les partenaires de Pékin). Pendant les exposés des autres groupes, remplis ta \cle{fiche d'écoute} : note 2 mots du glossaire entendus et 1 comparaison avec \zh{比} relevée.

\item \textbf{Vote et débat (10 min).} La classe vote pour l'exposé le plus clair et le plus respectueux des deux cultures. Justifie ton vote en une phrase en chinois : \zh{我觉得第……组最好，因为……} (\emph{Wǒ juéde dì ... zǔ zuì hǎo, yīnwèi ...}).
\end{enumerate}

\courssub{Ressources à mobiliser}

\begin{aretenir}[Structures utiles pour l'exposé]
\begin{itemize}
\item \textbf{Comparaison} : A \zh{比} B + adjectif. Ex. : \zh{中国的网上购物比喀麦隆的多。} (\emph{Zhōngguó de wǎngshang gòuwù bǐ Kāmàilóng de duō.}) « Les achats en ligne en Chine sont plus nombreux qu'au Cameroun. »
\item \textbf{Opinion} : \zh{我觉得} + phrase. Ex. : \zh{我觉得直播很有意思。} (\emph{Wǒ juéde zhíbō hěn yǒuyìsi.}) « Je trouve le livestreaming très intéressant. »
\item \textbf{Moyen} : \zh{用} + outil + \zh{来} + action. Ex. : \zh{商人用广播来做广告。} « Les commerçants utilisent la radio pour faire de la publicité. »
\item \textbf{Concession} : \zh{虽然……但是……} « bien que... mais... ».
\end{itemize}
\end{aretenir}

\begin{center}
\begin{tabularx}{\linewidth}{|X|X|X|X|}
\hline
\rowcolor{popblueL} Caractères & Pinyin & Nature & Traduction \\
\hline
市场 & shìchǎng & nom & marché \\
\hline
广告 & guǎnggào & nom & publicité \\
\hline
顾客 & gùkè & nom & client \\
\hline
吸引 & xīyǐn & verbe & attirer \\
\hline
直播 & zhíbō & nom/verbe & livestreaming ; diffuser en direct \\
\hline
打折 & dǎzhé & verbe & faire une réduction, solder \\
\hline
手机支付 & shǒujī zhīfù & nom & paiement mobile \\
\hline
流行 & liúxíng & adjectif & populaire, à la mode \\
\hline
商品 & shāngpǐn & nom & produit, marchandise \\
\hline
越来越 & yuè lái yuè & adverbe & de plus en plus \\
\hline
\end{tabularx}
\end{center}

\begin{attention}[Pièges fréquents]
Ne dis pas \zh{比……很} : après \zh{比}, l'adjectif seul suffit (\zh{中国的人口比喀麦隆多}, sans \zh{很}). Attention aussi à l'ordre : c'est toujours A \zh{比} B + adjectif, jamais l'inverse. Enfin, \zh{觉得} exprime une opinion personnelle, pas un fait : réserve-le à ton jugement.
\end{attention}

\courssub{Proposition de production et corrigé commenté}

\begin{exoresolu}[Modèle de mini-exposé : « Le marketing chez moi et en Chine »]
Rédige un exposé de 8 à 10 phrases comparant une pratique camerounaise et une pratique chinoise, avec 5 mots du glossaire, une comparaison en \zh{比} et une opinion en \zh{我觉得}.
\tcblower
\textbf{Production modèle :}\\[4pt]
\zh{大家好！今天我们谈谈喀麦隆和中国的市场营销。}\\
\emph{Dàjiā hǎo! Jīntiān wǒmen tántan Kāmàilóng hé Zhōngguó de shìchǎng yíngxiāo.}\\
« Bonjour à tous ! Aujourd'hui, nous parlons du marketing au Cameroun et en Chine. »\\[3pt]
\zh{在喀麦隆，很多公司在广播上做广告，市场里的商人用叫卖来吸引顾客。}\\
\emph{Zài Kāmàilóng, hěn duō gōngsī zài guǎngbō shang zuò guǎnggào, shìchǎng li de shāngrén yòng jiàomài lái xīyǐn gùkè.}\\
« Au Cameroun, beaucoup d'entreprises font de la publicité à la radio, et au marché les commerçants attirent les clients en criant leurs marchandises. »\\[3pt]
\zh{现在，年轻人也用 WhatsApp 卖东西，比如卖衣服和食品。}\\
\emph{Xiànzài, niánqīngrén yě yòng WhatsApp mài dōngxi, bǐrú mài yīfu hé shípǐn.}\\
« Aujourd'hui, les jeunes vendent aussi sur WhatsApp, par exemple des vêtements et de la nourriture. »\\[3pt]
\zh{在中国，网上购物越来越流行，很多人看直播买商品。}\\
\emph{Zài Zhōngguó, wǎngshang gòuwù yuè lái yuè liúxíng, hěn duō rén kàn zhíbō mǎi shāngpǐn.}\\
« En Chine, les achats en ligne sont de plus en plus populaires, et beaucoup de gens achètent des produits en regardant du livestreaming. »\\[3pt]
\zh{双十一的时候，很多商品打折，人们用手机支付。}\\
\emph{Shuāng Shíyī de shíhou, hěn duō shāngpǐn dǎzhé, rénmen yòng shǒujī zhīfù.}\\
« Pendant le Double 11, beaucoup de produits sont soldés et les gens paient par téléphone mobile. »\\[3pt]
\zh{中国的网上购物比喀麦隆的多，但是喀麦隆的市场比中国的热闹。}\\
\emph{Zhōngguó de wǎngshang gòuwù bǐ Kāmàilóng de duō, dànshì Kāmàilóng de shìchǎng bǐ Zhōngguó de rènao.}\\
« Les achats en ligne sont plus nombreux en Chine qu'au Cameroun, mais les marchés camerounais sont plus animés que ceux de Chine. »\\[3pt]
\zh{我觉得两种方法都很有意思：直播很方便，叫卖很热情。}\\
\emph{Wǒ juéde liǎng zhǒng fāngfǎ dōu hěn yǒuyìsi: zhíbō hěn fāngbiàn, jiàomài hěn rèqíng.}\\
« Je trouve les deux méthodes intéressantes : le livestreaming est pratique, le cri des commerçants est chaleureux. »\\[3pt]
\zh{谢谢大家！}\\
\emph{Xièxie dàjiā!} « Merci à tous ! »\\[6pt]
\textbf{Commentaire des choix de langue :}
\begin{itemize}
\item \textbf{Glossaire} : 7 mots de la leçon sont employés (市场， 广告， 顾客， 吸引， 直播， 打折， 手机支付， 流行）, soit plus que le minimum exigé.
\item \textbf{Comparaisons} : deux structures \zh{A 比 B + adjectif} correctes, sans \zh{很} après \zh{比} ; la seconde est équilibrée et valorise le Cameroun, ce qui respecte les deux cultures.
\item \textbf{Opinion} : \zh{我觉得} introduit un jugement personnel nuancé (\zh{两种方法都……}), avec la structure parallèle \zh{很……，很……}.
\item \textbf{Cohésion} : connecteurs temporels (\zh{现在}, \zh{……的时候}) et de concession (\zh{但是}) ; ouverture et clôture polies adaptées à un exposé oral.
\end{itemize}
\end{exoresolu}

\courssub{Bilan de l'activité}

Cette activité t'a permis de mobiliser en situation réelle les trois savoir-faire de la leçon : comparer des pratiques camerounaises et chinoises, employer le vocabulaire socioculturel du marketing et exprimer un point de vue simple en chinois. La fiche d'écoute a développé ta compréhension orale, et le vote final t'a entraîné à justifier un choix en une phrase.

\begin{methode}[Auto-évaluation : ai-je réussi ?]
\begin{enumerate}
\item Mon exposé contient-il 8 à 10 phrases et au moins 5 mots du glossaire ?
\item Ma comparaison avec \zh{比} est-elle correcte (A \zh{比} B + adjectif, sans \zh{很}) ?
\item Ai-je exprimé une opinion personnelle avec \zh{我觉得} ?
\item Ma prononciation des tons était-elle claire (shìchǎng, guǎnggào, dǎzhé) ?
\item Ai-je présenté les deux cultures de façon factuelle et respectueuse ?
\end{enumerate}
\end{methode}

Pour aller plus loin, tu peux rédiger à l'écrit une version longue de ton exposé (15 phrases) ou enregistrer une capsule audio d'une minute pour l'exposition « Deux marchés, un monde » : c'est une excellente préparation à l'épreuve orale du baccalauréat.

\newpage

\coursec{Méthodes et exercices résolus}

\coursec{Leçon 33 — Éléments socioculturels : le marketing en Chine et au Cameroun}

\courssub{Introduction : qu'est-ce que le marketing ?}

\begin{definition}[Le marketing \zh{营销}]
Le marketing (\zh{营销} \emph{yíngxiāo}, litt. « gérer la vente ») désigne l'ensemble des techniques et actions visant à \cle{connaître les besoins des consommateurs}, \cle{créer de la valeur} (produits, services, image) et \cle{communiquer} pour vendre. En chinois, le verbe \zh{营} (gérer, exploiter) + \zh{销} (vendre, écouler) résume bien la démarche : \zh{做营销} (faire du marketing), \zh{营销策略} (stratégie marketing).
\end{definition}

\begin{textebox}[Dialogue : Qu'est-ce que le marketing ?]
\zh{A: 什么是营销？} (Shénme shì yíngxiāo?) \\
\zh{B: 营销就是让更多的人知道产品，想买产品，最后买下产品。} (Yíngxiāo jiùshì ràng gèng duō de rén zhīdào chǎnpǐn, xiǎng mǎi chǎnpǐn, zuìhòu mǎi xià chǎnpǐn.) \\
\zh{A: 所以广告、打折、直播带货都是营销？} (Suǒyǐ guǎnggào, dǎzhé, zhíbò dàihuò dōu shì yíngxiāo?) \\
\zh{B: 对，都是营销手段。} (Duì, dōu shì yíngxiāo shǒuduàn.) \\
\emph{Traduction :} A : Qu'est-ce que le marketing ? B : Le marketing, c'est faire en sorte que plus de gens connaissent le produit, veuillent l'acheter et finissent par l'acheter. A : Donc la pub, les soldes, la vente en direct sont du marketing ? B : Oui, ce sont tous des moyens de marketing.
\end{textebox}

\courssub{Le marketing au Cameroun : faits et pratiques}

\begin{textebox}[Le marketing au Cameroun : réalités du terrain]
\zh{在喀麦隆，传统营销依然很重要。} (Zài Kāmàilóng, chuántǒng yíngxiāo yīrán hěn zhòngyào.) \\
\zh{很多公司在广播、电视和路边大牌做广告。} (Hěnduō gōngsī zài guǎngbō, diànshì hé lùbiān dàpái zuò guǎnggào.) \\
\zh{市场像莫科洛市场，商人直接和顾客讲价，人情很重。} (Shìchǎng xiàng Mòkēluò shìchǎng, shāngrén zhíjiē hé gùkè jiǎngjià, rénqíng hěn zhòng.) \\
\zh{现在手机普及了，MTN 和 Orange 推送短信广告，WhatsApp 群发促销信息也很常见。} (Xiànzài shǒujī pǔjí le, MTN hé Orange tuīsòng duǎnxìn guǎnggào, WhatsApp qūnfā cùxiāo xìnxī yě hěn chángjiàn.) \\
\zh{但是，网购还不如中国普遍，很多人担心质量，喜欢看实物再买。} (Dànshì, wǎnggòu hái bùrú Zhōngguó pǔbiàn, hěnduō rén dānxīn zhìliàng, xǐhuan kàn shíwù zài mǎi.) \\
\emph{Traduction :} Au Cameroun, le marketing traditionnel reste important. Beaucoup d'entreprises font de la pub à la radio, la télé et sur les panneaux routiers. Dans les marchés comme Mokolo, les commerçants négocient directement avec les clients, la relation humaine est primordiale. Maintenant que les téléphones sont répandus, MTN et Orange envoient des SMS publicitaires, les promos par groupes WhatsApp sont aussi courantes. Mais le commerce en ligne n'est pas aussi répandu qu'en Chine, beaucoup craignent la qualité et préfèrent voir l'objet avant d'acheter.
\end{textebox}

\begin{aretenir}[Spécificités camerounaises]
\begin{itemize}
\item \cle{Supports dominants} : Radio (zone rurale), TV (urbain), Panneaux (\zh{路边大牌}), SMS / WhatsApp (mobile money \zh{Mobile Money} / \zh{MoMo}).
\item \cle{Relationnel} : \zh{讲价} (marchander) est une norme culturelle ; \zh{人情} (relations, faveurs) influence l'achat.
\item \cle{Freins au e-commerce} : Livraison lente/coûteuse hors Yaoundé/Douala, paiement à la livraison (\zh{货到付款}) préféré, confiance limitée.
\end{itemize}
\end{aretenir}

\courssub{Le marketing en Chine : faits et pratiques}

\begin{textebox}[Le marketing en Chine : l'ère du numérique]
\zh{中国是全球最大的网购市场，营销早已全面数字化。} (Zhōngguó shì quánqiú zuìdà de wǎnggòu shìchǎng, yíngxiāo zǎoyǐ quánmiàn shùzìhuà.) \\
\zh{微信、抖音、淘宝、京东是主要平台，直播带货极其火爆。} (Wēixìn, Dǒuyīn, Táobǎo, Jīngdōng shì zhǔyào píngtái, zhíbò dàihuò jíqí huǒbào.) \\
\zh{关键意见领袖 和 关键意见消费者 推荐产品，粉丝下单极快。} (Guānjiàn yìjiàn lǐngxiù hé guānjiàn yìjiàn xiāofèizhě tuījiàn chǎnpǐn, fěnsī xiàdān jíkuài.) \\
\zh{双十一、六一八 是全民购物节，打折力度大，预售、定金、尾款 玩法复杂。} (Shuāng shíyī, Liù yāo bā shì quánmín gòuwùjié, dǎzhé lìdù dà, yùshòu, dìngjīn, wěikuǎn wánfǎ fùzá.) \\
\zh{公司重视大数据分析，精准推送广告，甚至「一千人一千面」。} (Gōngsī zhòngshì dàshùjù fēnxī, jīngzhǔn tuīsòng guǎnggào, shènzhì « yīqiān rén yīqiān miàn ».) \\
\emph{Traduction :} La Chine est le plus grand marché mondial d'achat en ligne, le marketing est déjà entièrement numérisé. WeChat, Douyin, Taobao, JD sont les plateformes principales, la vente en direct (livestreaming) est explosif. Les KOL et KOC recommandent les produits, les fans commandent très vite. Le 11.11 et le 6.18 sont des fêtes nationales d'achat, grosses réductions, prévente, acompte, solde : mécaniques complexes. Les entreprises valorisent le big data pour pousser des pubs ultra-ciblées, jusqu'au « mille personnes, mille visages » (personnalisation totale).
\end{textebox}

\begin{savaistu}[Vocabulaire spécifique Chine : les grands rendez-vous]
\begin{itemize}
\item \zh{双十一} \emph{Shuāng Shíyī} (11 novembre) : « Fête des célibataires » → plus grand jour de soldes au monde.
\item \zh{六一八} \emph{Liù Yāo Bā} (18 juin) : Anniversaire de JD.com → 2e grand rendez-vous.
\item \zh{直播带货} \emph{zhíbò dàihuò} : Vente en direct (livestreaming e-commerce). Animateur = \zh{主播} \emph{zhǔbō}.
\item \zh{种草} \emph{zhǒngcǎo} (litt. « planter l'herbe ») : Recommander un produit pour donner envie (influence marketing).
\item \zh{拔草} \emph{bácǎo} : Acheter le produit recommandé (« arracher l'herbe »).
\end{itemize}
\end{savaistu}

\courssub{Vocabulaire associé : glossaire et caractères}

\begin{center}
\begin{tabularx}{\linewidth}{|X|X|X|X|}
\hline
\rowcolor{popblueL} \textbf{Caractères} & \textbf{Pinyin} & \textbf{Nature} & \textbf{Traduction} \\ \hline
营销 & yíngxiāo & n./v. & marketing ; faire du marketing \\ \hline
广告 & guǎnggào & n. & publicité, annonce \\ \hline
做广告 & zuò guǎnggào & v. & faire de la publicité \\ \hline
推销 & tuīxiāo & v. & promouvoir, démarcher, vendre activement \\ \hline
打折 & dǎzhé & v. & faire une réduction, soldes \\ \hline
打八折 & dǎ bā zhé & phr. & 20\% de réduction (payer 80\%) \\ \hline
顾客 & gùkè & n. & client (clientèle) \\ \hline
消费者 & xiāofèizhě & n. & consommateur \\ \hline
商人 & shāngrén & n. & commerçant, marchand \\ \hline
公司 & gōngsī & n. & entreprise, société \\ \hline
产品 & chǎnpǐn & n. & produit \\ \hline
质量 & zhìliàng & n. & qualité \\ \hline
价格 & jiàgé & n. & prix \\ \hline
讲价 & jiǎngjià & v. & marchander, négocier le prix \\ \hline
网购 & wǎnggòu & v./n. & achat en ligne, e-commerce \\ \hline
网站 & wǎngzhàn & n. & site web \\ \hline
平台 & píngtái & n. & plateforme (numérique) \\ \hline
社交媒体 & shèjiāo méitǐ & n. & réseaux sociaux \\ \hline
直播 & zhíbò & v./n. & direct (vidéo), streaming \\ \hline
带货 & dàihuò & v. & vendre des produits (en direct) \\ \hline
手机 & shǒujī & n. & téléphone portable \\ \hline
应用 / 软件 & yìngyòng / ruǎnjiàn & n. & application, logiciel \\ \hline
市场 & shìchǎng & n. & marché (physique ou économique) \\ \hline
策略 & cèlüè & n. & stratégie \\ \hline
目的 & mùdì & n. & but, objectif \\ \hline
方式 & fāngshì & n. & manière, mode, méthode \\ \hline
普遍 & pǔbiàn & adj. & répandu, généralisé, courant \\ \hline
方便 & fāngbiàn & adj. & pratique, commode \\ \hline
真实 & zhēnshí & adj. & vrai, authentique, honnête \\ \hline
创意 & chuàngyì & n./adj. & créativité ; créatif \\ \hline
虽然 & suīrán & conj. & bien que, quoique \\ \hline
但是 & dànshì & conj. & mais, cependant \\ \hline
因为 & yīnwèi & conj. & parce que \\ \hline
所以 & suǒyǐ & conj. & donc, c'est pourquoi \\ \hline
觉得 & juéde & v. & penser, trouver (avis) \\ \hline
认为 & rènwéi & v. & considérer, estimer \\ \hline
应该 & yīnggāi & v. mod. & devoir, devrait \\ \hline
越来越 & yuè lái yuè & adv. & de plus en plus \\ \hline
\end{tabularx}
\end{center}

\begin{methode}[Mémoriser et réemployer le lexique du marketing]
\begin{enumerate}
\item \cle{Apprendre par familles de mots} : Acteurs (\zh{公司}, \zh{顾客}, \zh{商人}, \zh{消费者}) — Actions (\zh{做广告}, \zh{推销}, \zh{打折}, \zh{讲价}, \zh{直播带货}) — Supports (\zh{广告}, \zh{网站}, \zh{平台}, \zh{社交媒体}, \zh{手机软件}) — Indicateurs (\zh{质量}, \zh{价格}, \zh{目的}, \zh{方式})。
\item \cle{Apprendre chaque nom avec son classificateur} : \zh{一家公司} (entreprise), \zh{一则广告} (annonce), \zh{一位顾客} (client - respect), \zh{一个商人}, \zh{两个国家} (\textbf{jamais} \zh{二个国家}), \zh{一种产品}, \zh{一个市场}。
\item \cle{Vérifier les tons dans les collocations} : \zh{做广告} (zuò guǎnggào, 4-3-4), \zh{打折} (dǎ zhé, 3-2), \zh{网上购物} (wǎngshàng gòuwù, 3-4 4-4), \zh{直播带货} (zhíbò dàihuò, 2-4 4-4).
\item \cle{Réemployer dans une phrase personnelle} : \zh{我妈妈在市场卖菜，她的顾客很多，不讲价不行。}
\end{enumerate}
\end{methode}

\begin{exoresolu}[Vocabulaire : texte à trous]
Énoncé — Complète avec les mots suivants : \zh{广告、顾客、打折、网站、推销}.\par
(1) 这家\zh{公司}的\zh{\_\_\_\_}很有意思，很多人都喜欢。 \\
(2) 商店\zh{\_\_\_\_}的时候，东西比较便宜。 \\
(3) 服务员对\zh{\_\_\_\_}很热情。 \\
(4) 现在年轻人在\zh{\_\_\_\_}上买东西。 \\
(5) 他想\zh{\_\_\_\_}自己的新产品。
\tcblower
Solution détaillée :
\begin{enumerate}
\item (1) \zh{广告} (guǎnggào) : l'indice « 很有意思 » qualifie la pub. \zh{这家公司的广告很有意思，很多人都喜欢。}
\item (2) \zh{打折} (dǎzhé) : l'indice « 比较便宜 » (moins cher). Attention : \zh{折} se prononce \emph{zhé} (2e ton) dans \zh{打折}. \zh{商店打折的时候，东西比较便宜。}
\item (3) \zh{顾客} (gùkè) : c'est envers les clients que le serveur est \zh{热情} (chaleureux). \zh{服务员对顾客很热情。}
\item (4) \zh{网站} (wǎngzhàn) : structure \zh{在} + lieu + \zh{上} + verbe. \zh{现在年轻人在网站上买东西。}
\item (5) \zh{推销} (tuīxiāo) : verbe transitif, objet \zh{自己的新产品}. \zh{想推销} = a l'intention de promouvoir. \zh{他想推销自己的新产品。}
\end{enumerate}
\end{exoresolu}

\courssub{Comparer et donner son avis en chinois}

\begin{propriete}[Structures de comparaison : \zh{比}, \zh{跟……一样}, \zh{没有}, \zh{不如}]
\begin{center}
\begin{tabularx}{\linewidth}{|X|X|X|}
\hline
\rowcolor{popblueL} \textbf{Type} & \textbf{Structure} & \textbf{Exemple (Caractères + Pinyin + Traduction)} \\ \hline
Supériorité & A \zh{比} B + Adj (pas de \zh{更}) & \zh{中国的网购比喀麦隆普遍。} (Zhōngguó de wǎnggòu bǐ Kāmàilóng pǔbiàn.) Le e-commerce en Chine est plus répandu qu'au Cameroun. \\ \hline
Égalité & A \zh{跟/和} B \zh{一样} + Adj & \zh{两个国家的广告一样有创意。} (Liǎng ge guójiā de guǎnggào yíyàng yǒu chuàngyì.) Les pubs des deux pays sont aussi créatives. \\ \hline
Infériorité (négation) & A \zh{没有} B + Adj & \zh{喀麦隆的网购没有中国普遍。} (Kāmàilóng de wǎnggòu méiyǒu Zhōngguó pǔbiàn.) Le e-commerce au Cameroun n'est pas aussi répandu qu'en Chine. \\ \hline
Infériorité (formel) & A \zh{不如} B + Adj & \zh{喀麦隆的物流不如中国快。} (Kāmàilóng de wùliú bùrú Zhōngguó kuài.) La logistique au Cameroun n'est pas aussi rapide qu'en Chine. \\ \hline
\end{tabularx}
\end{center}
\emph{Rappel :} \zh{比} + Adj suffit. \zh{更} / \zh{最} sont interdits après \zh{比}. Classificateur \zh{两} (liǎng) pour « deux » + classif., \zh{二} (èr) seulement pour compter ou numéros.
\end{propriete}

\begin{propriete}[Ordre des mots : Complément de lieu \zh{在} + Lieu + \zh{上/里} + Verbe]
En chinois, le lieu de l'action se place \textbf{AVANT} le verbe.
\end{propriete}

\begin{exemplebox}[Exemples]
\zh{在中国，很多人\underline{在网上}买东西。} (En Chine, beaucoup achètent \underline{sur Internet}.) \\
\zh{在喀麦隆，公司\underline{在广播里}做广告。} (Au Cameroun, les entreprises font de la pub \underline{à la radio}.) \\
\zh{*错误：} \zh{很多人买东西在网上。} (Calque du français « achètent sur Internet » — \textbf{interdit}).
\end{exemplebox}


\begin{propriete}[Concession et cause : paires obligatoires]
Contrairement au français, les deux termes s'emploient \textbf{ensemble} dans la phrase complexe.
\begin{center}
\begin{tabularx}{\linewidth}{|X|X|}
\hline
\rowcolor{popblueL} \zh{虽然} … \zh{但是/却} … & Bien que … , (pourtant) … \\ \hline
\zh{因为} … \zh{所以} … & Parce que … , donc … \\ \hline
\end{tabularx}
\end{center}
\end{propriete}

\begin{exemplebox}[Exemples]
\zh{虽然方式不一样，但是目的都一样。} (Bien que les méthodes diffèrent, le but est le même.) \\
\zh{因为网购方便，所以我喜欢。} (Parce que c'est pratique, je l'aime.)
\end{exemplebox}


\begin{methode}[Donner son avis structuré (Expression orale/écrite)]
\begin{enumerate}
\item \cle{Amorce} : \zh{我觉得……} / \zh{我认为……} / \zh{对我来说，……}
\item \cle{Justification} : \zh{因为……所以……} (garde les deux).
\item \cle{Nuance} : \zh{但是/不过} + contre-argument.
\item \cle{Bilan / Suggestion} : \zh{所以我觉得……} / \zh{公司应该……} / \zh{建议……}
\item \cle{Marqueurs de degré} : \zh{越来越} + Adj (\zh{网购越来越普遍}) ; \zh{比较} / \zh{很} / \zh{特别}.
\end{enumerate}
\end{methode}

\begin{exoresolu}[Expression d'un point de vue]
Énoncé — Réponds en 4-5 phrases : \zh{你觉得广告有用吗？为什么？}
\tcblower
Solution détaillée — Modèle de réponse :\par
\zh{我觉得广告很有用。} (Wǒ juéde guǎnggào hěn yǒuyòng.)\par
\zh{因为广告让我们知道新产品，所以对公司和顾客都很重要。} (Yīnwèi guǎnggào ràng wǒmen zhīdào xīn chǎnpǐn, suǒyǐ duì gōngsī hé gùkè dōu hěn zhòngyào.)\par
\zh{但是，有些广告不太真实，我们不应该都相信。} (Dànshì, yǒuxiē guǎnggào bú tài zhēnshí, wǒmen bù yīnggāi dōu xiāngxìn.)\par
\zh{所以，我觉得我们应该小心地看广告。} (Suǒyǐ, wǒ juéde wǒmen yīnggāi xiǎoxīn de kàn guǎnggào.)\par
\emph{Points clés :} \zh{让} (faire que) + COI + V ; \zh{对……很重要} ; \zh{有些} (indéfini, sans classif.) ; \zh{应该} + V ; \zh{小心地} (adv \zh{地}) + V.
\end{exoresolu}

\begin{exoresolu}[Comparaison guidée]
Énoncé — Écris trois phrases correctes avec : (1) \zh{网购 / 中国 / 普遍 / 喀麦隆} ; (2) \zh{广告 / 两个国家 / 有创意} ; (3) \zh{虽然 / 方式 / 不一样 / 目的 / 一样}.
\tcblower
Solution détaillée :
\begin{enumerate}
\item \zh{中国的网购比喀麦隆普遍。} (Supériorité \zh{比} ; pas de \zh{更}).
\item \zh{两个国家的广告一样有创意。} (Égalité \zh{一样} ; \zh{两个} + classif.).
\item \zh{虽然两个国家的方式不一样，但是目的都一样。} (Concession \zh{虽然……但是……}).
\end{enumerate}
\end{exoresolu}

\courssub{Questions de réflexion et préparation du mini-exposé}

\begin{methode}[Répondre aux questions de compréhension]
\begin{enumerate}
\item \cle{Lis les questions AVANT le texte} : repère \zh{谁} (qui), \zh{什么} (quoi), \zh{哪儿} (où), \zh{为什么} (pourquoi), \zh{怎么样} (comment).
\item \cle{Surline les mots-clés du thème} dans le texte (\zh{市场}, \zh{广告}, \zh{网购}, \zh{顾客}, \zh{手机}, \zh{直播}).
\item \cle{Localise la réponse} : ordre chronologique (Q1 → début, Q3 → fin).
\item \cle{Rédige une phrase complète} : reprends le sujet de la question ; à \zh{为什么} réponds par \zh{因为……}.
\end{enumerate}
\end{methode}

\begin{exoresolu}[Compréhension de texte comparatif]
Énoncé — Lis le texte et réponds en phrases complètes.\par
\zh{在中国，很多公司用手机软件做广告，年轻人常常在网上买东西。在喀麦隆，公司常常在广播和电视上做广告，很多人还是喜欢去市场买东西，因为可以看到产品，还可以讲价。虽然两个国家的方式不一样，但是公司的目的都一样：让更多的顾客买他们的产品。}\par
(1) \zh{中国的公司用什么做广告？} \\
(2) \zh{喀麦隆人为什么喜欢去市场买东西？} \\
(3) \zh{两个国家的公司的目的是什么？}
\tcblower
Solution détaillée :
\begin{enumerate}
\item \zh{中国的公司用手机软件做广告。} (Extraction directe : \zh{用} + outil + \zh{做广告}).
\item \zh{因为在市场可以看到产品，还可以讲价。} (\zh{因为} obligatoire ; \zh{还} = de plus / aussi).
\item \zh{他们的目的是让更多的顾客买他们的产品。} (\zh{更多} = plus nombreux ; \zh{让} + personne + V = faire en sorte que).
\end{enumerate}
\end{exoresolu}

\begin{methode}[Structurer un mini-exposé oral (2-3 min)]
\begin{enumerate}
\item \cle{Salutation + Sujet} : \zh{大家好，今天我要谈一谈中国和喀麦隆的营销。}
\item \cle{Développement symétrique} : \zh{先说说中国……} (ex: \zh{直播带货、双十一}) → \zh{再说说喀麦隆……} (ex: \zh{莫科洛市场、WhatsApp群发}).
\item \cle{Comparaison synthétique} : \zh{中国的网购比喀麦隆普遍，但是喀麦隆的市场更有生活气息。} + \zh{虽然方式不一样，但是目的都一样。}
\item \cle{Avis personnel} : \zh{我觉得……因为……所以……}
\item \cle{Remerciement} : \zh{谢谢大家！}
\end{enumerate}
\emph{Critères oraux :} Tons justes (\zh{mǎi} vs \zh{mài}), débit lent, regard public, notes \textbf{interdites} (seuls mots-clés autorisés).
\end{methode}

\begin{exoresolu}[Mini-exposé : remise en ordre et justification]
Énoncé — Remets dans l'ordre : (a) \zh{在喀麦隆，广播广告和市场很重要。} (b) \zh{谢谢大家！} (c) \zh{大家好，今天我要谈一谈中国和喀麦隆的营销。} (d) \zh{在中国，很多公司在网上做广告。} (e) \zh{虽然方式不一样，但是目的都一样。} (f) \zh{我觉得两个国家的营销都很有意思。}
\tcblower
Ordre : (c) → (d) → (a) → (e) → (f) → (b).\par
Justification : (c) Introduction ; (d)/(a) Développement symétrique (ordre annoncé : Chine puis Cameroun) ; (e) Synthèse comparative (\zh{虽然……但是……}) ; (f) Avis personnel (\zh{我觉得}) ; (b) Clôture rituelle.
\end{exoresolu}

\begin{experience}[Activité de classe : Débat « Marché physique vs Live-streaming »]
\cle{Consigne} : En binômes. L'un défend le modèle camerounais (marché Mokolo, contact humain, \zh{讲价}), l'autre le modèle chinois (Douyin/Taobao Live, \zh{直播带货}, rapidité, \zh{种草}). Utilisez : \zh{我觉得……因为……但是……所以……}. Durée : 5 min préparation, 3 min débat par binôme. L'enseignant note : justesse tons, structures \zh{比/没有/虽然}, vocabulaire thème.
\end{experience}

\begin{attention}[Les 5 erreurs qui coûtent des points au Bac]
\begin{enumerate}
\item \cle{Confondre \zh{mǎi} (买, acheter) et \zh{mài} (卖, vendre)} : même syllabe, tons opposés (3 vs 4). Caractères : \zh{买} (pas de trait horizontal en haut) vs \zh{卖} (avec \zh{十} « dix/croix » en haut = « sortir/ vendre »). \zh{公司卖，顾客买。}
\item \cle{Ajouter \zh{更} après \zh{比}} : \textbf{Faux} \zh{比喀麦隆更普遍} → \textbf{Juste} \zh{比喀麦隆普遍}. \zh{比} marque déjà le comparatif.
\item \cle{Mal placer le complément de lieu \zh{在}} : \textbf{Faux} \zh{我买东西在网上} → \textbf{Juste} \zh{我在网上买东西}. Règle : Sujet + \zh{在} Lieu + Verbe.
\item \cle{Oublier le classificateur} : \textbf{Faux} \zh{两个国家} (correct) vs \zh{二个国家} (faux) ; \zh{一则广告} (pas \zh{一个广告} pour une annonce), \zh{一位顾客} (respect).
\item \cle{Confondre \zh{的} (nom) et \zh{地} (verbe)} : \zh{真实的广告} (adj + \zh{的} + N) vs \zh{小心地看} (adv + \zh{地} + V). Et rappel : \zh{因为……所以……} / \zh{虽然……但是……} sont inséparables en chinois.
\end{enumerate}
\end{attention}

\begin{exercice}[Entraînement complet — Type Bac]
\begin{enumerate}
\item \cle{Traduction (Version) :} Traduis en français : \zh{虽然喀麦隆的网购没有中国普遍，但是手机广告越来越多，因为公司想推销产品。}
\item \cle{Traduction (Thème) :} Traduis en chinois (caractères + pinyin) : « Au Cameroun, les gens aiment aller au marché pour marchander, mais en Chine, les jeunes préfèrent acheter sur leur portable. »
\item \cle{Expression écrite :} Écris un paragraphe de 60-80 caractères : \zh{谈谈你对“直播带货”的看法。} (Donne ton avis sur la vente en direct : utilité, risques, ton expérience).
\item \cle{Grammaire :} Complète avec \zh{比 / 跟……一样 / 没有 / 不如 / 虽然 / 但是 / 因为 / 所以} : \\
(1) 这家公司\_\_\_\_那家公司大。 \\
(2) \_\_\_\_今天下雨，\_\_\_\_我们去公园。 \\
(3) 喀麦隆的物流\_\_\_\_中国快。 \\
(4) 两个平台\_\_\_\_\_\_\_\_方便。
\end{enumerate}
\end{exercice}

\begin{corrige}[Exercice complet]
\begin{enumerate}
\item \emph{Version :} « Bien que le commerce en ligne au Cameroun ne soit pas aussi répandu qu'en Chine, les publicités sur mobile sont de plus en plus nombreuses, car les entreprises veulent promouvoir leurs produits. »
\item \emph{Thème :} \zh{在喀麦隆，人们喜欢去市场讲价，但是在中国，年轻人比较喜欢用手机买东西。} (Zài Kāmàilóng, rénmen xǐhuan qù shìchǎng jiǎngjià, dànshì zài Zhōngguó, niánqīng rén bǐjiào xǐhuan yòng shǒujī mǎi dōngxi.)
\item \emph{Expression (modèle) :} \zh{我觉得直播带货很方便，因为可以直接看产品效果，价格有时也便宜。但是，有些主播不诚实，产品质量不好，退货还麻烦。所以，我们看直播买东西要小心，不要只听广告，要看评价。} (Env. 75 caractères).
\item \emph{Grammaire :} (1) \zh{比} ; (2) \zh{虽然} … \zh{但是} ; (3) \zh{不如} (ou \zh{没有}) ; (4) \zh{跟……一样}.
\end{enumerate}
\end{corrige}

\newpage

\coursec{Exercices}

\courssub{Je vérifie mes connaissances}

\begin{exercice}[QCM : le marketing en Chine et au Cameroun]
Pour chaque question, choisis la bonne réponse (a, b ou c).

\begin{enumerate}
\item En Chine, le « Double 11 » (\zh{双十一}) est :
\begin{enumerate}
\item une fête traditionnelle de la famille ;
\item une grande journée de soldes en ligne ;
\item une fête nationale chinoise.
\end{enumerate}
\item \zh{直播带货} (zhíbō dàihuò) signifie :
\begin{enumerate}
\item vendre des produits en direct sur Internet ;
\item regarder un match de football en direct ;
\item apprendre le chinois en ligne.
\end{enumerate}
\item Au Cameroun, beaucoup de petits commerçants font leur publicité surtout :
\begin{enumerate}
\item par des affiches, la radio et le bouche-à-oreille ;
\item uniquement par des applications chinoises ;
\item uniquement à la télévision nationale.
\end{enumerate}
\item \zh{打折} (dǎzhé) veut dire :
\begin{enumerate}
\item augmenter les prix ;
\item faire une réduction ;
\item fermer le magasin.
\end{enumerate}
\end{enumerate}
\end{exercice}

\begin{exercice}[Vrai ou faux ? Justifie ta réponse]
Dis si chaque phrase est vraie ou fausse, puis justifie en une phrase en français.

\begin{enumerate}
\item En Chine, le paiement par téléphone portable (WeChat Pay, Alipay) est très répandu.
\item Au Cameroun, personne n'achète jamais en ligne.
\item Le « Double 11 » a lieu le 11 novembre.
\item Les vendeurs en direct (\zh{主播}) présentent des produits et répondent aux questions des clients.
\item Au Cameroun, les marchés traditionnels ont complètement disparu à cause d'Internet.
\end{enumerate}
\end{exercice}

\begin{exercice}[Vocabulaire : appariement]
Associe chaque mot chinois à sa traduction française.

\begin{center}
\begin{tabularx}{\linewidth}{|X|X|X|X|}
\hline
\rowcolor{popblueL} N° & Chinois & N° & Français \\
\hline
1 & \zh{广告} (guǎnggào) & a & le client \\
\hline
2 & \zh{顾客} (gùkè) & b & la marque \\
\hline
3 & \zh{品牌} (pǐnpái) & c & la publicité \\
\hline
4 & \zh{打折} (dǎzhé) & d & acheter en ligne \\
\hline
5 & \zh{网购} (wǎnggòu) & e & faire une réduction \\
\hline
6 & \zh{市场} (shìchǎng) & f & le prix \\
\hline
7 & \zh{价格} (jiàgé) & g & le marché \\
\hline
8 & \zh{直播} (zhíbō) & h & la qualité \\
\hline
9 & \zh{质量} (zhìliàng) & i & le direct (diffusion en ligne) \\
\hline
10 & \zh{便宜} (piányi) & j & bon marché \\
\hline
\end{tabularx}
\end{center}
\end{exercice}

\begin{exercice}[Pinyin et caractères]
Écris le pinyin avec les tons sous chaque caractère, puis donne la traduction française.

\begin{enumerate}
\item \zh{广告}
\item \zh{顾客}
\item \zh{打折}
\item \zh{直播}
\item \zh{品牌}
\item \zh{网购}
\end{enumerate}
\end{exercice}

\courssub{Je m'entraîne}

\begin{exercice}[Phrases à trous]
Complète chaque phrase avec le mot qui convient : \zh{广告}、\zh{顾客}、\zh{打折}、\zh{直播}、\zh{品牌}、\zh{网购}。

\begin{enumerate}
\item 双十一的时候，很多商店都\_\_\_\_\_\_，东西很便宜。
\item 这个\_\_\_\_\_\_很有名，年轻人都喜欢。
\item 现在，很多中国人喜欢\_\_\_\_\_\_，不去商店。
\item 电视上的\_\_\_\_\_\_太多了，我不想看。
\item 这位主播在\_\_\_\_\_\_里卖衣服，很多人看。
\item 商店对\_\_\_\_\_\_很友好，所以他们常常来买东西。
\end{enumerate}
\end{exercice}

\begin{exercice}[Transforme avec 比]
Transforme chaque paire de phrases en une phrase comparative avec \zh{比}。

Modèle : 网购方便。传统市场不方便。→ 网购比传统市场方便。

\begin{enumerate}
\item 中国的网购很多。喀麦隆的网购不太多。
\item 这件衣服贵。那件衣服不贵。
\item 网上买东西快。去市场买东西不快。
\item 这个品牌有名。那个品牌没有名。
\item 今天的价格便宜。昨天的价格不便宜。
\end{enumerate}
\end{exercice}

\begin{exercice}[Remets les mots en ordre]
Remets les mots dans le bon ordre pour faire une phrase correcte, puis traduis-la en français.

\begin{enumerate}
\item 比 / 网购 / 方便 / 去市场 / 。
\item 很多 / 在 / 中国人 / 网上 / 东西 / 买 / 。
\item 顾客 / 这个 / 的 / 商店 / 很多 / 有 / 。
\item 直播 / 我 / 看 / 喜欢 / 卖东西的 / 。
\item 双十一 / 打折 / 的时候 / 商店 / 都 / 。
\end{enumerate}
\end{exercice}

\begin{exercice}[Traduction (thème)]
Traduis les phrases suivantes en chinois (caractères obligatoires).

\begin{enumerate}
\item Au Cameroun, beaucoup de gens achètent au marché ; en Chine, beaucoup de gens achètent en ligne.
\item Pendant le « Double 11 », les magasins font des réductions.
\item Cette marque est très connue, mais elle est un peu chère.
\item Les vendeurs en direct présentent des produits et répondent aux questions des clients.
\item Je pense que la publicité à la radio est très utile au Cameroun.
\end{enumerate}
\end{exercice}

\courssub{Je me prépare au Bac}

\begin{exercice}[Compréhension écrite : le « Double 11 »]
Lis le texte suivant, puis réponds aux questions.

\begin{textebox}[双十一]
双十一是中国最大的购物节。每年的十一月十一日，很多网上商店都打折，东西比平时便宜很多。这一天，几亿人在网上买东西。有的人买衣服，有的人买手机，还有的人买吃的。现在，很多主播在直播里卖东西：他们介绍产品，回答顾客的问题。顾客看了直播，常常马上就买。双十一也让很多外国品牌进入了中国市场。\\
\textit{Shuāng shíyī shì Zhōngguó zuì dà de gòuwù jié. Měi nián de shíyī yuè shíyī rì, hěn duō wǎngshàng shāngdiàn dōu dǎzhé, dōngxi bǐ píngshí piányi hěn duō. Zhè yī tiān, jǐ yì rén zài wǎngshàng mǎi dōngxi. Yǒu de rén mǎi yīfu, yǒu de rén mǎi shǒujī, hái yǒu de rén mǎi chī de. Xiànzài, hěn duō zhǔbō zài zhíbō lǐ mài dōngxi : tāmen jièshào chǎnpǐn, huídá gùkè de wèntí. Gùkè kàn le zhíbō, chángcháng mǎshàng jiù mǎi. Shuāng shíyī yě ràng hěn duō wàiguó pǐnpái jìnrù le Zhōngguó shìchǎng.}\\
« Le Double 11 est la plus grande fête du shopping de Chine. Chaque année, le 11 novembre, de nombreuses boutiques en ligne font des réductions et les produits sont bien moins chers que d'habitude. Ce jour-là, des centaines de millions de personnes achètent en ligne. Certains achètent des vêtements, d'autres des téléphones, d'autres encore de la nourriture. Aujourd'hui, de nombreux animateurs vendent des produits en direct : ils présentent les produits et répondent aux questions des clients. Après avoir regardé le direct, les clients achètent souvent immédiatement. Le Double 11 a aussi permis à de nombreuses marques étrangères d'entrer sur le marché chinois. »
\end{textebox}

\textbf{Questions} (réponds en chinois, en phrases complètes) :

\textbf{A. Vrai ou faux} (recopie 对 ou 不对) :
\begin{enumerate}
\item 双十一是十一月十一日。
\item 双十一的时候，东西比平时贵。
\item 主播在直播里介绍产品。
\end{enumerate}

\textbf{B. Questions ouvertes} :
\begin{enumerate}
\item 双十一的时候，人们买什么东西？
\item 顾客看了直播，常常做什么？
\item 你觉得网购好不好？为什么？(réponds avec 我觉得……因为……)
\end{enumerate}
\end{exercice}

\begin{exercice}[Thème et version]
\textbf{A. Thème.} Traduis en chinois (caractères obligatoires) :

\begin{enumerate}
\item En Chine, le commerce en ligne est plus développé qu'au Cameroun.
\item Au Cameroun, les vendeurs du marché parlent directement avec les clients.
\item Cette marque de jus camerounais est bonne et pas chère.
\item Pendant les fêtes, les magasins font beaucoup de publicité.
\end{enumerate}

\textbf{B. Version.} Traduis en français :

\zh{在中国，很多年轻人喜欢在网上买东西。他们觉得网购比去商店方便，也比去商店便宜。但是，有些老人还是喜欢去市场，因为他们可以看看产品的质量。}

\textit{Zài Zhōngguó, hěn duō niánqīngrén xǐhuan zài wǎngshàng mǎi dōngxi. Tāmen juéde wǎnggòu bǐ qù shāngdiàn fāngbiàn, yě bǐ qù shāngdiàn piányi. Dànshì, yǒuxiē lǎorén háishi xǐhuan qù shìchǎng, yīnwèi tāmen kěyǐ kànkan chǎnpǐn de zhìliàng.}
\end{exercice}

\begin{exercice}[Expression écrite type Bac]
\textbf{Sujet :} 喀麦隆和中国的广告

Rédige un court texte de \textbf{6 à 8 phrases} (au moins 60 caractères chinois) comparant une pratique publicitaire au Cameroun et en Chine.

Consignes :
\begin{itemize}
\item utilise au moins une comparaison avec \zh{比} ;
\item utilise au moins deux mots du glossaire de la leçon (\zh{广告}、\zh{顾客}、\zh{品牌}、\zh{打折}、\zh{直播}、\zh{网购}…) ;
\item donne au moins une opinion personnelle avec \zh{我觉得……} ou \zh{我认为……} ;
\item termine par une conclusion d'une phrase.
\end{itemize}

\textbf{Expression orale (en binôme) :} jeu de rôle « vendeur en direct » (\zh{直播}). Un élève présente un produit camerounais (jus, tissu, arachides) en chinois : nom du produit, prix, qualité. L'autre élève joue le client et pose au moins trois questions (prix, qualité, réduction). Puis échangez les rôles.
\end{exercice}

\newpage

\coursec{Corrigés des exercices}

\begin{corrige}[Exercice 1 — QCM : le marketing en Chine et au Cameroun]
\begin{enumerate}
\item \textbf{b} — Le « Double 11 » (\zh{双十一}) est la plus grande journée de soldes en ligne en Chine (équivalent du Black Friday), organisée chaque 11 novembre par les plateformes comme Taobao, JD.com, Pinduoduo.
\item \textbf{a} — \zh{直播带货} (zhíbō dàihuò) = « diffusion en direct + apporter des marchandises » : vendre des produits en direct sur Internet (live‑shopping).
\item \textbf{a} — Au Cameroun, les petits commerçants privilégient encore l'affichage, la radio locale et le bouche‑à‑oreille ; le numérique progresse mais n'est pas exclusif.
\item \textbf{b} — \zh{打折} (dǎzhé) = « frapper la remise » → faire une réduction, soldes.
\end{enumerate}
\end{corrige}

\begin{corrige}[Exercice 2 — Vrai ou faux ? Justifie ta réponse]
\begin{enumerate}
\item \textbf{Vrai.} Le paiement mobile (WeChat Pay, Alipay) couvre quasi tous les commerces en Chine, des supermarchés aux étals de rue.
\item \textbf{Faux.} Le commerce en ligne existe au Cameroun (Jumia, Kaymu, groupes WhatsApp/Marketplace Facebook) et progresse vite, surtout chez les jeunes urbains.
\item \textbf{Vrai.} « Double 11 » = 11 novembre (11/11), date choisie pour ses « 1 » solitaires devenus fête du shopping.
\item \textbf{Vrai.} Les \zh{主播} (zhǔbō, animateurs/vendeurs en direct) présentent les produits, font des démonstrations, répondent aux questions en temps réel et lancent des coupons limités.
\item \textbf{Faux.} Les marchés traditionnels (Marché Central, Mokolo, Mfoundi, etc.) restent le cœur du commerce de proximité au Cameroun ; ils s'adaptent au numérique mais ne disparaissent pas.
\end{enumerate}
\end{corrige}

\begin{corrige}[Exercice 3 — Vocabulaire : appariement]
\begin{center}
\begin{tabularx}{\linewidth}{|X|X|X|X|}
\hline
\rowcolor{popblueL} N° & Chinois & N° & Français \\
\hline
1 & \zh{广告} (guǎnggào) & c & la publicité \\
\hline
2 & \zh{顾客} (gùkè) & a & le client \\
\hline
3 & \zh{品牌} (pǐnpái) & b & la marque \\
\hline
4 & \zh{打折} (dǎzhé) & e & faire une réduction \\
\hline
5 & \zh{网购} (wǎnggòu) & d & acheter en ligne \\
\hline
6 & \zh{市场} (shìchǎng) & g & le marché \\
\hline
7 & \zh{价格} (jiàgé) & f & le prix \\
\hline
8 & \zh{直播} (zhíbō) & i & le direct (diffusion en ligne) \\
\hline
9 & \zh{质量} (zhìliàng) & h & la qualité \\
\hline
10 & \zh{便宜} (piányi) & j & bon marché \\
\hline
\end{tabularx}
\end{center}
\end{corrige}

\begin{corrige}[Exercice 4 — Pinyin et caractères]
\begin{enumerate}
\item \zh{广告} — \textit{guǎnggào} — la publicité
\item \zh{顾客} — \textit{gùkè} — le client
\item \zh{打折} — \textit{dǎzhé} — faire une réduction
\item \zh{直播} — \textit{zhíbō} — le direct (diffusion en ligne)
\item \zh{品牌} — \textit{pǐnpái} — la marque
\item \zh{网购} — \textit{wǎnggòu} — acheter en ligne
\end{enumerate}
\end{corrige}

\begin{corrige}[Exercice 5 — Phrases à trous]
\begin{enumerate}
\item 双十一的时候，很多商店都\zh{打折}，东西很便宜。\\
\textit{Shuāng shíyī de shíhou, hěn duō shāngdiàn dōu dǎzhé, dōngxi hěn piányi.}
\item 这个\zh{品牌}很有名，年轻人都喜欢。\\
\textit{Zhège pǐnpái hěn yǒumíng, niánqīngrén dōu xǐhuan.}
\item 现在，很多中国人喜欢\zh{网购}，不去商店。\\
\textit{Xiànzài, hěn duō Zhōngguórén xǐhuan wǎnggòu, bù qù shāngdiàn.}
\item 电视上的\zh{广告}太多了，我不想看。\\
\textit{Diànshì shàng de guǎnggào tài duō le, wǒ bù xiǎng kàn.}
\item 这位主播在\zh{直播}里卖衣服，很多人看。\\
\textit{Zhè wèi zhǔbō zài zhíbō lǐ mài yīfu, hěn duō rén kàn.}
\item 商店对\zh{顾客}很友好，所以他们常常来买东西。\\
\textit{Shāngdiàn duì gùkè hěn yǒuhǎo, suǒyǐ tāmen chángcháng lái mǎi dōngxi.}
\end{enumerate}
\end{corrige}

\begin{corrige}[Exercice 6 — Transforme avec 比]
Structure : \cle{A + 比 + B + Adj} (A est plus Adj que B). Pas de \zh{很} après l'adjectif dans la comparative.
\begin{enumerate}
\item 中国的网购\zh{比}喀麦隆的网购\zh{多}。\\
\textit{Zhōngguó de wǎnggòu bǐ Kāmàilóng de wǎnggòu duō.}
\item 这件衣服\zh{比}那件衣服\zh{贵}。\\
\textit{Zhè jiàn yīfu bǐ nà jiàn yīfu guì.}
\item 网上买东西\zh{比}去市场买东西\zh{快}。\\
\textit{Wǎngshàng mǎi dōngxi bǐ qù shìchǎng mǎi dōngxi kuài.}
\item 这个品牌\zh{比}那个品牌\zh{有名}。\\
\textit{Zhège pǐnpái bǐ nàge pǐnpái yǒumíng.}
\item 今天的价格\zh{比}昨天的价格\zh{便宜}。\\
\textit{Jīntiān de jiàgé bǐ zuótiān de jiàgé piányi.}
\end{enumerate}
\end{corrige}

\begin{corrige}[Exercice 7 — Remets les mots en ordre]
\begin{enumerate}
\item \zh{网购比去市场方便。}\\
\textit{Wǎnggòu bǐ qù shìchǎng fāngbiàn.} — Le shopping en ligne est plus pratique qu'aller au marché.
\item \zh{很多中国人在网上买东西。}\\
\textit{Hěn duō Zhōngguórén zài wǎngshàng mǎi dōngxi.} — Beaucoup de Chinois achètent des choses en ligne.
\item \zh{这个商店有很多顾客。}\\
\textit{Zhège shāngdiàn yǒu hěn duō gùkè.} — Ce magasin a beaucoup de clients.
\item \zh{我喜欢看直播卖东西的。}\\
\textit{Wǒ xǐhuan kàn zhíbō mài dōngxi de.} — J'aime regarder les directs qui vendent des choses.
\item \zh{双十一的时候商店都打折。}\\
\textit{Shuāng shíyī de shíhou shāngdiàn dōu dǎzhé.} — Pendant le Double 11, les magasins font tous des réductions.
\end{enumerate}
\end{corrige}

\begin{corrige}[Exercice 8 — Traduction (thème)]
\begin{enumerate}
\item \zh{在喀麦隆，很多人在市场买东西；在中国，很多人在网上买东西。}\\
\textit{Zài Kāmàilóng, hěn duō rén zài shìchǎng mǎi dōngxi; zài Zhōngguó, hěn duō rén zài wǎngshàng mǎi dōngxi.}
\item \zh{双十一的时候，商店都打折。}\\
\textit{Shuāng shíyī de shíhou, shāngdiàn dōu dǎzhé.}
\item \zh{这个品牌很有名，但是有点儿贵。}\\
\textit{Zhège pǐnpái hěn yǒumíng, dànshì yǒudiǎnr guì.}
\item \zh{主播在直播里介绍产品，回答顾客的问题。}\\
\textit{Zhǔbō zài zhíbō lǐ jièshào chǎnpǐn, huídá gùkè de wèntí.}
\item \zh{我觉得在喀麦隆，广播广告很有用。}\\
\textit{Wǒ juéde zài Kāmàilóng, guǎngbō guǎnggào hěn yǒuyòng.}
\end{enumerate}
\end{corrige}

\begin{corrige}[Exercice 9 — Compréhension écrite : le « Double 11 »]
\textbf{A. Vrai ou faux}
\begin{enumerate}
\item \zh{对} — Le texte dit : « 每年的十一月十一日 ».
\item \zh{不对} — Le texte dit : « 东西比平时便宜很多 » (bien moins chers qu'habituellement), pas plus chers.
\item \zh{对} — Le texte dit : « 他们介绍产品 ».
\end{enumerate}

\textbf{B. Questions ouvertes}
\begin{enumerate}
\item \zh{双十一的时候，人们买衣服、手机、吃的东西等。}\\
\textit{Shuāng shíyī de shíhou, rénmen mǎi yīfu, shǒujī, chī de dōngxi děng.}
\item \zh{顾客看了直播，常常马上就买。}\\
\textit{Gùkè kàn le zhíbō, chángcháng mǎshàng jiù mǎi.}
\item \zh{我觉得网购好，因为方便、便宜，还有很多选择。} (ou : \zh{我觉得网购不好，因为看不见质量，可能有假货。})\\
\textit{Wǒ juéde wǎnggòu hǎo, yīnwèi fāngbiàn, piányi, hái yǒu hěn duō xuǎnzé.}
\end{enumerate}
\end{corrige}

\begin{corrige}[Exercice 10 — Thème et version]
\textbf{A. Thème}
\begin{enumerate}
\item \zh{在中国，网购比喀麦隆更发达。}\\
\textit{Zài Zhōngguó, wǎnggòu bǐ Kāmàilóng gèng fādá.}
\item \zh{在喀麦隆，市场的卖家直接跟顾客说话。}\\
\textit{Zài Kāmàilóng, shìchǎng de mài jiā zhíjiē gēn gùkè shuōhuà.}
\item \zh{这个喀麦隆果汁品牌很好，也不贵。}\\
\textit{Zhège Kāmàilóng guǒzhī pǐnpái hěn hǎo, yě bù guì.}
\item \zh{过节的时候，商店做很多广告。}\\
\textit{Guòjié de shíhou, shāngdiàn zuò hěn duō guǎnggào.}
\end{enumerate}

\textbf{B. Version}
« En Chine, beaucoup de jeunes aiment acheter des choses en ligne. Ils trouvent que le shopping en ligne est plus pratique qu'aller au magasin, et aussi moins cher qu'aller au magasin. Mais certaines personnes âgées aiment encore aller au marché, parce qu'elles peuvent voir la qualité des produits. »
\end{corrige}

\begin{corrige}[Exercice 11 — Expression écrite type Bac]
\textbf{Barème indicatif (sur 10 pts) :}
\begin{itemize}
\item Respect des consignes (longueur 6–8 phrases, ≥ 60 caractères, \zh{比}, 2 mots glossaire, opinion, conclusion) : 4 pts
\item Correction linguistique (grammaire, vocabulaire, ordre des mots) : 3 pts
\item Cohérence, pertinence du comparatif Cameroun/Chine : 2 pts
\item Qualité de l'expression (richesse, connecteurs) : 1 pt
\end{itemize}

\textbf{Proposition de rédaction modèle (≈ 85 caractères) :}
\begin{textebox}[Modèle : 喀麦隆和中国的广告]
喀麦隆和中国的广告不一样。在喀麦隆，广告常常在广播、电视和海报上，顾客听广播知道品牌。在中国，直播广告很流行，主播在网上卖东西，很多人看。我觉得中国的直播广告比喀麦隆的传统广告快，因为顾客马上能买。但是喀麦隆的广播广告也很有用，因为很多人听广播。我认为两个国家的广告都有特点。以后，喀麦隆也可以多用直播卖东西。\\
\textit{Kāmàilóng hé Zhōngguó de guǎnggào bù yīyàng. Zài Kāmàilóng, guǎnggào chángcháng zài guǎngbō, diànshì hé hǎibào shàng, gùkè tīng guǎngbō zhīdào pǐnpái. Zài Zhōngguó, zhíbō guǎnggào hěn liúxíng, zhǔbō zài wǎngshàng mài dōngxi, hěn duō rén kàn. Wǒ juéde Zhōngguó de zhíbō guǎnggào bǐ Kāmàilóng de chuántǒng guǎnggào kuài, yīnwèi gùkè mǎshàng néng mǎi. Dànshì Kāmàilóng de guǎngbō guǎnggào yě hěn yǒuyòng, yīnwèi hěn duō rén tīng guǎngbō. Wǒ rènwéi liǎng ge guójiā de guǎnggào dōu yǒu tèdiǎn. Yǐhòu, Kāmàilóng yě kěyǐ duō yòng zhíbō mài dōngxi.}
\end{textebox}

\textbf{Points de vigilance (erreurs fréquentes) :}
\begin{itemize}
\item Oublier \zh{比} ou placer l'adjectif avant \zh{比} (*\zh{方便比去市场网购} → faux).
\item Confondre \zh{广告} (publicité) et \zh{广播} (radio) ; \zh{主播} (animateur live) et \zh{顾客} (client).
\item Omettre le sujet dans la phrase d'opinion (*\zh{觉得网购好} → préférer \zh{我觉得网购好}).
\item Ne pas terminer par une vraie conclusion (une phrase qui clôture le propos).
\item Caractères mal formés : \zh{品} (pǐn, trois bouches) ≠ \zh{品} ; \zh{牌} (pái, radical 片) ; \zh{直} (zhí, dix + œil + trait) ; \zh{播} (bō, radical 水 + radical 扁).
\end{itemize}

\textbf{Expression orale (jeu de rôle « vendeur en direct ») — grille d'évaluation :}
\begin{itemize}
\item Vendeur : nom du produit (\zh{这个产品叫……}), prix (\zh{价格是……元}), qualité (\zh{质量很好/纯天然}), appel à l'action (\zh{欢迎下单}).
\item Client : au moins 3 questions — \zh{多少钱？} / \zh{有打折吗？} / \zh{质量怎么样？} / \zh{怎么发货？}.
\item Prononciation des tons (surtout \zh{直播 zhíbō}, \zh{打折 dǎzhé}, \zh{便宜 piányi}), fluidité, interaction naturelle.
\end{itemize}
\end{corrige}

\newpage

\coursec{Fiche bilan}

\begin{popfigure}
\begin{center}\tcbox[colback=popdark,colframe=popdark,coltext=white,arc=9pt,boxsep=4pt,left=6pt,right=6pt]{\bfseries\small Le marketing Cameroun / Chine}\end{center}
\vspace{-2pt}
\begin{tcbraster}[raster columns=3,raster equal height=rows,raster column skip=6pt,raster row skip=6pt,enhanced,blankest]
\begin{tcolorbox}[enhanced,colback=popblue,colframe=popblue,coltext=white,boxrule=0.9pt,arc=3pt,left=4pt,right=4pt,top=3pt,bottom=3pt,boxsep=1pt,halign=left]\bfseries\scriptsize Définition 市场 shìchǎng 营销 yíngxiāo\end{tcolorbox}
\begin{tcolorbox}[enhanced,colback=poporange,colframe=poporange,coltext=white,boxrule=0.9pt,arc=3pt,left=4pt,right=4pt,top=3pt,bottom=3pt,boxsep=1pt,halign=left]\bfseries\scriptsize Marketing au Cameroun marchés, radio, bouche à oreille\end{tcolorbox}
\begin{tcolorbox}[enhanced,colback=popgreen,colframe=popgreen,coltext=white,boxrule=0.9pt,arc=3pt,left=4pt,right=4pt,top=3pt,bottom=3pt,boxsep=1pt,halign=left]\bfseries\scriptsize Marketing en Chine WeChat, live, e-commerce\end{tcolorbox}
\begin{tcolorbox}[enhanced,colback=poppurple,colframe=poppurple,coltext=white,boxrule=0.9pt,arc=3pt,left=4pt,right=4pt,top=3pt,bottom=3pt,boxsep=1pt,halign=left]\bfseries\scriptsize Vocabulaire 广告 guǎnggào 顾客 gùkè 品牌 pǐnpái\end{tcolorbox}
\begin{tcolorbox}[enhanced,colback=poppink,colframe=poppink,coltext=white,boxrule=0.9pt,arc=3pt,left=4pt,right=4pt,top=3pt,bottom=3pt,boxsep=1pt,halign=left]\bfseries\scriptsize Grammaire 比较 bǐjiào 比 bǐ 越来越 yuè lái yuè\end{tcolorbox}
\begin{tcolorbox}[enhanced,colback=popgold,colframe=popgold,coltext=white,boxrule=0.9pt,arc=3pt,left=4pt,right=4pt,top=3pt,bottom=3pt,boxsep=1pt,halign=left]\bfseries\scriptsize Donner son avis 我觉得 wǒ juéde 因为 yīnwèi\end{tcolorbox}
\begin{tcolorbox}[enhanced,colback=popblue!70,colframe=popblue!70,coltext=white,boxrule=0.9pt,arc=3pt,left=4pt,right=4pt,top=3pt,bottom=3pt,boxsep=1pt,halign=left]\bfseries\scriptsize Mini-exposé comparer + argumenter\end{tcolorbox}
\end{tcbraster}

\legende{Carte mentale de la leçon : le marketing en Chine et au Cameroun}
\end{popfigure}

\begin{bilan}
\textbf{1. Lexique clé à connaître par cœur}
\begin{center}
\begin{tabularx}{\linewidth}{|X|X|X|X|}
\hline
\rowcolor{popblueL} Caractères & Pinyin & Nature & Traduction \\
\hline
市场 & shìchǎng & nom & marché \\
\hline
营销 & yíngxiāo & nom / verbe & marketing, commercialiser \\
\hline
广告 & guǎnggào & nom & publicité, annonce \\
\hline
顾客 & gùkè & nom & client \\
\hline
品牌 & pǐnpái & nom & marque \\
\hline
产品 & chǎnpǐn & nom & produit \\
\hline
价格 & jiàgé & nom & prix \\
\hline
质量 & zhìliàng & nom & qualité \\
\hline
打折 & dǎzhé & verbe & faire une réduction \\
\hline
网上购物 & wǎngshàng gòuwù & expr. & achats en ligne \\
\hline
直播 & zhíbō & nom / verbe & diffusion en direct (live) \\
\hline
影响 & yǐngxiǎng & verbe / nom & influencer, influence \\
\hline
\end{tabularx}
\end{center}

\textbf{2. Grammaire à maîtriser}
\begin{center}
\begin{tabularx}{\linewidth}{|X|X|X|}
\hline
\rowcolor{popblueL} Structure & Exemple & Traduction \\
\hline
A + 比 + B + adj. & 网购比市场便宜。Wǎnggòu bǐ shìchǎng piányi. & Les achats en ligne sont moins chers que le marché. \\
\hline
越来越 + adj. & 网购越来越流行。Wǎnggòu yuè lái yuè liúxíng. & Les achats en ligne sont de plus en plus populaires. \\
\hline
我觉得 + phrase & 我觉得广告很重要。Wǒ juéde guǎnggào hěn zhòngyào. & Je pense que la publicité est très importante. \\
\hline
因为 … 所以 … & 因为质量好，所以顾客多。Yīnwèi zhìliàng hǎo, suǒyǐ gùkè duō. & Parce que la qualité est bonne, il y a beaucoup de clients. \\
\hline
对 … 有影响 & 广告对顾客有影响。Guǎnggào duì gùkè yǒu yǐngxiǎng. & La publicité a une influence sur les clients. \\
\hline
\end{tabularx}
\end{center}

\textbf{3. Faits socioculturels à retenir}
\begin{itemize}
\item Au Cameroun : le marketing passe surtout par les marchés, la radio, la télévision, les affiches et le bouche à oreille ; les réseaux sociaux (Facebook, WhatsApp) progressent vite.
\item En Chine : le marketing est très numérique — applications comme WeChat (微信 Wēixìn), vente en direct (直播带货 zhíbō dàihuò), paiement mobile généralisé.
\item Point commun : dans les deux pays, la confiance et la relation personnelle avec le vendeur jouent un grand rôle.
\end{itemize}

\textbf{4. Méthode express : le mini-exposé comparatif}
\begin{enumerate}
\item Introduction : présenter le sujet (今天我说说喀麦隆和中国的营销。Jīntiān wǒ shuōshuo Kāmàilóng hé Zhōngguó de yíngxiāo.)
\item Décrire chaque pays en 2-3 phrases simples.
\item Comparer avec 比 et 越来越.
\item Conclure avec son avis : 我觉得…
\end{enumerate}
\end{bilan}

\begin{aretenir}[Checklist avant le Bac]
\begin{itemize}
\item[$\square$] Je connais le lexique clé : 市场, 营销, 广告, 顾客, 品牌, 产品, 价格, 质量.
\item[$\square$] Je sais écrire ces caractères dans le bon ordre des traits.
\item[$\square$] Je maîtrise la comparaison avec 比 (A + 比 + B + adjectif).
\item[$\square$] Je sais utiliser 越来越 + adjectif pour exprimer une évolution.
\item[$\square$] Je sais donner mon avis avec 我觉得 et le justifier avec 因为 … 所以 ….
\item[$\square$] Je sais construire 对 … 有 / 没有影响 sans oublier 有.
\item[$\square$] Je place correctement les compléments de temps et de lieu avant le verbe.
\item[$\square$] Je peux citer deux faits sur le marketing au Cameroun et deux sur la Chine.
\item[$\square$] Je fais attention aux tons de 买 mǎi (acheter) et 卖 mài (vendre).
\item[$\square$] Je suis capable de tenir un mini-exposé de 5 phrases en chinois.
\end{itemize}
\end{aretenir}

\findecours

\end{document}
